\documentclass[12pt,letterpaper]{report}
\usepackage{graphicx,amssymb,amsthm,mathtools,tikz-cd,stmaryrd}
\usepackage[english]{babel}
\usepackage[unicode=True]{hyperref}
\usepackage[a4paper]{geometry}

\newtheorem{definition}{Definition}[section]
\newtheorem{lemma}{Lemma}[section]
\newtheorem{theorem}{Theorem}[section]
\newtheorem{proposition}{Proposition}[section]
\newtheorem{corollary}{Corollary}[section]

\title{Representation Theory (Appendix)}
\date{}
\begin{document}
\maketitle

\appendix





\chapter{Representation Theory}
\section{Abstract Representation Theory}
\subsection{Semisimple Rings}
\begin{definition}Let $R$ be a ring. A left/right $R$-module is a commutative group $M$ with a biadditive map $\begin{array}{rrr}R\times M&\rightarrow&M\\(a,x)&\mapsto&ax/xa\end{array}$ s.t. $\\\left\{\begin{array}{rr}1x=x,\,\,\,\forall x\in M&x1=x,\,\,\,\,\forall x\in M\\(ab)x=a(bx),\,\,\forall a,b\in R,\forall x\in M&x(ab)=(xa)b,\,\,\,\forall a,b\in R,\forall x\in M\end{array}\right.$.\end{definition}
$\rightarrow$ if the ring $R$ is commutative, the notions of left and right $R$-module coincide.
\\\\e.g. the ring $R$ with left/right multiplication by itself is a left/right $R$-module, called the \textbf{left/right $R$-module.}
\begin{definition}Let $M$ be a $R$-module. A $R$-submodule of $M$ is a subgroup $N$ of $M$ s.t. $ax\in N,\,\,\,\forall a\in R,\,\,\,x\in N$.\end{definition}
e.g. a submodule of the left $R$-module $M=R$ is just a left ideal of $R$.
\begin{definition}If $M$ is a $R$-module and $N$ is a submodule of $M$, then the quotient group $M/N$ has a structure of $R$-module given by $a(x+N)=ax+N,\,\,\,\forall a\in R,\,\,\,x\in M$. This is called a quotient $R$-module.\end{definition}
\begin{definition}Let $M$ be a $R$-module and $(M_{i})_{i\in I}$ be a family of submodules of $M$. Then $M$ is the sum of the $M_{i}$ and write $M=\sum_{i\in I}M_{i}$ if $\forall x\in M,\,\,\,\exists x_{i}\in M_{i}$, all but finitely many equal to 0, s.t. $x=\sum_{i\in I}x_{i}$. Additionally $M$ is the direct sum of $M_{i}$ and write $M=\bigoplus_{i\in I}M_{i}$ if $\forall x\in M,\,\,\,\exists !x_{i}\in M_{i}$ s.t. $x=\sum_{i\in I}x_{i}$.\end{definition}
\begin{definition}Let $M,N$ be $R$-modules. A $R$-linear map from $M$ to $N$ is a morphism of abelian groups $\phi:M\rightarrow N$ s.t. $\phi(ax)=a\phi(x),\,\,\,\forall a\in R,x\in M$.\end{definition}
\begin{definition}An exact sequence of $R$-modules is an exact sequence of abelian groups where all the abelian groups are $R$-modules and all the maps are $R$-linear.\end{definition}
e.g. for $\phi:M\rightarrow N$ a $R$-linear map, $0\rightarrow\ker\phi\rightarrow M\rightarrow N\rightarrow\text{coker}\phi\rightarrow0$ is an exact sequence of $R$-modules.
\begin{definition}Let $R$ be a ring and $M$ be a $R$-module. Then $M$ is simple or irreducible if $M\neq0$ and if the only $R$-submodules of $M$ are 0 and $M$. Additionally $M$ is said to be semi-simple or completely reducible if for every submodule $N$ of $M$, there exists a submodule $N'$ of $M$ s.t. $M=N\oplus N'$. \end{definition}
\begin{lemma}If $M\neq0$ and $M$ is semisimple, then it has a simple submodule.\end{lemma}
pf. Let $x\in M\setminus\{0\}$ and set $M'=Rx$. Then $M'$ is semisimple and nonzero, so we may assume $M=M'$. Let $X$ be the set of submodules $N$ of $M$ s.t. $x\notin N$, ordered by inclusion. This set is nonempty because it contains the zero submodule. If $Y\subset X$ is totally ordered, $N=\cup_{N'\in Y}N'$ is a submodule of $M$. Also $x\notin N$ by definition, so $N$ is an element of $X$, so $\exists $ an upper bound of any totally ordered subset of $X$. Then by Zorn's lemma, $X$ has a maximal element $N$. As $x\notin N,\,\,\,N\neq M$. As $M$ is semisimple, there exists a submodule $N'$ of $M$ s.t. $M=N\oplus N'$ and $N'\neq 0$ because $N\neq M$. If $N'$ is not simple, choose a submodule $0\neq N''\subsetneq N'$, so that $N\oplus N''\notin X$ by maximality of $N$ in $X$, so $x\in N\oplus N''$, but then $N\oplus N''=M$ a contradiction. Thus $N'$ is simple and the lemma follows.
\begin{theorem}Let $R$ be a ring and $M$ be a $R$-module. Then TFAE\begin{enumerate}\item $M$ is semisimple\item $M$ is the direct sum of a family of simple submodules\item $M$ is the sum of a family of simple submodules\end{enumerate}\end{theorem}
pf.\begin{itemize}\item $2\rightarrow 3 $ direct sums are sums
\item$1\rightarrow 3$
\\\\assume $M$ is semisimple. Let $M'$ be the sum of all the simple submodules of $M$ and choose a submodule $M''$ of $M$ st. $M=M'\oplus M''$. If $M'\neq M,\,\,\,\,M''\neq0$, so by the previous lemma, $M''$has a simple submodule $N$, but then $N\subset M',$ a contradiction that $M',M''$ are in direct sum. So $M'=M$.
\item $3\rightarrow1$
\\\\let $(M_{i})$ be the family of all simple submodules of $M$. Assume $M=\sum_{i\in I}M_{i}$. Let $N$ be a submodule of $M$. Define the set $X$ be a subset $K$ of $I$ s.t. the number $N+\sum_{k\in K}M_{k}$ is direct, ordered by inclusion. This set is not empty since $\emptyset\in X$. If $Y\subset X$ is a totally ordered subset, let $K=\cup_{K'\in Y}K'$. Let $n\in N$ and $(m_{k})_{k\in K}\in \Pi_{k\in K}M_{k}$ s.t. $n+\sum m_{k}=0$. Let $K_{0}\subset K$ be a finite subset s.t. $m_{k}=0,\,\,\,\forall k\in K\setminus K_{0}$. Then $\forall k\in K_{0},\,\,\,\exists L_{k}\in Y$ s.t. $k\in L_{k}$. As $Y$ is totally ordered and $K_{0}$ is finite, $\exists L\in Y$ s.t. $\cup_{k\in K_{0}}L_{k}\subset L$. Then $L\in X$, so $n+\sum m_{k}=0$ implies $n=0$ and $m_{k}=0,\,\,\,\forall k\in K_{0}$. By the choice of $K_{0}$, we get $n$ and all $m_{k}$ are 0. So $K\in X$. Then by Zorn's lemma $X$ has a maximal element $J$. Let $M'=N\oplus \bigoplus_{j\in J}M_{j}$. Let $i\in I$. If $M_{i}\not\subset M',\,\,\,M'\cap M_{i}=0$ given $M_{i}$ being simple. Then the sum $M'+M_{i}=(N+\sum_{j\in J}M_{j})+M_{i}$ is direct, so $J\cup\{i\}\in X$, a contradiction to the maximality of $J$, so $M=M'$, so $M$ is semisimple. 
\end{itemize}
\begin{definition}Let $R$ be a ring and $M$ be a $R$-module. Then $M$ is said to be finitely generated if there exists a finite family $(x_{i})_{i\in I}$ of $M$ s.t. $M=\sum_{i\in I}Rx_{i}$. We say that $M$ is cyclic, if $\exists x\in M$ s.t. $M=Rx$.\end{definition}
\begin{theorem}Let $R$ be a ring. TFAE\begin{enumerate}\item all s.e.s. of $R$-modules split\item all $R$-modules are semisimple\item all finitely generated $R$-modules are semisimple\item all cyclic $R$-modules are semisimple\item the left regular $R$-module $R$ is semisimple\end{enumerate}\end{theorem}
pf.
\begin{itemize}\item $1\leftrightarrow 2$ by the definition of semisimple
\item $2\rightarrow 3\rightarrow 4\rightarrow 5$ is obvious
\item $5\rightarrow 2$
\\\\let $M$ be a $R$-module. If $x\in M, \,\,\,\,Rx\subset M$ is a quotient of $R$, so it is semisimple because $R$ is semisimple. As $M=\sum_{x\in M}Rx,$ by the previous theorem, $M$ is semisimple.

\end{itemize} 
\begin{definition}A $R$-module $P$ is called projective if for every surjective $R$-linear map $\pi:M\rightarrow N$ and every $R$-linear map $\phi:P\rightarrow N$, there exists a $R$-linear map $\psi:P\rightarrow M$ s.t. $\pi\psi=\phi$.\begin{tikzcd}&P\arrow[dl,dotted,"\psi"]\arrow[d,"\phi"]\\M\arrow[r,"\pi"]&N\end{tikzcd}\end{definition}
\begin{definition}A $R$-module $F$ is called free if there exists a family $(e_{i})_{i\in I}$ of elements of $F$ s.t. $F=\bigoplus_{i\in I}Re_{i}$ and s.t. $\forall i\in I$, the map $\begin{array}{rrr}R&\rightarrow &Re_{i}\\a&\mapsto&ae_{i}\end{array}$ is an isomorphism.\end{definition}

\begin{lemma}Let $P$ be a $R$-module. Then TFAE\begin{enumerate}\item $P$ is projective \item $P$ is a direct summand of a free $R$-module\end{enumerate}\end{lemma}
pf. \begin{itemize}\item $2\rightarrow 1$
\\\\\textbf{every free $R$-module is projective.}
\\pf. if $F=\bigoplus Re_{i},$ let $\pi:M\rightarrow N$ be a surjective $R$-linear map and $\phi:P\rightarrow N$ be a $R$-linear map. $\forall i\in I$ choose $x_{i}\in M$ s.t. $\pi(x_{i})=\phi(e_{i})$. Define a $R$-linear map $\begin{array}{rrr}\psi:P&\rightarrow&M\\e_{i}&\mapsto&x_{i}\end{array}$. Then $\pi\psi=\phi$ and it is projective.
\\\textbf{a direct summand of projective module is projective}
\\pf. let $P$ be projective and suppose $P=P'\oplus P''$. Let $\pi:M\rightarrow N$ be surjective $R$-linear map and $\phi':P'\rightarrow N$ be a $R$-linear map. Let $\phi=\phi'+0:P=P'\oplus P''\rightarrow N$. As $P$ is projective, $\exists \psi:P\rightarrow M$ s.t. $\pi\psi=\phi$. Write $\psi=\psi'\oplus\psi''$ where $\left\{\begin{array}{r}\psi':P'\rightarrow M\\\psi'':P''\rightarrow M\end{array}\right.$. Then $\pi\psi'=\phi'$ and the stmt follows.
\\Then free module is a projective module and its direct summand is also projective and the lemma follows.


\item $1\rightarrow 2$
\\\\suppose $P$ is projective. Let $M=\bigoplus_{x\in P}R$ and define $\begin{array}{rrr}\phi:M&\rightarrow& P\\(a_{x})_{x\in P}&\mapsto&\sum_{x\in P}a_{x}x\end{array}$. Then $\phi$ is $R$-linear and surjective, so there exists a $R$-linear map $\psi:P\rightarrow M$ s.t. $\phi\psi=id_{P}$. So $M\cong P\oplus \ker\phi$ and we have found a free $R$-module $M$ s.t. $P$ is a direct summand of $M$.

\end{itemize}
\begin{theorem}Let $R$ be a ring. TFAE\begin{enumerate}\item $R$ is a semisimple ring.\item all $R$-modules are projective\item all finitely generated $R$-modules are projective\item all cyclic $R$-modules are projective\end{enumerate}\end{theorem}
pf.\begin{itemize}\item $1\leftrightarrow 2$ by the definition
\item $2\rightarrow3\rightarrow 4$ is obvious
\item $4\rightarrow 1
$\\\\suppose all cyclic $R$-modules are projective. Let $I\subset R$ be a left ideal, then $R/I$ is a cyclic $R$-module, hence projective. So there exists a s.e.s. $0\rightarrow I\rightarrow R\rightarrow R/I\rightarrow0$ which splits. So $R$ is semisimple. 



\end{itemize}
\begin{theorem}[Schur's Lemma]Let $R$ be a ring, $M,N$ be $R$-modules and $u:M\rightarrow N$ be a $R$-linear map. \begin{enumerate}\item if $M$ is simple, then $u=0$ or $u$ is injective\item if $N$ is simple, then $u=0$ or $u$ is surjective\end{enumerate}\end{theorem}
pf. since $\ker u\subseteq M$, given $M$ being simple, $\ker u$ is either $M$ or 0 and the first part of the theorem follows. Similarly, Im$u\subseteq N$, given $N$ being simple, Im $u$ is either $N$ or 0 and the last part follows. 




\begin{definition}Let $R$ be a ring. Write $S(R)$ for the set of isomorphism classes of simple $R$-modules.\end{definition}
\begin{lemma}Let $M$ be a $R$-module and suppose there exists a simple $R$-module $S$ and a set $I$ and an isomorphism $\phi:M\rightarrow\bigoplus_{i\in I}S$. Then every simple $R$-submodule of $M$ is isomorphic to $S$.\end{lemma}
pf. Let $N$ be a simple $R$-submodule of $M$ and suppose $N$ is not isomorphic to $S$. Then $\forall i\in I,$ the composition of the projection on the $i$th summand $\bigoplus S\rightarrow S$ and of $\phi$ is a $R$-linear map $N\rightarrow S$ which has to be zero by Schur's lemma. Then $\phi$ is zero on $N$, which implies $N=0,$ a contradiction to the fact that $N$ is simple.

\begin{theorem}Let $R$ be a ring and $M$ be a semisimple $R$-module. $\forall S\in S(R)$, let $M_{S}$ be the sum of all the submodules of $M$ that are isomorphic to $S$. Then $\left\{\begin{array}{r}M=\bigoplus_{S\in S(R)}M_{S}\\\exists I_{S}$ s.t. $M_{S}\cong \bigoplus_{i\in I_{S}}S,\,\,\,\forall S\in S(R)\end{array}\right.$. Let $N$ be a submodule of $M$. If we write $N\cong\bigoplus_{S\in S(R)}N_{S}$ and $N_{S}\cong\bigoplus_{S\in J_{S}}S$ as in the previous statements, then $\forall S\in S(R),\,\,\,N_{S}=M_{S}\cap N,$ and we can find an injection $J_{S}\hookrightarrow I_{S}$.
\end{theorem}
pf.\begin{itemize}
\item $\sum_{S\in S(R)}M_{S}$ is the sum of all the simple submodules of $M$. As $M$ is semisimple, $M=\sum_{S\in S(R)}M_{S}$. Let $S\in S(R)$ and $X=S(R)\setminus \{S\}$. Also, let $N\equiv M_{S}\cap(\sum_{S'\in X}M_{S'})$. As $N$ is a submodule of $M$, it is semisimple. So if it is nonzero, it is a simple submodule of $M_{S}$, hence isomorphic to $S$ and also a simple submodule of $\sum_{S'\in X}M_{S'}$, hence isomorphic to an element of $X$, a contradiction. So $N=0$ and the stmt follows.


\item As $M_{S}$ is a submodule of the semisimple $R$-module $M$, it is semisimple. Then $M_{S}$ is the direct sum of a family of simple submodules by the theorem A.1.1. But every simple submodules of $M_{S}$ is isomorphic to $S$. 


\item
Let $S\in S(R)$. Then $N_{S}$ is a sum of simple $R$-modules isomorphic to $S$, so $N_{S}\subset M_{S}$. On the other hand, $N\cap M_{S}$ is a direct sum of simple submodules of $N$ and all these simple modules have to be isomorphic to $S$, so $N\cap M_{S}\subset N_{S}$. So we may assume $\left\{\begin{array}{r}M=M_{S}\\N=N_{S}\end{array}\right.$ for some $S\in S(R)$. Then we can reduce the case into the one s.t. if $S$ is simple $R$-module and $I,J$ are two sets s.t. we have an injective $R$-linear map $u:N\equiv\bigoplus_{j\in J}S\hookrightarrow M\equiv\bigoplus_{i\in I}S$, then there exists an injection $J\hookrightarrow I$. 
\\For any $R$-modules $M_{1},M_{2}$, let $\hom(M_{1},M_{2})\subset\hom_{R}(M_{1},M_{2})$ be the subgroup of $R$-linear maps that are 0 outside a finite length submodule of $M_{1}$. Then $\left\{\begin{array}{r}\hom(N,S)=\bigoplus_{j\in J}\text{End}_{R}(S)\subset\hom_{R}(N,S)=\Pi_{j\in J}\text{End}_{R}(S)\\\hom(M,S)=\bigoplus_{i\in I}\text{End}_{R}(S)\subset\hom_{R}(M,S)=\Pi_{i\in I}\text{End}_{R}(S)\end{array}\right.$. Here $\mathbb{D}\equiv\text{End}_{R}(S)$ is a division ring by Schur's lemma, and $\hom_{R}(M,S),\,\,\,\hom_{R}(N,S)$ are naturally $\mathbb{D}$-modules and $\hom(M,S),\hom(N,S)$ are $\mathbb{D}$-submodules. Moreover, $\exists\begin{array}{rrr}\phi:\hom(M,S)&\rightarrow&\hom(N,S)\\f&\mapsto&f\circ u\end{array}$ is $\mathbb{D}$-linear. Let $g\in \hom(N,S)$. As $M$ is a semisimple $R$-module, $\exists N'$ a $R$-submodule of $M$ s.t. $M=u(N)\oplus N'$. Define $f:M\rightarrow S$ by $f(u(x)+y)=g(x)$ if $x\in N,\,\,\,y\in N'$, since $u$ is injective, and clearly $R$-linear. So $\phi(f)=g$ so it is surjective. Then the stmt follows by the incomplete basis theorem. 
\end{itemize}
$\rightarrow $ the nonzero $M_{S}$ are called the \textbf{isotypic components of $M$} and $M=\bigoplus_{S\in S(R)}M_{S}$ is called \textbf{the isotypic decomposition of $M$. }
\begin{definition}A ring $R$ is called simple if $R\neq0$ and if the only ideals of $R$ are 0 and $R$.\end{definition}
\begin{proposition}If $R$ is a ring and $n\geq1$ is an integer, then any ideal $I$ of $M_{n}(R)$ is of the form $I=M_{n}(J)$, where $J$ is a uniquely determined ideal of $R$.\end{proposition}
pf. If $J$ is an ideal of $R,\,\,\,M_{n}(J)$ is an ideal of $M_{n}(R)$. Also if $J,J'$ are two ideals of $R$ s.t. $M_{n}(J)=M_{n}(J')$, then $J=J'$. So we just need to prove any ideal of $M_{n}(R)$ is of the form $M_{n}(J)$. Let $I$ be an ideal of $M_{n}(R)$ and let $J$ be the set of $(1,1)$-entries of elements of $I$. Let $x,y\in J$ and let $a\in R$. Choose matrices $X,Y\in I$ s.t. $(1,1)$-entries of $X,Y$ are $x,y$ respectively. Then $aX,Xa,X+Y$ are in $I$ and their respective $(1,1)$-entries are $ax,xa,x+y$ so $ax,xa,x+y\in J$. So $J$ is an ideal of $R$. Let's denote by $E_{ij}$ the elementary matrices in $M_{n}(R)$. If $X=(x_{ij})\in M_{n}(R)$, then $E_{ij}XE_{kl}=x_{jk}E_{il}$. If $X\in I,\,\,\,\,\forall j,k\in\{1,\dots,n\},\,\,\,E_{1j}XE_{k1}=x_{jk}E_{11}\in I$ so $x_{jk}\in J$ so $X\in M_{n}(J)$. This implies $I\subset M_{n}(J)$. Let $X=(x_{ij})\in M_{n}(J)$. Then $X=\sum_{ij}x_{ij}E_{ij}$. Fix $i,j\in\{1,\dots,n\}$. Choose $Y\in I$ s.t. $(1,1)$-entry of $Y$ is $x_{ij}$. Then $E_{i1}YE_{1j}=x_{11}E_{ij}\in I$ so $M_{n}(J)\subset I$. Thus the proposition follows. 
\begin{corollary}If $R$ is a simple ring, then so is $M_{n}(R),\,\,\,\forall n\geq1$.\end{corollary}


\begin{theorem}Let $\mathbb{D}$ be a division ring, $n\geq1$ be an integer and $R=M_{n}(\mathbb{D})$. Let $V=\mathbb{D}^{n}$. Then $V$ is viewed as a left $R$-module by considering $\mathbb{D}^{n}$ as a space of $n\times 1$ matrices, and also viewed as the right $\mathbb{D}$-module $\mathbb{D}_{\mathbb{D}}^{n}$. Then \begin{enumerate}\item the ring $R$ is simple, semisimple, left Artinian and left Noetherian.
\item $R$ has a unique simple left module which is $V$. As a left $R$-module, $R$ and $V^{n}$ are isomorphic. \item End$_{\mathbb{D}}(V)=R$\item End$_{R}(V)=\mathbb{D}$\end{enumerate}\end{theorem}
pf.\begin{itemize}
\item By the previous proposition, $R$ is simple because $\mathbb{D}$ is. As a left $\mathbb{D}$-vector space, $R$ is finite dimensional. Since left ideals of $R$ are $\mathbb{D}$-vector subspaces of $R$, $R$ is left Artinian and left Noetherian. Take a nonzero $R$-submodule $W$ of $V$. Let $w\in W\setminus\{0\}$ and write $w=\begin{bmatrix}w_{1}\\\vdots\\w_{n}\end{bmatrix}.$ Choose $i_{0}\in\{1,\dots,n\}$ s.t. $w_{i_{0}}\neq0$. Then $\begin{bmatrix}1\\0\\\vdots\\0\end{bmatrix}=(w_{i_{0}}^{-1}E_{1i_{0}})w\in W$. If $v=\begin{bmatrix}v_{1}\\\vdots\\v_{n}\end{bmatrix}$ is any element of $V$, then $v=\begin{bmatrix}v_{1}&0&\hdots&0\\\vdots&\vdots&&\vdots\\v_{n}&0&\hdots&0\end{bmatrix}\begin{bmatrix}1\\0\\\vdots\\0\end{bmatrix}\in W$ so $W=V$. As the left $R$-module, $R\cong V^{n}$, the ring $R$ is semisimple.
\item Let $M$ be a simple $R$-module. Choose $x\in M\setminus\{0\}$. Then the map $\begin{array}{rrr}u:R&\rightarrow &M\\a&\mapsto&ax\end{array}$ is surjective, so $M$ is isomorphic to a quotient of the $R$-module $R$. As $R\cong V^{n},\,\,\,\,M\cong V$. 
\item Consider the map $\begin{array}{rrr}\phi:R&\rightarrow&\text{End}_{\mathbb{D}}(V_{\mathbb{D}})\\a&\mapsto&(x\mapsto ax)\end{array}$. This map is well-defined because for all $a\in R,\,\,\,\forall x\in V,\,\,\,\lambda\in\mathbb{D},\,\,\,\phi(a)(x\lambda)=a(x\lambda)=(ax)\lambda=(\phi(a)(x))\lambda$, so $\phi(a)$ is $\mathbb{D}$-linear. Let $a=(a_{ij})\in R=M_{n}(\mathbb{D})$ s.t. $\phi(a)=0$. Then $\forall j\in\{1,\dots,n\}$, if $e_{j}\in V$ is the element with $j$th entry equal to 1 and all other entries equal to 0, we have $0=\phi(a)(e_{j})=\begin{bmatrix}a_{1j}\\\vdots\\a_{nj}\end{bmatrix}$ so $a=0$, which implies $\phi$ is injective. Let $u\in \text{End}_{\mathbb{D}}(V_{\mathbb{D}}).\,\,\,\forall j\in\{1,\dots,n\}$ if $e_{j}\in V$ is defined as above, $u(e_{j})=\begin{bmatrix}a_{1j}\\\vdots\\a_{nj}\end{bmatrix}. $ Let $a=(a_{ij})\in M_{n}(\mathbb{D})$. Then $\phi(a)(e_{j})=u(e_{j}),\,\,\,\forall j\in\{1,\dots,n\}$. As $(e_{1},\dots,e_{n})$ is a basis of $V_{\mathbb{D}}$ over $\mathbb{D}$, this implies $\phi(a)=u$ so $\phi$ is surjective.
\item Let $E=\text{End}_{R}(V)$. Write the action of the ring $E$ on $V$ on the right, that is, write $x(v)=vx,\,\,\,\forall v\in V,\,\,\,x\in E$. Let $\begin{array}{rrr}\psi:\mathbb{D}&\rightarrow& E\\\lambda&\mapsto&(v\mapsto v\lambda)\end{array}$. As for $\lambda$, this map is well-defined, i.e. $\psi(\lambda)$ is $R$-linear $\forall \lambda\in\mathbb{D}$ and it is a morphism of rings. Because $\mathbb{D}$ is a division algebra, $\psi$ is injective. Let $x\in E$ and define $\lambda\in\mathbb{D}$ by $e_{1}x=\lambda e_{1}+\sum_{j=2}^{n}\mu_{j}e_{j}$ s.t. $\mu_{j}\in\mathbb{D},\,\,\,\,e_{1},\dots,e_{n}\in V$ are as in the proof above. Let $v=\begin{bmatrix}v_{1}\\\vdots\\v_{n}\end{bmatrix}\in V$ and let $a=\begin{bmatrix}v_{1}&0&\hdots&0\\\vdots&\vdots&&\vdots\\v_{n}&0&\hdots&0\end{bmatrix}\in M_{n}(\mathbb{D})=R$. Then $v=ae_{1}$ so $vx=(ae_{1})x=a(e_{1}x)=a\begin{bmatrix}\lambda\\\mu_{2}\\\vdots\\\mu_{n}\end{bmatrix}=\begin{bmatrix}v_{1}\lambda\\\vdots\\v_{n}\lambda\end{bmatrix}=v\lambda$ so $x=\psi(\lambda)$ and $\psi $ is surjective. 



\end{itemize}
\begin{corollary}If $\mathbb{D,D}'$ are two division rings and $n,n'\geq1$ are two integers s.t. $M_{n}(\mathbb{D})\cong M_{n'}(\mathbb{D}')$ as rings, then $\mathbb{D}\cong\mathbb{D}'$ and $n=n'$.\end{corollary}
pf. Let $R=M_{n}(\mathbb{D})\cong M_{n'}(\mathbb{D}')$ and let $M$ be the unique simple $R$-module by the previous theorem 2. Then by the previous theorem 3, $\mathbb{D}'\cong\text{End}_{R}(M)\cong\mathbb{D}$. Hence $n=\dim_{\mathbb{D}}(M)=\dim_{\mathbb{D}'}(M)=n'$. 

\begin{definition}If $R$ is a ring, we denote by $R^{op}$ its opposite ring; it is isomorphic to $R$ as an additive group and its multiplication is given by $ab(R^{op})=ba(R)$.\end{definition}
\begin{theorem}[Double Centraliser Property]Let $R$ be a simple ring and $I$ be a nonzero left ideal of $R$. Let $D=\text{End}_{R}(I)$ and make $I$ a right $D^{op}$-module by setting $xu=u(x)$ if $\left\{\begin{array}{r}x\in I\\u\in D\end{array}\right.$. Then the map $\begin{array}{rrr}f:R&\rightarrow&\text{End}_{D^{op}}(I)\\a&\mapsto&(x\mapsto ax)\end{array}$ is an isomorphism of rings.\end{theorem}
pf. As $I\neq0,\,\,\,\ker f\neq R$. As $\ker f$ is an ideal of $R$ and $R$ is simple, $\ker f=0$, i.e. $f$ is injective. Let $E=\text{End}_{D^{op}}(I)$. Make $I$ a left $E$-module by setting $hx=h(x),\,\,\,\,\,\forall x\in I,\,\,h\in E$. Then $\forall x\in I,\,\,\,\forall h\in E,\,\,\,\,hf(x)=f(hx)$. Indeed if $a\in I,\,\,\,\begin{array}{rrr}r_{a}:I&\rightarrow&I\\y&\mapsto&ya\end{array}$ is in $D$ so $h(xa)=h(r_{a}(x))=h(xr_{a})=h(x)r_{a}=h(x)a$ so $(hf(x))(a)=h(xa)=h(x)a=f(h(x))(a)$. Thus $Ef(I)\subset f(I)$. Since $I$ is a nonzero left ideal of $R,\,\,\,IR$ is a nonzero ideal of $R$, hence $IR=R$ as $R$ is simple. Then $f(I)f(R)=f(R)$ and $E=Ef(R)=Ef(I)f(R)\subset f(I)f(R)=f(R)$, so $f$ is surjective. 
\begin{corollary}If $R$ is a simple ring with a minimal nonzero left ideal, then there exists a unique division ring $\mathbb{D}$ and a unique integer $n\geq1$ s.t. $R\cong M_{n}(\mathbb{D})$. In particular, a simple ring is left Artinian iff it is of the form $M_{n}(\mathbb{D})$ with $\mathbb{D}$ a division ring and $n\in\mathbb{Z}_{++}$.\end{corollary}
pf. Let $I\subset R$ be a minimal nonzero left ideal. Then $I$ is simple $R$-module, so $\mathbb{D}\equiv \text{End}_{R}(I)$ is a division ring by Schur's lemma. Make $I$ a right $\mathbb{D}^{op}$-module as in the previous theorem. Then $R\cong \text{End}_{\mathbb{D}^{op}}(I)$. Let $E=\text{End}_{\mathbb{D}^{op}}(I)$ and let $E_{f}=\{u\in E|rk(u)<\infty\}$ where $rk(u)$ is the dimension over $\mathbb{D}^{op}$ of the image of $u$. Then $E_{f}\neq0$ and $E_{f}$ is an ideal of $E$. Since $E\cong R$ is a simple ring, $E_{f}=E$, hence the identity of $I$ is in $E_{f}$ and $I$ is a finite dimensional right $\mathbb{D}^{op}$-vector space. The uniqueness follows by the previous corollary.

\begin{proposition}Let $R_{1},\dots,R_{n}$ be rings and let $R=R_{1}\times\dots\times R_{n}$. Then every $R$-module is of the form $M=M_{1}\times\dots\times M_{n}$ with $M_{i}$ a $R_{i}$-module for $1\leq i\leq n$.\end{proposition}
pf. In $R$, write $1=e_{1}+\dots+e_{n},\,\,\,\,\forall e_{i}\in R_{i}$. As an element of $R_{1}\times\dots\times R_{n},\,\,\,\, e_{i}$ is the $n$-tuple with $i$th entry equal to 1$\in R_{i}$ and 0 for all other entries. Then all $e_{i}$ are central in $R$. Let $M$ be a $R$-module, set $M_{i}=e_{i}M$. Then $R$ acts on $M_{i}$ through the obvious projection $R\rightarrow R_{i}$, so we need to show $M\cong M_{1}\times\dots\times M_{n}$. Consider the map $\begin{array}{rrr}M_{1}\times\dots\times M_{n}&\rightarrow&M\\(x_{1},\dots,x_{n})&\mapsto&\sum_{i}x_{i}\end{array}$ which is the $R$-linear map. If $x\in M,\,\,\,x=xe_{1}+\dots +xe_{n}$ with $xe_{i}\in M_{i},\,\,\,\forall i$ so $u$ is surjective. If $\sum x_{i}=0,\,\,\,x_{i}\in M_{i}$ then $\forall j\in\{1,\dots,n\},\,\,\,0=e_{j}(x_{1}+\dots+x_{n})=e_{j}x_{j}=x_{j}$ so $u$ is injective. 
\begin{corollary}Suppose $R=R_{1}\times\dots\times R_{n}$ as in the previous proposition.\begin{enumerate}\item $R$ is semisimple iff all the $R_{i}$ are semisimple\item Every simple $R$-module is of the form $0\times\dots\times0\times M_{i}\times0\times\dots\times0,\,\,\,i\in\{1,\dots,n\}$ and $M_{i}$ a simple $R_{i}$-module\end{enumerate}\end{corollary}\begin{corollary}Let $\mathbb{D}_{1},\dots,\mathbb{D}_{r}$ be division rings, and $n_{1},\dots,n_{r}\geq1$ be integers. Then $R\equiv M_{n_{1}}(\mathbb{D}_{1})\times\dots M_{n_{r}}(\mathbb{D}_{r})$ is a semisimple ring and its simple modules are $\mathbb{D}_{1}^{n_{1}},\dots,\mathbb{D}_{r}^{n_{r}}$.\end{corollary}

Let $R$ be a ring and $I$ be a left ideal of $R$. Denote $\mathcal{I}_{I}\subset R$ the sum of all the left ideals $I'$ of $R$ that are isomorphic to $I$ as $R$-modules. Then it is a left ideal of $R$. 
\begin{lemma}Let $R$ be a ring and $I$ be a minimal nonzero left ideal of $R$. Then $\mathcal{I}_{I}$ is an ideal of $R$. Moreover if $I,J$ are two nonisomorphic minimal nonzero left ideals of $R$, then $\mathcal{I}_{I}\mathcal{I}_{J}=0$.\end{lemma}
pf. Let $I'$ be a left ideal of $R$ s.t. $I'\cong I$ as $R$-modules, and let $a\in R$. Then we have a surjective $R$-linear map $\begin{array}{rrr}I'&\rightarrow& I'a\\x&\mapsto&xa\end{array}$ and $I'$ is a simple $R$-module, so $I'a=0$ or $I'a\cong I'$. If $I'a=0,\,\,\,I'a\subset\mathcal{I}_{I}$, otherwise $I'a$ is another left ideal of $R$ isomorphic to $I$, so $I'a\subset\mathcal{I}_{I}$. thus $\mathcal{I}_{I}$ is a right ideal. 
\\Let $J$ be another minimal nonzero left ideal of $R$, and suppose $\mathcal{I}_{I}\mathcal{I}_{J}\neq0$. Then there are left ideals $I',J'$ of $R$ and an element $a\in J'$ s.t. $I'\cong I,\,\,\,J'\cong J,\,\,\,I'a\neq 0$. As $J'$ is a simple $R$-module and $I'a\subset J'$ is a nonzero submodule of $J',\,\,\,\,I'a=J'$. As $I'$ is a simple $R$ module, the surjective map $\begin{array}{rrr}I'&\rightarrow&I'a\\x&\mapsto&xa\end{array}$ is an isomorphism. So $I\cong I'\cong I'a=J'\cong J$. 
\begin{theorem}[Artin-Wedderburn] Let $R$ be a semisimple ring. Let $(I_{i})_{i\in A}$ be a set of representatives of the isomorphism classes of minimal nonzero left ideals of $R$, and let $R_{i}=\mathcal{I}_{I_{i}}\subset R$. Then $A$ is finite so we choose an identification $A=\{1,\dots,r\}$. Moreover, all the $R_{i}$ are rings with units, $R\cong R_{1}\times\dots\times R_{r}$ as rings and $\forall i\in\{1,\dots,r\}$, there exists a unique division ring $\mathbb{D}_{i}$ and a unique integer $n_{i}\geq1$ s.t. $R_{i}\cong M_{n_{i}}(\mathbb{D}_{i})$.\end{theorem}
pf. As $R$ is a semisimple ring, $R$ is a direct sum of simple submodules by the theorem A.1.1. So $R=\bigoplus_{i\in A}R_{i}$ by the theorem A.1.5. Then $A$ is finite. 
\\By the previous lemma, every $\mathcal{I}_{i}$ is an ideal of $R$. As $R=\bigoplus_{i=1}^{r}R_{i}$, all the $R_{i}$ are rings with units and $R\cong R_{1}\times\dots\times R_{r}$ as rings. Fix $i\in\{1,\dots,r\}$. Let $J\neq0$ be an ideal of $R_{i}$. As $J$ is also an ideal of $R$, it contains a minimal nonzero left ideal $I$ of $R$. By definition of $I_{1},\dots,I_{r},\,,\,\,\,\exists j\in\{1,\dots,r\}$ s.t. $I\cong I_{j}$ as $R$-modules. As $I\subset J\subset R_{i},\,\,\,j=i$. So $R_{i}=\mathcal{I}_{I_{i}}=\mathcal{I}_{I}$. Hence it suffices to show if $I'$ is a left ideal of $R$ s.t. $I\cong I'$, then $I'\subset J$. Fix such a $I'$ and let $\phi:I\rightarrow I'$ be an isomorphism of $R$-modules. As $R$ is a semisimple ring, there exists a left ideal $I''$ of $R$ s.t. $R=I\oplus I''$. Write $1=e+e'',\,\,\,e\in I,\,\,\,e''\in I''$. Then $e^{2}=e,\,\,\,I=Re$. So $I'=\phi(I)=\phi(Re)=\phi(Re^{2})=\phi((Re)e)=\phi(Ie)=I\phi(e)\subset J$. 
\\\\\\e.g. let $K$ be a field and suppose $R$ is a semisimple $K$-algebra that is finite dimensional as a $K$-vector space. Then its simple factors $R_{i}$ are also $K$-algebras and so are the division rings $\mathbb{D}_{i}$ with $\dim_{K}\mathbb{D}_{i}<\infty$. In particular, if $K$ is algebraically closed, all the $\mathbb{D}_{i}$ are equal to $K$, so $R\cong M_{n_{1}}(K)\times\dots\times M_{n_{r}}(K)$.  

\begin{definition}Let $R$ be a ring. The Jacobson radical of $R$ is the intersection of all the maximal left ideals of $R$, denoted by $rad(R)$.\end{definition}



\subsection{The Representation Ring}
\begin{definition}Let $R$ be a ring. Define $K(R)$ to be the quotient of the free abelian group on the basis elements $[M]$, for $M$ a finite length $R$-module, by all the relations of the form $[M]=[M']+[M'']$, where $0\rightarrow M'\rightarrow M\rightarrow M''\rightarrow0$ is an exact sequence of $R$-modules. It is called the Grothendieck group of the category of finite length $R$-modules. Define $PK(R)$ to be the quotient of the free abelian group on the basis elements $[P]$, for $P$ a finite length projective $R$-module, by all the relations of the from $[P]=[P']+[P'']$ where $0\rightarrow P'\rightarrow P\rightarrow P''\rightarrow0$ is an exact sequence of $R$-modules. It is called the Grothendieck group of the category of finite length projective $R$-modules. \end{definition}
$\rightarrow 0\rightarrow0\rightarrow0\rightarrow0\rightarrow0$ gives $[0]+[0]=[0]$ in $K(R),PK(R)$ so $[0]=0$. Then for an isomorphism of $R$-modules $M\rightarrow M',\,\,\,0\rightarrow M\rightarrow M'\rightarrow0$ gives $[M]=[M']+[0]=[M']$ in $K(R)$, and similarly for $PK(R)$. So if $0=M_{n}\subset M_{n-1}\subset\dots\subset M_{0}=M,\,,\,\,[M]=\sum_{i=1}^{n}[M_{i-1}/M_{i}]$ in $K(R)$. 
\begin{proposition}As a group, $K(R)$ is the free abelian group with basis $\{[M]:M\in S(R)\}$, where $S(R)$ is the set of isomorphism classes of simple $R$-modules.\end{proposition}
pf. Let $S$ be the free abelian group on the set $S(R)$, and denoted by $(e_{M})_{M\in S(R)}$ its canonical basis. Define $\begin{array}{rrr}\phi:S&\rightarrow&K(R)\\e_{M}&\mapsto&[M]\end{array},\,\,\,\,\,\begin{array}{rrr}\psi:K(R)&\rightarrow&S\\M&\mapsto&\sum e_{M_{i-1}/M_{i}}\end{array}$ where $M=M_{0}\supset M_{1}\supset\dots\supset M_{n}=0$. If $0\rightarrow M'\rightarrow M\rightarrow_{u}M''\rightarrow0$ is a s.e.s. of finite length $R$-modules, choose Jordan-H\"older series $\left\{\begin{array}{rr}M'=M_{0}'\supset\dots\supset M_{n}'=0\\M''=M_{0}''\supset\dots\supset M_{m}''=0\end{array}\right.$. $\forall 0\leq j\leq m,\,\,\,M_{j}=u^{-1}(M_{j}'')$. Then $M=M_{0}\supset\dots\supset M_{m}=M_{0}'\supset\dots\supset M_{n}'=0$ is a Jordan-H\"older series for $M$, which gives the desired equality. So $\psi$ is well-defined and $\psi,\phi$ are inverses of each other, so the proposition follows.
\begin{proposition}If $R$ is a semisimple ring, the $K(R)=PK(R)$ and for all finite length $R$-modules $M,M',\,\,\,\,\,[M]=[M']$ in $K(R)$ iff $M\cong M'$.\end{proposition}
pf. If $R$ is semisimple, every $R$-module is projective, so $K(R)=PK(R)$.
\\Let $M,M'$ be two $R$-modules of finite length. As $R$ is semisimple, $\left\{\begin{array}{r}M\cong\bigoplus_{N\in S(R)}N^{\oplus r_{N}}\\M'\cong\bigoplus_{N\in S(R)}N^{\oplus r_{N}'}\end{array}\right.$. Then $[M]=\sum r_{N}[N],\,\,\,\,[M']=\sum r_{N}'[N]$. Then by the previous proposition, $[M]=[M']$ iff $r_{N}=r_{N}',\,\,\,\forall N\in S(R)$, which is equivalent to $M\cong M'$.
\begin{definition}The representation ring of $G$ over $k$, denoted by $R_{k}(G)$ is defined as $R_{k}(G)=K(k[G])$. Also $P_{k}(G)=PK(k[G]),\,\,\,\,S_{k}(G)=S(k[G])$.\end{definition}
$\rightarrow \forall [M],[M']\in R_{k}(G),\,\,\,[M][M']=[M\otimes_{k}M']$ defines a biadditive map $R_{k}(G)\otimes R_{k}(G)\rightarrow R_{k}(G)$ with unit element equal to $[\mathbb{1}]$. By the properties of tensor products, it is commutative. 
\begin{proposition}If $M$ is a $k[G]$-module and $N$ is a projective $k[G]$-module, then $k[G]$-module $M\otimes_{k}N$ is also projective.\end{proposition}
pf. As $N$ is projective over $k[G]$, it is a direct summand of some free $k[G]$-module $k[G]^{\oplus I}$ and then $M\otimes_{k}N$ is a direct summand of $M\otimes_{k}k[G]^{\oplus I}=(M\otimes_{k}k[G])^{\oplus I}$. So we can reduce into the case where $M\otimes_{k}k[G]$. Let $\underline{M}$ be the $k$-vector space $M$, considered as a representation of $G$ with trivial action. Define a $k$-linear map $\begin{array}{rrr}u:\underline{M}\otimes_{k}k[G]&\rightarrow&M\otimes_{k}k[G]\\x\otimes g&\mapsto&gx\otimes g\end{array}$. Then $k$-linear map $\\\begin{array}{rrr}v:M\otimes_{k}k[G]&\rightarrow&\underline{M}\otimes_{k}k[G]\\x\otimes g&\mapsto&g^{-1}x\otimes g\end{array}$ is an inverse of $u$, so $u$ is an isomorphism of $k$-modules. Also $\forall g,h\in G,\,\,\,\,\forall x\in \underline{M},\,\,\,\left\{\begin{array}{r}h\cdot u(x\otimes g)=h\cdot(gx\otimes g)=(hgx)\otimes(hg)\\u(h\cdot(x\otimes g))=u(x\otimes (hg))=(hgx)\otimes (hg)\end{array}\right.$ gives $u$ is $G$-equivariant. Since $\underline{M}\otimes_{k}k[G]$ is a direct sum of copies of $k[G]$, the $k[G]$-module $M\otimes_{k}k[G]$ is free, and the proposition follows. 
\\\\e.g. $G=\mathbb{Z}/n\mathbb{Z}$, $k$ a field containing all the $n$th root of unity, and char$(k)\not|n$. Then $R_{k}(G)\cong\mathbb{Z}^{n}$ as a group and $R_{k}(G)$ is isomorphic to the group algebra $\mathbb{Z}[\mathbb{Z}/n\mathbb{Z}]$ as a ring. 

\subsection{Induction and Frobenius Reciprocity}

Let $R$ be a ring and $\alpha:H\rightarrow G$ be a morphism of groups. Then $\alpha$ gives a morphism of rings $\phi:R[H]\rightarrow R[G]$. 
\begin{definition}If $M$ is a $R[G]$-module, we can see it as a $R[H]$-module by making $R[H]$ act via the morphism $\phi:R[H]\rightarrow R[G]$. The resulting $R[H]$-module is denoted by $\text{Res}_{H}^{G}M$ and called the restriction of $M$ to $H$.\end{definition}
$\rightarrow $if $u:M\rightarrow N$ is a morphism of $R[G]$-modules, Res$_{H}^{G}(u):\text{Res}_{H}^{G}M\rightarrow\text{Res}_{H}^{G}N$ for the same map $u$, seen as a morphism of $R[H]$-modules.

\begin{definition}If $M$ is a $R[H]$-module, set $\text{Ind}_{H}^{G}M=R[G]\otimes_{R[H]}M$ where $R[H]$ acts on the left on $R[G]$ via the morphism $\phi$. This is called the induction of $M$ from $H$ to $G$.
\end{definition}
$\rightarrow $ if $u:M\rightarrow N$ is a morphism of $R[H]$-modules, Ind$_{H}^{G}(u):\text{Ind}_{H}^{G}M\rightarrow\text{Ind}_{H}^{G}N$ for the $R[G]$-linear map $id\otimes u$. 
\begin{definition}If $M$ is a $R[H]$-module, we set CoInd$_{H}^{G}M=\hom_{R[H]}(R[G],M)$ where $R[G]$ is seen as a left $R[H]$-module via $\phi:R[H]\rightarrow R[G]$. We make this into a left $R[G]$-module using the right regular action of $R[G]$ on itself, i.e. $\forall x\in R[G],\,\,\,u\in\text{CoInd}_{H}^{G}M,\,\,\,(x\cdot u)(y)=u(yx),\,\,\,\forall y\in R[H]$. The $R[G]$-module CoInd$_{H}^{G}$ is called the coinduction of $M$ from $H$ to $G$.\end{definition}
$\rightarrow$ if $u:M\rightarrow N$ is a morphism of $R[H]$-modules, we write CoInd$_{H}^{G}(u):\text{CoInd}_{H}^{G}M\rightarrow \text{CoInd}_{H}^{G}N$ for the $R[G]$-linear map sending $v\in\hom_{R[H]}(R[G],M)$ to $u\circ v$.
\\\\\textbf{Let $M$ be a $R[H]$-module. Then restriction map $R[G]\rightarrow M$ along the inclusion $G\subset R[G]$ induces an isomorphism of $R$-modules}\begin{equation*}\psi:\text{CoInd}_{H}^{G}M\rightarrow\{f:G\rightarrow M\,\,|\,\,\forall h\in H,\,\,\,\forall g\in G,\,\,\,f(\alpha(h)g)=hf(g)\}\end{equation*}
pf. Because $G$ generates the $R[H]$-module $R[G]$, so a $R[H]$-linear map $u:R[G]\rightarrow M$ is uniquely determined by its restriction to $G$. So $\psi$ is injective. Let $f:G\rightarrow M$ satisfying the condition in the formula above, and define a $R$-linear map $\begin{array}{rrr}R[G]&\rightarrow&M\\\sum_{g\in G}c_{g}g&\mapsto&\sum c_{g}f(g)\end{array}$. Then $\forall h\in H,\,\,\,x=\sum_{g\in G}c_{g}c\in R[G],\,\,\,u(hx)=u(\sum c_{g}\alpha(h)g)=\sum c_{g}f(\alpha(h)g)=\sum c_{g}hf(g)=hu(x),$ i.e. $u$ is $R[H]$-linear. Then $\psi(u)=f$, so $\psi$ is surjective. 

\begin{proposition}If $\alpha:H\rightarrow G$ is injective, then $\phi:R[H]\rightarrow R[G]$ makes $R[G]$ into a free $R[H]$-module. More precisely, let $(g_{i})_{i\in I}$ be a complete set of representatives of $\alpha(H)\subset G$ in $G$, then $(g_{i})_{i\in I}$ is a basis of $R[G]$ as a left/right $R[H]$-module.\end{proposition}
\begin{definition}Suppose $G=\{1\}$. Then for every $R[H]$-module $M$, Ind$_{H}^{1}M$ is called the $R$-module of coinvariants of $M$ under $H$ and denote by $M_{H}$.\end{definition}
$\rightarrow M_{H}=M/(hm-m)$ the quotient of $M$ by the $R$-submodule generated by all $hm-m$. 
\begin{lemma}Let $M$ be a $R[H]$-module.\begin{enumerate}\item if $\alpha:H\rightarrow G$ is surjective, then Ind$_{H}^{G}M=M_{\ker \alpha}$, with the following action of $G$; $\forall g\in G,\,\,\,\,\,x\in M_{\ker\alpha}$, choose preimage $h\in H$ of $g$ by $\alpha$ and $m\in M$ of $x$ by the obvious quotient map, and then $gx$ is the image in $M_{\ker\alpha}$ of $hm$.\item Ind$_{H}^{G}M=\text{Ind}_{\alpha(H)}^{G}M_{\ker\alpha}$\end{enumerate}\end{lemma}
pf.\begin{itemize}
\item Let $u:\text{Ind}_{H}^{G}M=R[G]\otimes_{R[H]}M\rightarrow M_{\ker\alpha}$ be defined as follows; $\forall a\in R[G],\,\,\,m\in M$, choose $b\in R[H]$ s.t. $\phi(b)=a$ and take for $u(a\otimes m)$ the image of $bm$ in $M_{\ker\alpha}$. Because $\ker\phi$ is the $R[H]$-submodule of $R[H]$, generated by $h-1,\,\,\,h\in\ker\alpha$ this does not depend on the choice of $h$. Since the formula $u(a\otimes m)$ is additive in $a$ and $m$ and takes the same value on $(a\phi(x),m)$ and $(a,xm)$ if $x\in R[H]$, the map is well defined on $R[G]\otimes_{R[H]}M.\,\,\,\,\begin{array}{rrr}v':M&\rightarrow&R[G]\otimes_{R[H]}M\\m&\mapsto&1\otimes m\end{array}$. If $m\in M,\,\,\,h\in\ker\alpha,\,\,\,v'(hm-m)=1\otimes hm-1\otimes m=\alpha(h)\otimes m-1\otimes m=0$ so $v'$ defines a $R$-linear map $v:M_{\ker\alpha}\rightarrow\text{Ind}_{H}^{G}M$. Here $u$ and $v'$ are inverses to each other, so the statement follows.
\item $\begin{array}{rr}\text{Ind}_{H}^{G}M=\text{Ind}_{\alpha(H)}^{G}\text{Ind}_{H}^{\alpha(H)}M& $transitivity of $\otimes\\=\text{Ind}_{\alpha(H)}^{G}M_{\ker\alpha}& $ statement 1$\end{array}$

\end{itemize}
\begin{theorem}For every exact sequence $0\rightarrow M'\rightarrow M\rightarrow M''\rightarrow 0$ of $R[H]$-modules, the sequence \begin{equation*}\text{Ind}_{H}^{G}M'\rightarrow\text{Ind}_{H}^{G}M\rightarrow\text{Ind}_{H}^{G}M''\rightarrow0\end{equation*} is exact. Moreover $0\rightarrow\text{Ind}_{H}^{G}M'\rightarrow\text{Ind}_{H}^{G}M\rightarrow\text{Ind}_{H}^{G}M''\rightarrow0$ is exact if $\\\left\{\begin{array}{r}\alpha:H\rightarrow G$ is injective $\\\ker\alpha<\infty,\,\,\,R$ is semisimple, $|\ker\alpha|$ is invertible in $R\end{array}\right.$ \end{theorem}
pf. By the property of tensor product the first stmt follows.
\\The right $R[H]$-module $R[G]$ is free by the previous proposition and therefore the tensor products by $R[G]$ over $R[H]$ preserves exact sequence. 
\\By the previous lemma, Ind$_{H}^{G}M=\text{Ind}_{\alpha(H)}^{G}M_{\ker\alpha}$ for every $R[H]$-module $M$. Let $f:M\rightarrow N$ be an injective $R[G]$-linear map. As $R[G]$ is semisimple, there exists a $R[G]$-module $M'$ of $N$ s.t. $N=M\oplus M'$. Then $N_{G}=M_{G}\oplus M_{G}'$ and $M_{G}\rightarrow N_{G}$ is injective. So coinvariants preserves injection and therefore we have the exact sequence preserved.

\begin{definition}Suppose $G=\{1\}$. Then for every $R[H]$-module $M,\,\,\,\text{CoInd}_{H}^{1}M$ is called the $R$-module invariants of $M$ under $H$ and denoted by $M^{H}$.\end{definition}
$\rightarrow M^{H}=\{m\in M|\forall h\in H,\,\,\,hm=m\}$.
\begin{lemma}Let $M$ be a $R[H]$-module.\begin{enumerate}\item if $\alpha:H\rightarrow G$ is surjective, CoInd$_{H}^{G}M=M^{\ker\alpha}$ with the following action of $G$; $\forall g\in G,\,\,\,\,\forall m\in M^{\ker\alpha}$, choose a preimage $h\in H$ of $g$ by $\alpha$ and then $gm$ is defined to be $hm$.\item CoInd$_{H}^{G}M=\text{CoInd}_{\alpha(H)}^{G}M^{\ker\alpha}$\end{enumerate}\end{lemma}
pf. \begin{itemize}\item Let $u:\text{CoInd}_{H}^{G}M=\hom_{R[H]}(R[G],M)\rightarrow M^{\ker\alpha}$ be defined as follows; if $f:R[G]\rightarrow M$ is a $R[H]$-liner map, take $u(f)=f(1)$. Then it is well-defined because $\forall h\in\ker\alpha,\,\,\,\,hf(1)=f(\alpha(h)1)=f(1)$. Let $v:M^{\ker\alpha}\rightarrow \hom_{R[H]}(R[G],M)$ be the map defined follows; if $m\in M,\,\,\,a\in R[G]$, choose $b\in R[H]$ s.t. $\phi(b)=a$ and set $v(m)(a)=bm$. This does not depend on the choice of $b$ because $\ker\phi$ is the $R[H]$-submodule of $R[H]$ generated by all $h-1,\,\,\,h\in\ker\alpha$. Then $u$ and $v$ are inverses to each other and the stmt follows.
\item $\begin{array}{rr}\text{CoInd}_{H}^{G}M=\text{CoInd}_{\alpha(H)}^{G}\text{CoInd}_{H}^{\alpha(H)}M&$ properties of homomorphisms$\\=\text{CoInd}_{\alpha(H)}^{G}M^{\ker\alpha}&$ by the statement 1$\end{array}$. 

\end{itemize}
\begin{theorem}For every exact sequence $0\rightarrow M'\rightarrow M\rightarrow M''\rightarrow0$ of $R[H]$-modules, the sequence $0\rightarrow\text{CoInd}_{H}^{G}M'\rightarrow\text{CoInd}_{H}^{G}M\rightarrow\text{CoInd}_{H}^{G}M''$ is exact. Moreover $0\rightarrow\text{CoInd}_{H}^{G}M'\rightarrow\text{CoInd}_{H}^{G}M\rightarrow\text{CoInd}_{H}^{G}M''\rightarrow0$ if $\\\left\{\begin{array}{r}\alpha:H\rightarrow G$ is injective $\\\ker\alpha<\infty,\,\,\,R$ is semisimple, $|\ker\alpha|$ is invertible in $R\end{array}\right.$\end{theorem}
pf. By the left-exactness of hom, the first stmt follows.
\\By the previous proposition, the left $R[H]$-module $R[G]$ is free, so taking $\hom_{R[H]}(R[G],\cdot)$ preserves exact sequences.
\\By the previous lemma, $\text{CoInd}_{H}^{G}M=\text{CoInd}_{\alpha(H)}^{G}M^{\ker\alpha}$ for every $R[H]$-module $M$. Let $M\rightarrow N$ be a surjective $R[G]$-linear map. As $R[G]$ is semisimple, there exists a $R[G]$-submodule $N$ of $M$ s.t. $M=N\oplus N'$ and by the definition of invariants, $M^{G}=N^{G}\oplus N^{'G}$ and $M^{G}\rightarrow N^{G}$ is surjective. So invariants preserves surjection and therefore we get a s.e.s. 
\begin{definition}Let $M$ be a $R[H]$-module. We denote by $\begin{array}{rrr}\epsilon_{M}:\text{ Res}_{H}^{G}\text{CoInd}_{H}^{G}M&\rightarrow &M\\u&\mapsto&u(1)\end{array}$, where $u\in \hom_{R[H]}(R[G],M)$. Additionally denote by $\begin{array}{rrr}\eta_{M}:M&\rightarrow&\text{Res}_{H}^{G}\text{Ind}_{H}^{G}M\\m&\mapsto&1\otimes m\end{array}$.\end{definition}
$\rightarrow \forall x\in R[H],\,\,\forall u\in \hom_{R[H]}(R[G],M),\,\,\,\epsilon_{M}(x\cdot u)=(x\cdot u)(1)=u(x)=xu(1)$ so it is $R[H]$-linear. $\eta_{M}$ is clearly $R[H]$-linear. 
\begin{theorem}[Frobenius Reciprocity] Let $M$ be a $R[G]$-module and $N$ be a $R[H]$-module. \begin{enumerate}\item the morphism of groups \begin{equation*}\begin{array}{rrr}\Phi:\hom_{R[G]}(M,\text{CoInd}_{H}^{G}N)&\rightarrow&\hom_{R[H]}(\text{Res}_{H}^{G}M,N)\\u&\mapsto&\epsilon_{N}\circ \text{Res}_{H}^{G}(u)\end{array}\end{equation*} is an isomorphism. 

\item the morphism of groups \begin{equation*}\begin{array}{rrr}\Phi':\hom_{R[G]}(\text{Ind}_{H}^{G}N,M)&\rightarrow&\hom_{R[H]}(N,\text{Res}_{H}^{G}M)\\u&\mapsto&\text{Res}_{H}^{G}(u)\circ\eta_{N}\end{array}\end{equation*} is an isomorphism.\end{enumerate}\end{theorem}
pf. Let $R_{1}=R[H],\,\,\,R_{2}=R[G]$ for simplicity.
\begin{itemize}\item \begin{equation*}\begin{array}{rrr}\Psi:\hom_{R[H]}(\text{Res}_{H}^{G}M,N)=\hom_{R_{1}}(M,N)&\rightarrow&\hom_{R[G]}(M,\text{CoInd}_{H}^{G}N)=\hom_{R_{2}}(M,\hom_{R_{1}}(R_{2},N))\\u:M\rightarrow N&\mapsto&m\mapsto(a\mapsto(u(am))\end{array}\end{equation*} Let $u$ be a $R_{1}$-linear map from $M$ to $N$. Then $\Psi(u):M\rightarrow\hom_{R_{1}}(R_{2},N)$ is the $R_{2}$-linear map $m\mapsto(a\mapsto(u(am)))$ and $\Phi\Psi(u):M\rightarrow N$ is the $R_{1}$-linear map sending $m\in M$ to $(\Psi(u)(m))(1)=u(m)$. So $\Phi\Psi(u)=u$. Conversely, let $v$ be a $R_{2}$-linear map from $M$ to $\hom_{R_{1}}(R_{2},N)$. Then $\Phi(v)$ is the $R_{1}$-linear map $m\mapsto v(m)(1)$ and $\forall m\in M, \,\,\,\,a\in R_{2},\,\,\,(\Psi\Phi(v)(m))(a)=\Phi(v)(am)=(v(am))(1)=(a\cdot v(m))(1)=v(m)(a),\,\,\,\,\Psi\Phi(v)=v$. Thus they are inverses to each other, so isomorphism.
\item \begin{equation*}\begin{array}{rrr}\Psi':\hom_{R[H]}(N,\text{Res}_{H}^{G}M)=\hom_{R_{1}}(N,M)&\rightarrow&\hom_{R[G]}(\text{Ind}_{H}^{G}N,M)=\hom_{R_{2}}(R_{2}\otimes_{R_{1}}N,M)\\u:N\rightarrow M&\mapsto&a\otimes n\mapsto au(n)\end{array}\end{equation*}Let $u$ be a $R_{1}$-linear map from $N$ to $M$. Then $\forall n\in N,\,\,\,\Phi'\Psi'(u)(n)=\Psi'(u)(1\otimes n)=u(n),\,\,\,\,\Phi'\Psi'(u)=u$. Let $v$ be a $R_{2}$-linear map from $R_{2}\otimes_{R_{1}}N$ to $M$. Then $\forall a\in R_{2},\,\,\,\forall n\in N,\,\,\,\,\,\Psi'\Phi'(v)(a\otimes n)=a\Phi'(v)(n)=av(1\otimes n)=v(a\otimes n),\,\,\,\,\Psi'\Phi'(v)=v$. Thus they are inverses to each other, so isomorphism.

\end{itemize}



\section{Finite Group Representation}
\subsection{Introduction}
Let $R$ be a ring and $G$ be a group. A $R[G]$-module is a $R$-module $M$ with a $R$-linear action of $G$, i.e. a morphism of monoids $\rho:G\rightarrow \text{End}_{R}(M)$. This is called a $R$-linear representation of $G$ on the $R$-module $M$, denoted by $(M,\rho)$. The representation $(M,\rho)$ is called \textbf{irreducible} if the $R[G]$-module $M$ is simple and it is called\textbf{ faithful} if $\rho$ is injective. A $R[G]$-linear map is also called a $G$-equivariant map. A sub-$R[G]$-module is also called a $\textbf{subrepresentation}$. \textbf{The regular representation of $G$} is the representation corresponding to the left regular $R[G]$-module. 
\\\indent Let $\begin{array}{rrr}\epsilon:R[G]&\rightarrow&R\\\sum_{g\in G}\alpha_{g}g&\mapsto&\sum_{g\in G}\alpha_{g}\end{array}$ be called the augmentation map. It is a surjective $R$-linear ring homomorphism, and its kernel is called the augmentation ideal of $R[G]$.
\begin{theorem}Let $R$ be a ring and $G$ be a finite group. Then $R[G]$ is a semisimple ring iff $R$ is a semisimple ring and $|G|$ is invertible in $R$.\end{theorem}
pf. \begin{itemize}\item Let $M$ be a $R[G]$-module and let $N$ be a $R[G]$-submodule of $M$. As $R$ is semisimple, $\exists N'$, a $R$-submodule of $M$ s.t. $M=N\oplus N'$, so there exists a surjective $R$-linear map $f:M\rightarrow N$ s.t. $f|_{N}=id_{N}$. Define $\begin{array}{rrr}F:M&\rightarrow&N\\v&\mapsto&|G|^{-1}\sum_{g\in G}g^{-1}f(gv)\end{array}$. Then it is $R$-linear and $\forall g\in G,\,\,\,\forall v\in M,\,\,\,\,F(gv)=|G|^{-1}\sum_{h\in G}h^{-1}f(hgv)=|G|^{-1}g\sum_{h\in G}(hg)^{-1}f(hgv)=gF(v),$ i.e. $F(gv)=gF(v)$. Thus $F(v)=|G|^{-1}\sum g^{-1}f(gv)=|G|^{-1}\sum g^{-1}gv=v,$ i.e. $F|_{N}=id_{N}$. Here $\forall v\in M,\,\,\,F(v)\in N$ and therefore $FF(v)=F(v)$, hence $F(v-F(v))=0,\,\,\,v=(v-F(v))+F(v)\in \ker F+N$. If $v\in N\cap\ker F,\,\,\,\,v=F(v)=0$. Thus $M=N\oplus \ker F$.
\item Assume $R[G]$ is semisimple. Then there exists the augmentation map $\epsilon$. It makes $R$ into a $R[G]$-module, which is automatically semisimple by assumption. As $R[G]$ acts on $R$ through its quotient $R[G]/\ker\epsilon\cong R$, the $R$-module $R$ is semisimple, so $R$ is a semisimple ring. It remains to show $|G|$ is invertible in $R$; by the remark below it suffices to show $|G|\neq0$ in each $\mathbb{D}_{i}$ appearing in the Artin-Wedderburn decomposition of $R$. Suppose $|G|=0$ in $\mathbb{D}\equiv\mathbb{D}_{i}$. As $M_{n_{i}}(\mathbb{D})[G]\cong M_{n_{i}}(\mathbb{D}[G])$ is a quotient of $R[G]$, it is semisimple, hence so is $\mathbb{D}[G]$ by the proposition on ideals of matrix rings. Let $c=\sum_{g\in G}g\in\mathbb{D}[G]$. Then $c\neq0$, $c$ is central, and $c^{2}=|G|c=0$, so the left ideal $I=\mathbb{D}[G]c$ satisfies $Ic=\mathbb{D}[G]c^{2}=0$. By semisimplicity write $\mathbb{D}[G]=I\oplus J$ and $1=e+f$ with $e\in I,\,\,f\in J$. Then $c=ec+fc=cf\in J$, using $ec\in Ic=0$ and the centrality of $c$, while $c=1c\in I$, so $c\in I\cap J=0$, a contradiction.


\end{itemize}
By the Artin-Wedderburn theorem, $R=M_{n_{1}}(\mathbb{D}_{1})\times\dots\times M_{n_{r}}(\mathbb{D}_{r})$, where $\mathbb{D}_{j}$ are division rings and $n_{j}\geq1,\,\,\,\forall j$. For every $i,$ let $R_{i}=M_{n_{i}}(\mathbb{D}_{i})$. Then we have a surjective morphism of rings $R[G]\rightarrow R_{i}$ and therefore $R_{i}$ is also semisimple. Also $|G|$ is invertible in $R$ iff it is invertible in every  $R_{i}$, iff it is invertible in every $\mathbb{D}_{i}$, iff it is nonzero in $\mathbb{D}_{i}$.

\begin{theorem}\begin{enumerate}\item There are uniquely determined $k$-division algebras $\mathbb{D}_{1},\dots,\mathbb{D}_{r}$ and integers $n_{1},\dots,n_{r}\geq1$ s.t. $\dim_{k}(\mathbb{D}_{i})$ is finite for every $i$ and \begin{equation*}k[G]/\text{rad}(k[G])\cong M_{n_{1}}(\mathbb{D}_{1})\times\dots\times M_{n_{r}}(\mathbb{D}_{r})\end{equation*}A complete set of representatives of the isomorphism classes of irreducible representations of $G$ is given by $V_{1}\equiv\mathbb{D}_{1}^{n_{1}},\dots,V_{r}\equiv\mathbb{D}_{r}^{n_{r}}$ and we have a $G$-equivariant isomorphism \begin{equation*}k[G]/\text{rad}(k[G])\cong V_{1}^{n_{1}}\oplus \dots\oplus V_{r}^{n_{r}}\end{equation*}\item If $k$ is algebraically closed, then $\mathbb{D}_{i}=k,\,\,\,\forall i$, and \begin{equation*}\sum_{V}(\dim V)^{2}=\sum_{i=1}^{r}\dim(V_{i})^{2}=\sum_{i=1}^{r}n_{i}^{2}=\dim(k[G]/\text{rad}(k[G]))\leq\dim(k[G])=|G|\end{equation*} where the first sum is taken over the isomorphism classes of irreducible representations of $G$. The inequality is an equality iff the characteristic of $k$ does not divide $|G|$.\end{enumerate}\end{theorem}



\begin{definition}If $R$ is a ring and $G$ is a group, denote by $S_{R}(G)$ the set of isomorphism classes of irreducible representations of $G$ on $R$-modules.\end{definition}
e.g.\begin{itemize}\item Fix $R,G$ as above. The trivial representation of $G$ over $R$ is the representation of $G$ on $R$ given by the augmentation map $\epsilon:R[G]\rightarrow R$ by the trivial action of $G$ on $R$, denoted by $\mathbb{1}$. 
\item If $G$ is the symmetric group $\mathcal{G}_{n}$ and $R$ is any ring, we denote by $sgn$ the representation of $G$ on $R$ given by the sign morphism $G\rightarrow\{\pm1\}$ composed with the obvious map $\{\pm1\}\rightarrow R.$
\item Let $G$ be the quaternion group $Q=\{\pm1,\pm i,\pm j,\pm k\}\subset\mathbb{H}$ with the multiplication given by that of $\mathbb{H}$. Then $\mathbb{R}[G]\cong \mathbb{R}\times\mathbb{R}\times\mathbb{R}\times\mathbb{R}\times\mathbb{H}$ so $S_{\mathbb{R}}(G)$ has 5 elements.  
\end{itemize}

\begin{proposition}If $G$ is a finite abelian group and $k$ is an algebraically closed field, then every simple $k[G]$-module is of dimension 1 over $k$.\end{proposition}
pf. Let $V$ be a simple $k[G]$-module. Then any nonzero element $v\in V$ gives a surjective $k$-linear map $k[G]\rightarrow V$ sending $x\in k[G]$ to $xv$. Then $V$ is a finite dimensional vector space, and so is $\text{End}_{k[G]}(V)$. The finite dimensional $k$-algebra $\text{End}_{k[G]}(V)$ is also a $k$-division algebra by the Schur's lemma, $\text{End}_{k[G]}(V)=k$. Moreover, as $G$ is abelian, $k[G]$ is a commutative ring, so the action of any element of $k[G]$ on $V$ is $k[G]$-linear. Thus the action of $k[G]$ on $V$ is given by a map of $k$-algebras $k[G]\rightarrow\text{End}_{k[G]}(V)=k\subset\text{End}_{k}(V)$, So $\forall v\in V,$ the subspace $kv$ is a $k[G]$-submodule of $V$. As $V$ is irreducible, $\dim V=1$. 
\\\\e.g. Let $G=\mathbb{Z}/n\mathbb{Z},\,\,\,n\geq1,\,\,\,k$ be an algebraically closed field.\begin{itemize}\item char $k\not|n\,\,:$ fix a primitive $n$th root $\zeta_{n}$ of 1 in $k$. The $k$-algebra $k[G]$ is semisimple, so $k[G]=k\times\dots\times k,$ where $n$ factors $k$ correspond to the $n$ irreducible representations of $G$ given by the maps $\begin{array}{rrr}G=\mathbb{Z}/n\mathbb{Z}&\rightarrow&k^{\times}\\i&\mapsto&\zeta_{n}^{i}\end{array}$.
\item char $k|n\,\,:$ let $p=\text{char}(k)$ and write $n=p^{r}m$ with $(m,p)=1$. Then $G=\mathbb{Z}/p^{r}\mathbb{Z}\times\mathbb{Z}/m\mathbb{Z}.\,\,\,\,(1,0),(0,1)\in\mathbb{Z}/p^{r}\mathbb{Z}\times\mathbb{Z}/m\mathbb{Z}$ have images $a,b\in k^{\times}$ under the map $G\rightarrow k^{\times}$ corresponding to the representation. Then $a^{p^{r}}=1$ so $a=1$ given char$(k)=p$. Also $b$ can be any of the $m$ solutions of the equation $x^{m}=1$ in $k$. So we can get $m$ irreducible representations of $G$, and $k[G]/\text{rad}(k[G])\cong k^{m}$. \end{itemize}



\subsection{Characters in char$(k)=0$}

\begin{definition}A function $f:G\rightarrow k$ is called central if $\forall g,h\in G,\,\,\,\,\,f(gh)=f(hg)$. We write $\mathcal{C}(G,k)$ for the $k$-algebra of central functions from $G$ to $k$.\end{definition}
\begin{definition}Let $(V,\rho)$ be a representation of $G$. The character of $V$ is the function \begin{equation*}\begin{array}{rrr}\chi_{V}:G&\rightarrow&k\\g&\mapsto&\text{T}(\rho(g))\end{array}\end{equation*}\end{definition}

\begin{proposition}Let $(V,\rho_{V}),\,\,\,\,(W,\rho_{W})$ be representations of $G$. Then $\left\{\begin{array}{r}\chi_{V}(1)=\dim_{k}V\\\chi_{V}\in\mathcal{C}(G,k)\\\chi_{V\oplus W}=\chi_{V}+\chi_{W}\\\chi_{V\otimes_{k}W}=\chi_{V}\chi_{W}\end{array}\right.$.\end{proposition}
pf. Choose $k$-bases $\left\{\begin{array}{rr}(e_{1},\dots,e_{n})&V\\(f_{1},\dots,f_{m})&W\end{array}\right.$. Let $g\in G$. In the chosen bases of $V,W,$ write $\left\{\begin{array}{r}\rho_{V}(g)\\\rho_{W}(g)\end{array}\right.$ as matrices $\left\{\begin{array}{r}(x_{ij})\\(y_{ij})\end{array}\right.$. Then $(e_{i}\otimes f_{j})$ is a basis for $V\otimes_{k}W$, and the matrix $\rho_{V}(g)\otimes \rho_{W}(g)$ in this basis is $(x_{ii'}y_{jj'})$. Hence $\chi_{V\otimes_{k}W}(g)=\sum_{i=1}^{n}\sum_{j=1}^{m}x_{ii}y_{jj}=(\sum_{i}x_{ii})(\sum_{j}y_{jj})=\chi_{V}(g)\chi_{W}(g)$.
\begin{corollary}The map $V\mapsto\chi_{V}$ induces a morphism of rings $R_{k}(G)\rightarrow\mathcal{C}(G,k)$.\end{corollary}
For $K/k$ an extension of fields, given every representation of $G$ over $k,\,\,\,\chi_{V\otimes_{k}K}=\chi_{V}$, so \begin{tikzcd}R_{k}(G)\arrow[r]\arrow[d]&\mathcal{C}(G,k)\arrow[d]\\R_{K}(G)\arrow[r]&\mathcal{C}(G,K)\end{tikzcd} where the left vertical arrow is given by $[V]\mapsto[V\otimes_{k}K]$ and the right vertical arrow is the obvious inclusion. 
\begin{definition}Let $V,W$ be representations of $G$. The $k$-vector space $\hom_{k}(V,W)$ becomes a representation of $G$ if we make $g\in G$ act by $(g\cdot f)(v)=gf(g^{-1}v),\,\,\,\forall f\in\hom_{k}(V,W)$ and $\forall v\in V$. In particular, $V^{*}\equiv\hom_{k}(V,k)$ is a representation of $G$ and $(g\cdot f)(v)=f(g^{-1}v),\,\,\,\,\forall g\in G,\,\,\,f\in V^{*},\,\,\,v\in V$.\end{definition}
\begin{definition}For every representation $V$ of $G$, we set \begin{equation*}V^{G}=\{v\in V|\forall g\in G,\,\,\,gv=v\}\end{equation*} This is the space of invariants of $G$ in $V$.\end{definition}
$\rightarrow \hom_{G}(V,W)=\hom_{k}(V,W)^{G}$
\begin{proposition}Let $V,W$ be representations of $G$. Then the map $\\\begin{array}{rrr}\phi:V^{*}\otimes_{k}W&\rightarrow&\hom_{k}(V,W)\\v\otimes w&\mapsto&(v\mapsto f(v)w)\end{array}$ is a $G$-equivariant isomorphism.\end{proposition}
pf. It is well-defined, because the map $\begin{array}{rrr}V\times W&\rightarrow&\hom_{k}(V,W)\\(v,w)&\mapsto&(v\mapsto f(v)w)\end{array}$ is $k$-bilinear. Let $f\in V^{*},\,\,\,w\in W,\,\,\,g\in G,\,\,\,v\in V$. Then $\\\left\{\begin{array}{r}(\phi(g(f\otimes w)))(v)=(\phi((gf)\otimes(gw))(v)=(gf)(v)gw=f(g^{-1}v)gw\\(g\phi(f\otimes w))(v)=g(\phi(f\otimes w)(g^{-1}v))=gf(g^{-1}v)w\end{array}\right.$ so $\phi$ is $G$-equivariant. Let $(e_{i})$ be a basis of $W$ as a $k$-vector space. Then we have $k$-linear isomorphisms $\begin{array}{rrr}u:\oplus_{i}\hom_{k}(V,k)&\rightarrow&\hom_{k}(V,W)\\(f_{i})&\mapsto&(v\mapsto\sum f_{i}(v)e_{i})\end{array},\,\,\,\,\begin{array}{rrr}v:\oplus V^{*}&\rightarrow&V^{*}\otimes_{k}W\\(f_{i})&\mapsto&\sum f_{i}\otimes e_{i}\end{array}$. Then $\phi=u\circ v^{-1}$ is an isomorphism. 
\begin{proposition}Let $V,W$ be representations of $G$. Then $\forall g\in G,\,\,\,\\\left\{\begin{array}{r}\chi_{V^{*}}(g)=\chi_{V}(g^{-1})\\\chi_{\hom_{k}(V,W)}(g)=\chi_{V}(g^{-1})\chi_{W}(g)\end{array}\right.$.\end{proposition}
pf. Write $\rho$ for the action morphism $G\rightarrow\text{End}_{k}(V)$. Choose a basis $\mathcal{B}$ of $V$ as a $k$-vector space, and let $M$ be the matrix of $\rho(g)^{-1}$ in $\mathcal{B}$.Then $g$ acts by the matrix $M'$ on $V^{*}$, so the result follows. Then the first stmt implies the second stmt. 
\begin{theorem}[Schur's Lemma]Let $V,W$ be irreducible representations of $G$. Then $\hom_{G}(V,W)=0$ unless $V\cong W$, and $\text{End}_{G}(V)$ is a finite-dimensional $k$-division algebra. If moreover $k$ is algebraically closed, then $\text{End}_{G}(V)=k$.\end{theorem}
pf. The first part follows from the first Schur's lemma.
\\We have a map of rings $\begin{array}{rrr}K&\rightarrow&\text{End}_{G}(V)\\\lambda&\mapsto&\lambda\cdot 1\end{array}$ which is injective, given $k$ a field. Let $a\in R$. Then $\begin{array}{rrr}m_{a}:\text{End}_{G}(V)&\rightarrow&\text{End}_{G}(V)\\x&\mapsto&ax\end{array}$ is $k$-linear. Here $\text{End}_{G}(V)$ is finite dimensional $k$-vector space and $k$ is algebraically closed, so $m_{a}$ has at least one eigenvalue, i.e. $\exists \lambda\in k$ s.t. $\ker(m_{a}-\lambda\cdot Id)=\ker(m_{a-\lambda})\neq0$. This implies $a-\lambda$ is not invertible. Since $\text{End}_{G}(V)$ is a division algebra, $a-\lambda=0\Leftrightarrow a=\lambda\in k$.  


\begin{lemma}Suppose $k$ is algebraically closed and let $(V,\rho)$ be an irreducible representation of $G$. Then \begin{equation*}\sum_{g\in G}\chi_{V}(g)=\left\{\begin{array}{rr}0&V\not\cong\mathbb{1}\\|G|&V\cong\mathbb{1}\end{array}\right.\end{equation*}\end{lemma}
pf. Extend $\rho:G\rightarrow\text{End}_{k}(V)$ to $\rho:k[G]\rightarrow\text{End}_{k}(V)$ and let $c=\sum_{g\in G}g\in k[G]$. Then $c$ is central in $k[G]$, so $\rho(c)\in\text{End}_{k}(V)$ is a $G$-equivariant endomorphism of $V$. Then by the Schur's lemma, $\exists \lambda\in k$ s.t. $\rho(c)=\lambda id_{V}$, so $\sum_{g\in G}\chi_{V}(g)=T(\rho(c))=\lambda\dim V$. Moreover $\forall h\in G,\,\,\,\lambda\rho(h)=\rho(c)\rho(h)=\rho(ch)=\rho(c)=\lambda id_{V}$. So if $\lambda\neq0,\,\,\,V\cong\mathbb{1},\,\,\,\sum\chi_{V}(g)=\sum1=|G|$. 
\begin{lemma}Let $V$ be a representation of $G$. Then \begin{equation*}\sum_{g\in G}\chi_{V}(g)=|G|\dim_{k}(V^{G})\end{equation*}\end{lemma}
pf. Assume wlog $k$ is algebraically closed ( we has seen $\chi_{V\otimes_{k}K})=\chi_{V}$ and we can make $V$ algebraically closed by tensoring with algebraically closed $K$ ). Write $V=\oplus_{W\in S_{k}(G)}W^{\oplus n_{W}}$. Then $\sum_{g\in G}\chi_{V}(g)=\sum_{W\in S_{k}(G)}n_{W}\sum_{g\in G}\chi_{W}(g)=n_{\mathbb{1}}|G|$ by the previous lemma.
\\If $W\in S_{k}(G),\,\,\,W\cong\mathbb{1},\,\,\,W^{G}=0$ since $W^{G}$ is a subrepresentation of $W$. Then $V^{G}=\mathbb{1}^{n_{\mathbb{1}}},\,\,\,\dim_{k}(V^{G})=n_{\mathbb{1}}$.
\begin{theorem}Let $V,W$ be representations of $G$. Then \begin{equation*}\frac{1}{|G|}\sum_{g\in G}\chi_{V}(g^{-1})\chi_{W}(g)=\frac{1}{|G|}\sum_{g\in G}\chi_{V^{*}}(g)\chi_{W}(g)=\dim_{k}(\hom_{G}(V,W))\end{equation*}\end{theorem}
pf. $\frac{1}{|G|}\sum\chi_{V}(g^{-1})\chi_{W}(g)=\frac{1}{|G|}\sum\chi_{V^{*}}(g)\chi_{W}(g)=\frac{1}{|G|}\sum\chi_{\hom_{k}(V,W)}(g)=\dim_{k}(\hom_{k}(V,W)^{G})=\dim_{k}(\hom_{G}(V,W))$ by the previous lemma.
\begin{corollary}If $V,W$ are irreducible and nonisomorphic, then $\sum_{g\in G}\chi_{V^{*}}(g)\chi_{W}(g)=0$.\end{corollary}
pf. by the previous theorem and the Schur's lemma.
\begin{corollary}Suppose $k$ is algebraically closed and let $V$ be a representation of $G$. Then $V$ is irreducible iff $\sum_{g\in G}\chi_{V^{*}}(g)\chi_{V}(g)=|G|$.\end{corollary}
pf. By the previous theorem $\sum_{g\in G}\chi_{V^{*}}(g)\chi_{V}(g)=|G|\dim_{k}(\text{End}_{G}(V))$. If $V$ is irreducible, End$_{G}(V)=k$ by Schur's lemma so the one direction follows.
\\Conversely, suppose $\sum_{g\in G}\chi_{V^{*}}(g)\chi_{V}(g)=|G|$. Write $v=\oplus_{i\in I}V_{i}^{\oplus n_{i}}$ where $V_{i}$ are irreducible and pairwise nonisomorphic. Then by the previous corollary, $\sum\chi_{V^{*}}(g)\chi_{V}(g)=\sum n_{i}n_{j}\sum\chi_{V_{i}^{*}}(g)\chi_{V_{j}}(g)=|G|\sum n_{i}^{2}$. So only one of the $n_{i}$ can be nonzero, and moreover it has to be equal to 1, hence $V$ is irreducible.
\begin{corollary}The family $(\chi_{V})_{V\in S_{k}(G)}$ is linearly independent in $\mathcal{C}(G,k)$.\end{corollary}
pf. Suppose $\sum_{W\in S_{k}(G)}\alpha_{W}\chi_{W}=0$, with $\alpha_{W}\in k$. Then $\forall V\in S_{k}(G),\,\,\,0=\sum_{g\in G}(\sum_{W\in S_{k}(G)}\alpha_{W}\chi_{W}(g))\chi_{V^{*}}(g)\\=\sum\alpha_{W}\sum\chi_{W}(g)\chi_{V^{*}}(g)=\alpha_{V}|G|\rightarrow\alpha_{V}=0$, so LI follows.
\begin{corollary}For any representations $V,V'$ of $G,\,\,\,\,\,V\cong V'$ iff $\chi_{V}=\chi_{V'}$.\end{corollary}
pf. Write $\left\{\begin{array}{r}V=\bigoplus_{W\in S_{k}(G)}W^{\oplus n_{W}}\\V'=\bigoplus_{W\in S_{k}(G)}W^{\oplus n_{W}'}\end{array}\right.$. Then $\left\{\begin{array}{r}\chi_{V}=\sum n_{W}\chi_{W}\\\chi_{V'}=\sum n_{W}'\chi_{W}\end{array}\right.$. By the previous corollary, $\chi_{V}=\chi_{V'}$ iff $n_{W}=n_{W}',\,\,\,\forall W\in S_{k}(G)\Leftrightarrow V\cong V'$. 


\begin{lemma}If $f\in \mathcal{C}(G,k)$ is s.t. $\sum_{g\in G}f(g)\chi_{W^{*}}(g)=0,\,\,\,\forall W\in S_{k}(G),$ then $f=0$.\end{lemma}
pf. If $\rho:G\rightarrow \text{End}_{k}(V)$ is a representation of $G$, we set $\rho(f)=\sum f(g)\rho(g)\in\text{End}_{k}(V)$. Then $\forall g\in G,\,\,\,\,\rho(g)\rho(f)=\sum_{h\in G}f(h)\rho(gh)=\sum_{h\in G}f(h)\rho(ghg^{-1}g)=\sum f(ghg^{-1})\rho(ghg^{-1})\rho(g)=\rho(f)\rho(g)$ because $f$ is central. So if $V$ is irreducible, $\rho(f)\in\text{End}_{G}(V)=k$ by Schur's lemma, hence $\rho(f)=\lambda id_{V},\,\,\,\,\lambda\in k$. So $\lambda \dim(V)=T(\rho(f))=\sum_{g\in G}f(g)\chi_{V}(g)=0\rightarrow \lambda=0$ and $\rho(f)=0$. By the semisimplicity of $k[G]$, $\rho(f)=0$ for any representation $\rho$ of $G$. Then $0=\rho_{reg}(f)1=\sum f(g)g$ in $k[G]$, so $f(g)=0,\,\,\,\forall g\in G$.
\begin{theorem}Suppose that the field $k$ is algebraically closed. Then the family $(\chi_{W})_{W\in S_{k}(G)}$ is a basis of $\mathcal{C}(G,k)$.\end{theorem}
pf. They are L.I by the corollary A.2.4. so we just need to show it spans $\mathcal{C}(G,k)$. Let $f\in \mathcal{C}(G,k)$. Then $\forall W\in S_{k}(G),\,\,\,\,\alpha_{W}=\frac{1}{|G|}\sum_{g\in G}f(g)\chi_{W^{*}}(g)$. Set $f'=f-\sum_{W\in S_{k}(G)}\alpha_{W}\chi_{W}$. Then $\forall W\in S_{k}(G),\,\,\,\sum_{g\in G}f'(g)\chi_{W^{*}}(g)=\sum f(g)\chi_{W^{*}}(g)-\alpha_{W}\sum\chi_{W}(g)\chi_{W^{*}}(g)=0\rightarrow f'=0$ and $f$ corresponds to $\chi_{W}$ by the previous corollaries. So it spans. 
\begin{corollary}Suppose that $k$ is algebraically closed. Then $|S_{k}(G)|=\dim_{k}\mathcal{C}(G,k)$ and this is also equal to the number of conjugacy classes in $G$.\end{corollary}
\begin{corollary}Suppose $k$ is algebraically closed.Then the map $\begin{array}{rrr}R_{k}(G)&\rightarrow&\mathcal{C}(G,k)\\V&\mapsto&\chi_{V}\end{array}$ induces an isomorphism of $k$-algebra $R_{k}(G)\otimes_{\mathbb{Z}}k\rightarrow \mathcal{C}(G,k)$.\end{corollary}



Let $G_{1},G_{2}$ be two groups. If $V_{1}/V_{2}$ is a representation of $G_{1}/G_{2}$, then $V_{1}\otimes_{k}V_{2}$ becomes a representation of $G_{1}\times G_{2}$ with action $(g_{1},g_{2})(v_{1}\otimes v_{2})=(g_{1}v_{1})\otimes(g_{2}v_{2})$. 
\begin{theorem}Suppose $k$ is algebraically closed, and let $V_{1}/V_{2}$ be a representation of $G_{1}/G_{2}$. The representation $V_{1}\otimes_{k}V_{2}$ of $G_{1}\times G_{2}$ is irreducible iff $V_{1},V_{2}$ are both irreducible. Every irreducible representation of $G_{1}\times G_{2}$ is of the form $V_{1}\otimes_{k}V_{2}$.\end{theorem}
pf. The character of $V\equiv V_{1}\otimes_{k}V_{2}$ is the map $(g_{1},g_{2})\mapsto\chi_{V_{1}}(g_{1})\chi_{V_{2}}(g_{2})$. So $<V,V>_{G}=\frac{1}{|G_{1}\times G_{2}|}\sum_{(g_{1},g_{2})}\chi_{V_{1}}(g_{1})\chi_{V_{2}}(g_{2})=<V_{1},V_{1}>_{G_{1}}<V_{2},V_{2}>_{G_{2}}$. All three brackets are nonnegative integers, $<V,V>_{G}=1$ iff $<V_{1},V_{1}>_{G_{1}}=<V_{2},V_{2}>_{G_{2}}=1$. So if $V_{1},V_{2}$ are irreducible representation of $G_{1}/G_{2}$, then $V_{1}\otimes V_{2}$ is an irreducible representation of $G$. Thus we can define a map $S_{k}(G_{1})\times S_{k}(G_{2})\rightarrow S_{k}(G_{1}\times G_{2})$. Since the representations $V_{1}\otimes V_{2}\sim V_{1}'\otimes V_{2}'$ iff $V_{1}\sim V_{1}',\,\,\,V_{2}\sim V_{2}'$, the map is injective. Let $X(G_{1})$ be the set of conjugacy classes in $G_{1}$, and similarly for $X(G_{2}),X(G_{1}\times G_{2})$. Then we have a surjective map $u:G_{1}\times G_{2}\rightarrow X(G_{1}\times G_{2})$. Then $(g_{1},g_{2}),\,\,\,\,(g_{1}',g_{2}')\in G_{1}\times G_{2}$ are conjugate iff $g_{i}\sim g_{i}'$ are conjugate in $G_{i}$. So $u$ induces a bijection $X(G_{1})\times X(G_{2})\rightarrow X(G_{1}\times G_{2})$. As $|S_{k}(G)|=|X(G)|,\,\,\,\,\forall G\in\{G_{1},G_{2},G_{1}\times G_{2}\}$, so $|S_{k}(G_{1})\times S_{k}(G_{2})|=|S_{k}(G_{1}\times G_{2})|$. So source and target have the same cardinality and the map is a bijection and therefore the theorem follows. 
\\\\\indent Fix a finite group $G$ and we assume $k$ is any field s.t. char$(k)\not||G|$. If $f_{1},f_{2}\in\mathcal{C}(G,k)$, write $<f_{1},f_{2}>_{G}=\frac{1}{|G|}\sum_{g\in G}f_{1}(g)f_{2}(g^{-1})=\frac{1}{|G|}\sum_{g\in G}f_{1}(g^{-1})f_{2}(g)$. Then the previous theorem can be expressed as; for any representations $V,W$ of $G$, we have $<\chi_{V},\chi_{W}>=\dim_{k}(\hom_{G}(V,W))$. 


\begin{definition}Let $H$ be a subgroup of $G$. If $f\in\mathcal{C}(H,k)$, define $\\\begin{array}{rrr}\text{Ind}_{H}^{G}f:G&\rightarrow&k\\g&\mapsto&\frac{1}{|H|}\sum_{s\in G|_{s^{-1}gs\in H}}f(s^{-1}gs)\end{array}$.\end{definition}
\begin{theorem}$\begin{array}{r}\forall f\in \mathcal{C}(H,k),\,\,\,\,\,\text{Ind}_{H}^{G}f\in\mathcal{C}(G,k)\\\text{Ind}_{H}^{G}\chi_{V}=\chi_{\text{Ind}_{H}^{G}V}\end{array}$ for a representation $V$ of $H$.\end{theorem}
pf. Let $(V,\rho_{V})$ be a representation of $H$ and write $(W,\rho_{W})=\text{Ind}_{H}^{G}(V,\rho_{V})$. Let $g_{1},\dots,g_{r}$ be a system of representatives of $G/H$. Then $k[G]=\bigoplus_{i=1}^{r}g_{i}k[G]$, so $W=k[G]\otimes_{k[H]}V=\bigoplus_{i=1}^{r}W_{i}$ with $W_{i}=g_{i}k[H]\otimes_{k[H]}V$. Let $g\in G$. Then $\exists \sigma\in\mathcal{G}_{r}$ and $h_{1},\dots,h_{r}\in H$ s.t. $gg_{i}=g_{\sigma(i)}h_{i},\,\,\,\forall i$. Then $\rho_{W}(g)(W_{i})=W_{\sigma(i)},\forall i$. Choose a  basis $(e_{1},\dots,e_{n})$ of $V$. Then $\forall i,\,\,\,\,(g_{i}\otimes e_{1},\dots,g_{i}\otimes e_{n})$ is a basis of $W_{i}$, and $\rho_{W}(g)(g_{i}\otimes e_{s})=g_{\sigma(i)}\otimes\rho_{V}(h_{i})e_{s}$. Thus $T(\rho_{W}(g))=\sum_{i|_{\sigma(i)=i}}T(\rho_{V}(h_{i}))=\sum_{i|_{g_{i}^{-1}gg_{i}\in H}}f(g_{i}^{-1}gg_{i})$. But $G=\coprod_{i=1}^{r}g_{i}H$ and if $s\in g_{i}H,\,\,\,\left\{\begin{array}{r}s^{-1}gs\in H$ iff $g_{i}^{-1}gg_{i}\in H\\f(s^{-1}gs)=f(g_{i}^{-1}gg_{i})\end{array}\right.$.  So the second stmt follows. The first stmt follows by the first one. 
\\$\rightarrow$ if $f\in\mathcal{C}(G,k),\,\,\,\,\text{Res}_{H}^{G}f$ for $f|_{H}$. Then Res$_{H}^{G}\chi_{V}=\chi_{\text{Res}_{H}^{G}V}$. 
\begin{theorem}If $\left\{\begin{array}{r}f_{1}\in\mathcal{C}(H,k)\\f_{2}\in\mathcal{C}(G,k)\end{array}\right.,$ then $<f_{1},\text{Res}_{H}^{G}f_{2}>_{H}=<\text{Ind}_{H}^{G}f_{1},f_{2}>_{G}$.\end{theorem}
pf. $<f_{1},\text{Res}_{H}^{G}f_{2}>_{H}=\frac{1}{|H|}\sum_{h\in H}f_{1}(h)f_{2}(h)\\<\text{Ind}_{H}^{G}f_{1},f_{2}>_{G}=\frac{1}{|G|}\sum_{g\in G}\frac{1}{|H|}\sum_{s\in G|_{s^{-1}gs\in H}}f_{1}(s^{-1}gs)f_{2}(g)=\frac{1}{|G|}\sum\frac{1}{|H|}\sum f_{1}(s^{-1}gs)f_{2}(s^{-1}gs)=\frac{1}{|G|}\frac{1}{|H|}\sum_{h\in H}\sum_{s\in G}f_{1}(h)f_{2}(h)=\frac{1}{|H|}\sum_{h\in H}f_{1}(h)f_{2}(h)$. They are equivalent. 
\\\\\indent Let $R$ be a ring, $G$ be a finite group and $H,K$ be subgroups of $G$. Let $(V,\rho)$ be a representation of $H$ and $g_{1},\dots,g_{r}$ be a system of representatives of the double classes in $K\setminus G/H$, i.e. $G=\coprod_{i=1}^{r}Kg_{i}H$. For every $i\in\{1,\dots,r\},$ let $H_{i}=g_{i}Hg_{i}^{-1}\cap K$ and let $(V_{i},\rho_{i})$ be the representation of $H_{i}$ on $V$, given by $\rho_{i}(h)=\rho(g_{i}^{-1}hg_{i})$.
\begin{theorem}[Mackey's Formula]We have an isomorphism of $R[K]$-modules\begin{equation*}\text{Res}_{K}^{G}\text{Ind}_{H}^{G}V\cong\bigoplus_{i=1}^{r}\text{Ind}_{H_{i}}^{K}V_{i}\end{equation*}\end{theorem}
pf. $\forall j\in\{1,\dots,r\}$, let $(x_{i})_{i\in I_{j}}$ be a system of representatives of $K/H_{j}$. Then $\{x_{i}g_{j},1\leq j\leq r,i\in I_{j}\}$ is a system of representatives of $G/H\,\,\, ( G=\coprod_{j=1}^{r}Kg_{j}H=\coprod_{j=1}^{r}\coprod_{i\in I_{j}}x_{i}H_{j}g_{j}H=\coprod\coprod x_{i}g_{j}(g_{j}^{-1}H_{j}g_{j})H=\coprod\coprod x_{i}g_{j}H$ ). So $W\equiv \text{Ind}_{H}^{G}V=\bigoplus_{j=1}^{r}W_{j},\,\,\,\,\,\,W_{j}=\bigoplus_{i\in I_{j}}x_{i}g_{j}k[H]\otimes_{k[H]}V$. Fix $j$ and consider the $R$-linear maps $\begin{array}{rrr}\phi:\text{Ind}_{H_{j}}^{K}V_{j}&\rightarrow&W_{j}\\\sum_{i\in I_{j}}x_{i}\otimes v_{i}&\mapsto&\sum(x_{i}g_{j})\otimes v_{i}\end{array}$ and $\begin{array}{rrr}\psi:W_{j}&\rightarrow&\text{Ind}_{H_{j}}^{K}V_{j}\\\sum_{i\in I_{j}}(x_{i}g_{j})\otimes v_{i}&\mapsto&\sum_{i\in I_{j}}x_{i}\otimes v_{i}\end{array}$. Then $\phi,\psi$ are inverses to each other. Since $\phi$ is $K$-linear ($\forall y\in K,\exists \sigma\in\mathcal{G}_{I_{j}},h_{i}\in H_{j},i\in I_{j}$ s.t. $yx_{i}=x_{\sigma(i)}h_{i},\,\,\forall i\in I_{j},\,\,\,\,\phi(y(\sum x_{i}\otimes v_{i}))=\phi(\sum(yx_{i})\otimes v_{i})=\phi(\sum(x_{\sigma(i)}h_{i})\otimes v_{i})=\phi(\sum x_{\sigma(i)}\otimes(\rho_{j}(h_{i})v_{i}))=\phi(\sum x_{\sigma(i)}\otimes(\rho(g_{j}^{-1}h_{i}g_{j})v_{i}))=\sum(x_{\sigma(i)}g_{j})\otimes(\rho(g_{j}^{-1}h_{i}g_{j})v_{i})=\sum(x_{\sigma(i)}h_{i}g_{j})\otimes v_{i}=\sum(yx_{i}g_{j})\otimes v_{i}=y\phi(\sum x_{i}\otimes v_{i})$ ), the theorem follows. 
\begin{equation*}<V,W>_{G}=<\chi_{V},\chi_{W}>_{G}=\dim_{k}\hom_{G}(V,W)\end{equation*}
Let $H$ be a subgroup of $G$ and $(V,\rho)$ be a representation of $H$. For every $g\in G,\,\,\,\,H_{g}=gHg^{-1}\cap H$ and define two representations $\left\{\begin{array}{r}\rho_{g}(h)=\rho(g^{-1}hg)\\\text{Res}_{g}(\rho)(h)=\rho(h)\end{array}\right.$.
\begin{theorem}
TFAE
\begin{enumerate}\item $W\equiv\text{Ind}_{H}^{G}V$ is irreducible\item $V$ is irreducible, and $\forall g\in G\setminus H,\,\,\,\,\hom_{H_{g}}(\rho^{g},\text{Res}_{g}(\rho))=0$\end{enumerate}
\end{theorem}
pf.$\begin{array}{rr}<W,W>_{G}=<V,\text{Res}_{H}^{G}W>_{H}&$ by Frobenius reciprocity$\\=<V,\bigoplus_{g\in H\setminus G/H}\text{Ind}_{H_{g}}^{H}(\rho^{g})>_{H}&$ Mackey's formula$ \\=\sum_{g\in H\setminus G/H}<V,\text{Ind}_{H_{g}}^{H}(\rho^{g})>_{H}&=\sum<\text{Ind}_{H_{g}}^{H}(\rho^{g}),V>_{H}\\=\sum<\rho^{g},\text{Res}_{g}(\rho)>_{H}&$ by Frobenius reciprocity $\end{array}$ where all the summands are $\geq0$. If $g\in H,\,\,\,\,H_{g}=H$ and $\rho_{G}\cong\text{Res}_{g}(\rho)=\rho$, so $<\rho^{g},\text{Res}_{g}(V)>_{H}=<V,V>_{H}$. Then $W$ is irreducible $\Leftrightarrow <W,W>_{G}=1\Leftrightarrow <V,V>_{H}=1\\<\rho^{g},\text{Res}_{g}(\rho)>_{H}=0,\forall g\in G\setminus H\Leftrightarrow V$ is irreducible and $\hom_{H_{g}}(\rho^{g},\text{Res}_{g}(\rho)>=0,\,\,\,\forall g\in G\setminus H$. 

\begin{proposition}Write $R(G)_{\mathbb{Q}}=R(G)\otimes_{\mathbb{Z}}\mathbb{Q}$. Then the $\mathbb{Z}$-linear map $<\cdot,\cdot>_{G}:R(G)\times R(G)\rightarrow\mathbb{Z}$ induces a $\mathbb{Q}$-linear isomorphism\begin{equation*}\begin{array}{rrr}R(G)_{\mathbb{Q}}&\rightarrow&R(G)_{\mathbb{Q}}^{*}\equiv\hom_{\mathbb{Q}}(R(G)_{\mathbb{Q}},\mathbb{Q})\\x&\mapsto&(y\mapsto<x,y>_{G})\end{array}\end{equation*} Let $H$ be a subgroup of $G$. If we use the isomorphism above to identify $R(G)_{\mathbb{Q}}$ and $R(H)_{\mathbb{Q}}$ with their duals, then the transpose of $\text{Ind}_{G}^{H}$ is Res$_{H}^{G}$.\end{proposition}
pf.\begin{itemize}\item Denote by $\left\{\begin{array}{rr}(e_{V})_{V\in S_{k}(G)}&$ the basis $([V])_{V\in S_{k}(G)}$ of $R(G)_{\mathbb{Q}}\\(e_{V}^{*})_{V\in S_{k}(G)}&$ the dual basis $\end{array}\right.$. Then by the theorem A.2.4, and the corollary A.2.2., the map above sends each $e_{V}$ to a nonzero multiple of $e_{V}^{*}$ If $k$ is algebraically closed, the map sends each $e_{V}$ to $e_{V}^{*}$ itself by the corollary A.2.3., and it actually induces an isomorphism $R(G)\rightarrow\hom_{\mathbb{Z}}(R(G),\mathbb{Z})$.
\item by the Frobenius reciprocity


\end{itemize}
\begin{theorem}[Artin]Let $X$ be a set of subgroups of $G$. Consider the map $\text{Ind}_{X}=\bigoplus_{H\in X}\text{Ind}_{H}^{G}:\bigoplus_{H\in X}R(H)\rightarrow R(G)$. Then TFAE\begin{enumerate}\item $G=\cup_{H\in X}\cup_{g\in G}gHg^{-1}$\item Ind$_{X}\otimes_{\mathbb{Z}}\mathbb{Q}$ is surjective, that is $\forall x\in R(G),\,\,\,\exists d\geq1,\,\,\,\exists x_{H}\in R(H),\,\,\,\,H\in X$, st. $dx=\sum_{H\in X}\text{Ind}_{H}^{G}x_{H}$.\end{enumerate}\end{theorem}
pf.\begin{itemize}\item Let $S=\cup_{H\in X}\cup_{g\in G}gHg^{-1}$.Then if $H\in X,\,\,\,x_{H}\in R(H)_{\mathbb{Q}},\,\,\,\,\text{Ind}_{H}^{G}x_{H}$ is zero on $G\setminus S$. Then by 2, this implies $x=0$ on $G\setminus S,\,\,\,\,\forall x\in R(G)_{\mathbb{Q}}$. Then by A.2.5., $\forall f\in\mathcal{C}(G,k),\,\,\,f|_{G\setminus S}=0$. Hence $G\setminus S=\emptyset$.
\item By the previous proposition, the transpose of $\text{Ind}_{X}\otimes_{\mathbb{Z}}\mathbb{Q}$ is $\bigoplus_{H\in X}\text{Res}_{H}^{G}:R(G)_{\mathbb{Q}}\rightarrow\bigoplus_{H\in X}R(H)_{\mathbb{Q}}$. By 1, the sum of the restriction maps $\mathcal{C}(G,k)\rightarrow\bigoplus_{H\in X}\mathcal{C}(H,k)$ is injective, and therefore the map above is injective and 2 follows.  


\end{itemize}
\begin{corollary}Take for $X$ the set of cyclic subgroups of $G$. Then $\text{Ind}_{X}\otimes_{\mathbb{Z}}\mathbb{Q}$ is surjective.\end{corollary}
pf. Because $\forall g\in G$, it is an element of the cyclic subgroup that it generates, by the Artin's theorem, the corollary follows. 
\begin{definition}Let $p$ be a prime number. A finite group $H$ is called $p$-elementary if $H=C\times P$, with $C$ a cyclic group of order prime to $p$ and $P$ a $p$-group.\end{definition}
For every prime number $p$, let $X(p)$ be the set of $p$-elementary subgroups of $G$. Let $X$ be the union of all the $X(p)$ for $p$ prime.
\begin{theorem}Let $p$ be a prime number, write $|G|=p^{r}m$ with $(p,m)=1$, and let $V_{p}$ be the image of \begin{equation*}\text{Ind}_{p}\equiv\bigoplus_{H\in X(p)}\text{Ind}_{H}^{G}:\bigoplus_{H\in X(p)}R(H)\rightarrow R(G)\end{equation*} Then $m\in V_{p}$. In particular, $R(G)/V_{p}$ is a finite group of order prime to $p$.\end{theorem}
pf. Let $n=|G|$ and $\mathcal{O}=\mathbb{Z}[1,\zeta_{n},\dots,\zeta_{n}^{n-1}]\subset k$, where $\zeta_{n}$ is a primitive $n$th root of 1 in $k$. Then $\forall x\in R(G),\,\,\,x$ takes values in $\mathcal{O}$. Indeed for every representation $(V,\rho)$ of $G$ and $\forall g\in G,\,\,\,\rho(g)$ has eigenvalues in $\mathcal{O}$ by the previous proposition, so $\chi_{V}(g)\in\mathcal{O}$. We have $\mathcal{O}\cap\mathbb{Q}=\mathbb{Z}$ since $\forall \alpha\in \mathcal{O}\cap\mathbb{Q},\,\,\,f\in\mathbb{Q}[T]$ be its minimal polynomial over $\mathbb{Q}$, then $f\in\mathbb{Z}[T]$ as $\alpha$ is integral over $\mathbb{Z},\,\,\,\deg f=1$ as $\alpha\in\mathbb{Q}$. Then the map $\mathcal{O}/\mathbb{Z}\rightarrow(\mathcal{O}\otimes_{\mathbb{Z}}\mathbb{Q})/\mathbb{Q}$ is injective, so $\mathcal{O}/\mathbb{Z}$ is torsion free, so it is free as a $\mathbb{Z}$-module, so $\mathcal{O}$ has a $\mathbb{Z}$-basis of the form $(1,\alpha_{1},\dots,\alpha_{c})$. The image of $\mathcal{O}\otimes\text{Ind}_{p}:\bigoplus_{H\in X}\mathcal{O}\otimes_{\mathbb{Z}}R(H)\rightarrow\mathcal{O}\otimes_{\mathbb{Z}}R(G)$ is $\mathcal{O}\otimes_{\mathbb{Z}}V_{p}$ and we have $(\mathcal{O}\otimes_{\mathbb{Z}}V_{p})\cap R(G)=V_{p}$ since $\mathcal{O}\otimes_{\mathbb{Z}}V_{p}=V_{p}\oplus\bigoplus_{i=1}^{c}\alpha_{i}V_{p}\subset\mathcal{O}\otimes_{\mathbb{Z}}R(G)=R(G)\oplus\bigoplus\alpha_{i}R(G)$ and if $y=x+\sum\alpha_{i}x_{i}\in\mathcal{O}\otimes_{\mathbb{Z}}V_{p},\,\,\,x\in R(G)$ iff $\forall x_{i}=0$. Thus the theorem is reduced to the case when $m\in\mathcal{O}\otimes_{\mathbb{Z}}V_{p}$. 
\\If $C$ is a cyclic group of $G$ of order $c,\,\,\,\,\exists\begin{array}{rrr}\theta_{C}:C&\rightarrow&\mathbb{Z}\\x&\mapsto&\left\{\begin{array}{rr}c&C=<x>\\0&otherwise\end{array}\right.\end{array}$. Let $\theta_{C}'=\text{Ind}_{C}^{G}\theta_{C}$. Let $x\in G$. Then $\theta_{C}'(x)=\frac{1}{c}\sum_{s\in G|_{sxs^{-1}\in C}}\theta_{C}(sxs^{-1})=\sum_{s\in G|_{<sxs^{-1}>=C}}1.\,\,\,\,\forall s\in G,\,\,\,sxs^{-1}$ generates exactly one cyclic subgroup of $G$. Hence $\sum_{C\subset G,\,\,\,cyclic}\theta_{C}'(x)=\sum_{s\in G}1=|G|$. Let $C$ be any cyclic subgroup of $G$. 
\begin{itemize}\item \textbf{CLAIM : $\theta_{C}\in R(C)$}
\\pf. Let $c\equiv|C|$. If $c=1,\,\,\,\theta_{C}\in R(C)$. 
\\When $c>1,\,\,\,\,\,c=\sum_{B\subset C,\,\,\,cyclic}\text{Ind}_{B}^{C}\theta_{B}=\theta_{C}+\sum_{B\subsetneq C,\,\,\,cyclic}\text{Ind}_{B}^{C}\theta_{B}$. Since $c\in R(C)$ and $\forall \theta_{B}\in R(B),\,\,\,B\subsetneq C$ by the inductive hypothesis, so $\theta_{C}\in R(C)$. 
\end{itemize}
Let $f\in \mathcal{C}(G,\mathbb{Z})$ s.t. $f(G)\subset n\mathbb{Z}$. Then write $f=nf',\,\,\,f'\in\mathcal{C}(G,\mathbb{Z}).\,\,\,\,n=\sum_{C\subset G,\,\,cyclic}\text{Ind}_{C}^{G}\theta_{C}$, hence $f=\sum_{C\subset G,\,\,\,cyclic}\text{Ind}_{C}^{G}(\theta_{C})f'=\sum\text{Ind}_{C}^{G}(\theta_{C}\text{Res}_{C}^{G}f')$. Write $f_{C}=\theta_{C}\text{Res}_{C}^{G}f'$. Then $f_{C}\in\mathcal{C}(C,\mathbb{Z})$ and $f_{C}(C)\subset|C|\mathbb{Z}$. So $\forall \chi\in R(C),\,\,\,<f_{C},\chi>_{C}=\frac{1}{|C|}\sum_{x\in C}f_{C}(x)\chi(x)\in\mathcal{O}$ and therefore $f_{C}=\sum_{W\in S_{k}(C)}<f_{C},\chi_{W}>_{C}\chi_{W}\in\mathcal{O}\otimes_{\mathbb{Z}}R(C)$. So $f=\sum_{C\subset G,\,\,\,cyclic}\alpha_{c}\text{Ind}_{C}^{G}x_{C},$ where $\left\{\begin{array}{r}\alpha_{c}\in\mathcal{O}\\x_{C}\in R(C)\end{array}\right.$. 
\begin{definition}$x\in G$ is called $p$-unipotent if its order is a power of $p$. $x\in G$ is called $p$-regular if its order is prime to $p$. \end{definition}
\begin{itemize}\item\textbf{CLAIM : $\forall x\in G,\,\,\,\exists (x_{r},x_{u})$ where $x_{r},x_{u}\in G$ s.t. $\\\left\{\begin{array}{r}x_{r}$ is $p$-regular, and $x_{u}$ is $p$-unipotent$\\x=x_{r}x_{u}=x_{u}x_{r}\end{array}\right.$.}
\\pf. If $x_{r},x_{u}$ satisfies the conditions above, they have to be powers of $x$, since if $\left\{\begin{array}{r}a=|x_{r}|\\b=|x_{u}|\end{array}\right.,\,\,\,(a,b)=1$ so there exists an integer $N\geq1$ s.t. $a|N$ and $N\equiv1\mod b$ and then $\left\{\begin{array}{r}x^{N}=x_{r}^{N}x_{u}^{N}=x_{u}\\x^{1-N}=xx_{u}^{-1}=x_{r}\end{array}\right.$. Suppose we have two such pairs $(x_{r},x_{u}),(x_{r}',x_{u}')$ satisfying the above conditions. Then $\exists N\geq1$ s.t. $x_{u}^{p^{N}}=(x_{u}')^{p^{N}}=1,\,\,\,x_{r}^{p^{N}}=x_{r},\,\,\,(x_{r}')^{p^{N}}=x_{r}'$. Then $x^{p^{N}}=x_{r}=x_{r}'$ and therefore $x_{u}=x_{u}'$, so the uniqueness follows. Assume $G$ is generated by $x$, hence $G$ is cyclic. So we may assume $G=\mathbb{Z}/n\mathbb{Z}$ and $x=1$. Write $n=p^{r}m$ with $p\not|m$. Then by CRT, $G\cong\mathbb{Z}/p^{r}\mathbb{Z}\times\mathbb{Z}/m\mathbb{Z}$ and we can take $x_{r}=(0,1)$ and $x_{u}=(1,0)$ so such $x_{r},x_{u}$ exist. 
\item\textbf{CLAIM : Let $\chi\in\mathcal{O}\otimes_{\mathbb{Z}}R(G)$ be s.t. $\chi(G)\subset\mathbb{Z}$, and $x\in G$. Write $x=x_{r}x_{u}$. Then $\chi(x)=\chi(x_{r})\mod p$.}
\\pf. Assume $G=<x>$. Let $\chi_{1},\dots,\chi_{n}$ be the characters of the irreducible representations of $G$ over $k$, which are all 1-dimensional. Write $\chi=\sum_{i=1}^{n}a_{i}\chi_{i}$ with $a_{i}\in\mathcal{O}$. Let $q=p^{r}$ be the order of $x_{i}$. Then $x^{q}=x_{r}^{q}$ so $\forall i,\,\,\,\chi_{i}(x)^{q}=\chi_{i}(x_{r})^{q}$. So $\chi(x)^{q}=(\sum a_{i}\chi_{i}(x))^{q}=\sum a_{i}^{q}\chi_{i}(x)^{q}=\sum a_{i}^{q}\chi_{i}(x_{r})^{q}=\chi(x_{r})^{q}\mod p\mathcal{O}$. As $\left\{\begin{array}{r}\chi(x),\chi(x_{r})\in\mathbb{Z}\\p\mathcal{O}\cap\mathbb{Z}=p\mathbb{Z}\end{array}\right.,\,\,\,\,\chi(x)^{q}=\chi(x_{r})^{q}\mod p$. Also $a^{p}\equiv a\mod p,\,\,\,\,\forall a\in\mathbb{Z}$ implies $\chi(x)\equiv\chi(x_{r})\mod p$.
\item\textbf{CLAIM : Let $x\in G$ be $p$-regular. Let $H=C\times P$ be $p$-elementary with $C=<x>$ and $P$ a Sylow $p$-subgroup of $Z_{G}(x)$. Then $\exists \psi\in\mathcal{O}\otimes_{\mathbb{Z}}R(H)$ s.t. $\psi(H)\subset\mathbb{Z}$ and if $\psi'=\text{Ind}_{H}^{G}\psi,\,\,\,\,\\\left\{\begin{array}{r}\psi'(x)\neq0\mod p\\\psi'(s)=0,\,\,,\forall s\in G,$ a $p$-regular element that is not conjugate to $x\end{array}\right.$.}
\\pf. Let $\left\{\begin{array}{r}c=|C|\\p^{r}=|P|\end{array}\right.$. Let $\begin{array}{rrr}\psi_{C}:C&\rightarrow&\mathbb{Z}\\y&\mapsto&\left\{\begin{array}{rr}c&y=x\\0&otherwise\end{array}\right.\end{array}$. As $\psi_{C}(C)\subset c\mathbb{Z},\,\,\,\psi_{C}\in\mathcal{O}\otimes_{\mathbb{Z}}R(C)$ by the previous claim. Let $\begin{array}{rrr}\psi:H=C\times P&\rightarrow&\mathbb{Z}\\(x,y)&\mapsto&\psi_{C}(x)\end{array}$. Then $\psi\in\mathcal{O}\otimes_{\mathbb{Z}}R(H)$. Let $s\in G$ be $p$-regular. Then $\psi'(x)=\frac{1}{cp^{r}}\sum_{y\in G|_{ysy^{-1}\in H}}\psi(ysy^{-1})$. Let $y\in G$. If $ysy^{-1}\in H,\,\,\,ysy^{-1}\in C$ so $\psi(ysy^{-1})\neq0$ iff $ysy^{-1}=x$. Hence $\psi'(s)=0 $ if $s $ is not conjugate to $x$. Also $\psi'(x)=\frac{1}{cp^{r}}\sum_{y\in G|_{yxy^{-1}=x}}\psi(x)=\frac{1}{p^{r}}\sum_{y\in Z_{G}(x)}1=\frac{1}{p^{r}}|Z_{G}(x)|\neq0\mod p$. 

\item\textbf{$\exists\psi\in\mathcal{O}\otimes_{\mathbb{Z}}V_{p}$ s.t. $\psi(G)\subset\mathbb{Z}$ and $\psi(x)\neq0\mod p,\,\,\,\forall x\in G$.}


\end{itemize}
Write $n\equiv|G|=p^{r}m$ with $(p,m)=1$. Choose a $\psi\in\mathcal{O}\otimes_{\mathbb{Z}}V_{p}$ as previously. Let $N=\psi(p^{r})=|(\mathbb{Z}/p^{r}\mathbb{Z})^{\times}|$. Then $\forall l\in\mathbb{Z}$ prime to $p,\,\,\,l^{N}=1\mod p^{r}$. So $\forall x\in G,\,\,\,\psi(x)^{N}=1\mod p^{r}$, so $m(\psi^{N}-1)\in\mathcal{C}(G,\mathbb{Z})$ takes its values in $n\mathbb{Z}$. Then this implies $m(\psi^{N}-1)\in\mathcal{O}\otimes_{\mathbb{Z}}V_{p}$. AS it is an ideal of $\mathcal{O}\otimes_{\mathbb{Z}}R(G),\,\,\,m\psi^{N}\in\mathcal{O}\otimes_{\mathbb{Z}}V_{p}$. Then $m=m\psi^{N}-m(\psi^{N}-1)\in\mathcal{O}\otimes_{\mathbb{Z}}V_{p}$ and the theorem follows. 
\begin{corollary}For every representation $V$ of $G$, there exist monomial representations $V_{1},\dots,V_{r}$ of $G$ and integers $n_{1},\dots,n_{r}\in\mathbb{Z}$ s.t. in $R(G),\,\,\,[V]=\sum_{i=1}^{r}n_{i}[V_{i}]$.\end{corollary}


\begin{theorem}[Brauer]Let \begin{equation*}\text{Ind}_{X}=\bigoplus_{H\in X}\text{Ind}_{H}^{G}:\bigoplus_{H\in X}R(H)\rightarrow R(G)\end{equation*} Then it is surjective.\end{theorem}
pf. Im(Ind$_{X})=\sum_{p\,\,\text{prime}}V_{p}$ by the previous theorem, so $R(G)/\text{Im(Ind}_{X})$ is a finite group of order prime to every prime number, i.e. the trivial group. 
\begin{definition}A representation $V$ of $G$ is called monomial if there exists a subgroup $H$ of $G$ and a 1-dimensional representation $W$ of $H$ s.t. $V\cong\text{Ind}_{H}^{G}W$.\end{definition}
\begin{lemma}Let $P$ be a $p$-group and suppose $P$ is not abelian. Denote by $Z(P)$ the center of $P$. Then there exists an abelian normal subgroup $A$ of $P$ s.t. $Z(P)\subsetneq A$.\end{lemma}
pf. The quotient $P/Z(P)$ is a nontrivial $p$-group, so its center is nontrivial. Choose $A'\subset Z(P/Z(P))$ cyclic of order $p$, and let $A$ be its inverse image in $P$. Clearly $Z(P)\subsetneq A$ and $A$ is normal in $P$ because it is the inverse image of a normal subgroup of $P/Z(P)$. Also $A$ is abelian because it is generated by $Z(P)$ and by a lift of a generator of $A'$. 
\begin{proposition}Let $H$ be a $p$-elementary group. Then every irreducible representation of $H$ is monomial.\end{proposition}
pf. Write $H=C\times P$ with $C$ cyclic of order prime to $p$ and $P$ a $p$-group. Then irreducible representations of $H$ are all of the form $V_{1}\otimes_{k}V_{2}$ by the theorem A.2.6. where $V_{1},V_{2}$ is an irreducible representation of $C, P$ respectively. Then $V_{1}$ is 1-dimensional, so we just need to show $V_{2}$ is monomial. Assume $H=P$ is a $p$-group. If $P$ is abelian, every irreducible representation of $P$ is 1-dimensional so assume $P$ is not abelian. Let $(V,\rho)$ be an irreducible representation of $P$. If $\ker\rho\neq\{0\}$, by the inductive hypothesis to $P/\ker\rho,\,\,\,\rho$ is monomial. So assume $\rho$ is faithful. Then by the previous lemma, there exists a normal abelian subgroup $A$ of $P$ s.t. $Z(P)\subsetneq A$ where $Z(P)$ is the center of $P$. Let $V=V_{1}\oplus\dots\oplus V_{n}$ be the isotypic decomposition of Res$_{A}^{P}V$. Because $A$ is abelian, $A$ acts on each $V_{i}$ through a morphism of groups $\rho_{i}:A\rightarrow k^{\times}$. We cannot have $V_{1}=V$ because otherwise $\rho(A)$ would be contained in the center of $\rho(P)$, so $A$ would be contained in the center of $P$, a contradiction. Let $g\in P,\,\,\,i\in\{1,\dots,n\}$. Then if $v\in V_{i},\,\,\,y\in A,\,\,\,\,\,\,\rho(y)\rho(g)v=\rho(g)\rho(g^{-1}yg)v=\rho_{i}(g^{-1}yg)\rho(g)v$, so $\rho(g)$ sends $V_{i}$ bijectively to the isotypic component of Res$_{A}^{P}V$ corresponding to the map $\begin{array}{rrr}A&\rightarrow&k^{\times}\\y&\mapsto&\rho_{i}(g^{-1}yg)\end{array}$, i.e. the action of $P$ on $V$ permutes the $V_{i}$, so we get an action of $P$ on the set $\{V_{1},\dots,V_{n}\}$. As $V$ is irreducible, this action is transitive. Hence all the $V_{i}$ are isomorphic as $k$-vector space and so $\dim_{k}V_{1}=\dots=\dim_{k}V_{n}=\frac{1}{n}\dim_{k}V$. Let $Q=\{g\in P|\rho(g)V_{1}=V_{1}\}$, then $Q$ is a subgroup of $P$ and $|P/Q|=n>1$. So $Q\neq P$. As $Q$ stabilises $V_{1}$, we get a representation of $Q$ on $V_{1},$ denoted by $V_{Q}$. Define $\begin{array}{rrr}\phi:\text{Ind}_{Q}^{P}V_{Q}&\rightarrow&V\\g\otimes v&\mapsto&\rho(g)v\end{array}$, which is well-defined by the definition of $Q$. It is $P$-equivariant and surjective given $V=\sum_{g\in P}\rho(g)V_{1}$. Moreover, $\dim_{k}\text{Ind}_{Q}^{P}V_{Q}=|P/Q|\dim_{k}V_{Q}=n\dim_{k}V_{1}=\dim_{k}V$ so $\phi$ is an isomorphism and $V\cong\text{Ind}_{Q}^{P}V_{Q}$ as representations of $P$. Also since $V$ is irreducible, so is $V_{Q}$. By the inductive hypothesis, $V$ is monomial. 


\subsection{Characters in char$(k)=p$}


Let $R$ be a ring.
\begin{definition}If $M$ is a $R$-module, we write $rad(M)=rad(R)M$. This is a submodule of $M$.\end{definition}
$\rightarrow $ for $R$ left-Artinian, $R$-module $M$ is semisimple iff $rad(M)=0$
\\pf. If $rad(M)=0,\,\,\,\,M$ is $R/rad(R)$-module, so it is semisimple because $R/rad(R)$ is a semisimple ring.
\\If $M$ is semisimple $R$-module, $M=\bigoplus_{i\in I}M_{i}$ with all the $M_{i}$ simple. On each $M_{i},\,\,\,R$ acts through $R/rad(R)$ so $rad(R)M=0$.
\begin{definition}A $R$-module $M$ is called indecomposable if for every direct sum decomposition $M=M'\oplus M''$, we have either $M'=0$ or $M''=0$.\end{definition}
\begin{definition}A ring $S$ is called local if $S\neq\{0\}$ and $S$ has a unique maximal left ideal.\end{definition}
\begin{theorem}Let $S$ be a ring. TFAE\begin{enumerate}
\item $S$ is local\item $S$ has a unique maximal right ideal\item $rad(S)$ is a maximal left ideal of $S$\item $rad(S)$ is a maximal right ideal of $S$\item $S\neq\{0\}$ and for every $x\in S,$ either $x$ or $1-x$ is invertible\item$S/rad(S)$ is a division algebra\item $S\neq\{0\}$ and every $x\in S\setminus rad(S)$ is invertible\end{enumerate}
\end{theorem}
pf. \begin{itemize}
\item $i\leftrightarrow iii$ follows from the definition of $rad(S) 
\\ii\leftrightarrow iv$ follows since $rad(R)=rad(R^{op})$. 
\\$vii\leftrightarrow v$ follows since $x\in rad(R)$ and $\forall y,z,\in R,\,\,\,1-yxz$ is invertible are equivalent, when $R,R/rad(R)$ have the same simple modules. 
\\$vii\rightarrow vi$ trivially
\item $iii\rightarrow vii$
\\\\pf. Let $x\in S\setminus rad(S)$. Then $Sx\not\subset rad(S)$. As $rad(S)$ is a maximal left ideal of $S,\,\,\,\,\,Sx=S$ so $\exists y\in S$ s.t. $yx=1$. If $xy\in rad(S),\,\,\,1-yxy=0$ is invertible, a contradiction. Thus $xy\notin rad(S)$. Then $\exists z\in S$ s.t. $zxy=1$ so $y$ is a left and right invertible, hence invertible and $x=y^{-1}$ is also invertible.
\item $iv\rightarrow vii$ follows by the previous argument
\item $v\rightarrow iii,\,\,\,v\rightarrow iv$.
\\\\pf. Let $x\in S\setminus rad(S)$. Then $\exists y,z\in S$ s.t. $1-yxz$ is not invertible. Then by $v$, this implies $yxz$ is invertible, i.e. $x=y^{-1}z^{-1}$ is also invertible. So any left/right ideal of $S$ that strictly contains $rad(S)$ contains an invertible element, hence is equal to $S$.
\item $vi\rightarrow iii$
\\\\pf. Let $x\in S\setminus rad(S)$. Then $\exists y\in S$ s.t. $\left\{\begin{array}{r}yx\in 1+rad(S)\\xy\in 1+rad(S)\end{array}\right.$. Then $xy,yx$ are invertible hence $x,y$ are invertible. So any left ideal strictly containing $rad(S)$ contains an invertible element, which gives $iii$.

\end{itemize}


\begin{proposition}Let $S$ be a left Artinian ring. Then $rad(S)^{n}=0$ for $n$ big enough. In fact, if $n=lg(S),\,\,\,\,rad(S)^{n}=0$. \end{proposition}
pf. Let $S=I_{0}\supset\dots\supset I_{n}=0$ be a Jordon-H\"older series for $S$. Then $I_{i}$ are left ideals of $S$ and $I_{i}/I_{i+1}$ is a simple $S$-module $\forall i$. As $rad(S)$ annihilates every simple $S$-module. $rad(S)I_{i}\subset I_{i+1},\,\,\,\,\forall i$ and therefore $rad(S)^{n}=rad(S)^{n}I_{0}\subset I_{n}=0$. 
\begin{corollary}Let $S$ be a left Artinian ring. Then TFAE\begin{enumerate}\item $S$ is local\item every element of $S$ is nilpotent or invertible\end{enumerate}\end{corollary}
pf.\begin{itemize}\item $i\rightarrow ii$
\\\\pf. Let $x\in S$ and suppose $x$ is not nilpotent. Then $x\notin rad(S)$ by the previous proposition, so $x$ is invertible by the previous theorem
\item $ii\rightarrow i$
\\\\pf. If $x\in S$ is nilpotent, $\sum_{n\geq0}x^{n}$ is finite, hence defines an element of $S$, and this element is the inverse of $1-x$. So $ii$ implies $x$ or $1-x$ is invertible $\forall x\in S$, and by the previous theorem $S$ is local.\end{itemize}
\begin{proposition}[Fitting's Lemma]Let $R$ be a ring and $M$ be a $R$-module of finite length. Let $f\in\text{End}_{R}(M)$. Then for $n$ big enough, \begin{equation*}M=\ker f^{n}\oplus\text{Im}f^{n}\end{equation*}\end{proposition}
pf. As $M$ has finite length, every non-increasing or non-decreasing sequence of $R$-submodules of $M$ has to stabilise. Then for the non-decreasing sequence $\ker f^{n}$ and the non-increasing sequence Im$f^{n}$, we find an integer $N\geq0$ s.t. $\forall n\geq N,\,\,\,\left\{\begin{array}{r}\ker f^{n}=\ker f^{n+1}\\\text{Im}f^{n}=\text{Im}f^{n+1}\end{array}\right.$. Let $n\geq N,$ if $x\in M,\,\,\,f^{n}(x)\in \text{Im}f^{2n}$, so there exists $y\in M$ s.t. $f^{n}(x)=f^{2n}(y)$ and so $x-f^{n}(x)\in\ker f^{n}$ and $x=f^{n}(x)+(x-f^{n}(x))\in\text{Im}f^{n}+\ker f^{n}$. So $M=\text{Im}f^{n}+\ker f^{n}$. If $\exists 0\neq x\in \ker f^{n}\cap \text{Im}f^{n},\,\,\,x=f^{n}(y),\,\,\,\exists y\in M$. Then $y\in\ker f^{2n}=\ker f^{n}$ so $x=f^{n}(y)=0$. So $M=\ker f^{n}\oplus \text{Im}f^{n}$.
\begin{corollary}Let $R,M$ be as in the Fitting's lemma. Then $M$ is indecomposable iff $\text{End}_{R}(M)$ is local.\end{corollary}
pf. Suppose $M$ is indecomposable, and let $f\in \text{End}_{R}(M)$. Then there exists an integer $n\geq1$ s.t. $M=\ker f^{n}\oplus \text{Im}f^{n}$. As $M$ is indecomposable, either $\\\left\{\begin{array}{rr}\text{Im}f^{n}=0 & f$ is nilpotent  $\\\ker f^{n}=0& f^{n}$ is invertible$\end{array}\right.$. 
\\Suppose $M$ is not indecomposable and write $M=M'\oplus M''$ with $M',M''\neq0$. Let $\pi:M\rightarrow M'$ be the projection with kernel $M''$. Then $\pi$ is not invertible because $\ker\pi=M''\neq0$ and $1-\pi$ is not invertible because $\ker(1-\pi)=M'\neq0$ so $\text{End}_{R}M$ is not local.
\begin{theorem}[Krull-Schmidt-Remak Theorem]Let $M$ be a $R$-module of finite length. Then $M=M_{1}\oplus\dots\oplus M_{r}$ with the $M_{i}$ indecomposable. Moreover, the $M_{i}$ are uniquely determined up to reordering.\end{theorem}
pf.\begin{itemize}\item existence
\\\\pf. if $lg(M)=1,\,\,\,\,M$ is simple and the existence follows trivially.
\\if $lg(M)\geq2,$ and $M$ is not indecomposable, write $M=M'\oplus M''$ with $M',M''\neq0$. Then $lg(M'),lg(M'')<lg(M)$ and therefore by the inductive hypothesis, we get the decompositions for $M',M''$ and therefore for $M$.
\item uniqueness
\\\\pf. Assume the contrary. Then $M=M_{1}\oplus\dots\oplus M_{r}=M_{1}'\oplus\dots\oplus M_{s}'$ with all the $M_{i},M_{j}'$ are indecomposable. If $r=1,\,\,\,\,M_{1}=M_{1}'$ with $s=1$ is obvious. Suppose $r\geq2.$ Then $\forall j\in\{1,\dots,s\},$ let $u_{j}\in\text{End}_{R}(M_{1})$ be the composition $M_{1}\hookrightarrow M\rightarrow M_{j}'\hookrightarrow M\rightarrow M_{1}$, where all maps are obvious injections or projections. Then $id_{M_{1}}=\sum_{j=1}^{s}u_{j}$, so $\exists j$ s.t. $u_{j}\notin rad(\text{End}_{R}(M_{1}))$. Wlog assume $j=1$. As $M_{1}$ is indecomposable, End$_{R}(M_{1}$) is local by the previous corollary, so $u_{1}$ is invertible by the previous theorem. Write $u_{1}=vw$ where $w$ is the composition $M_{1}\hookrightarrow M\rightarrow M_{1}'$ and $v$ is the composition $M_{1}'\hookrightarrow M\rightarrow M_{1}$. Then $(u_{1}^{-1}v)w=id_{M_{1}}$ so $w$ is injective and $M_{1}'=w(M_{1})\oplus\ker(u_{1}^{-1}v)$. As $M_{1}'$ is indecomposable and $w(M_{1})\neq0,\,\,\,\,\ker(u_{1}^{-1}v)=0,$ so $w$ is an isomorphism, and so is $v$. Let $x\in M_{1}\cap(\bigoplus_{j\geq2}M_{j}')$. Then the projection of $x$ on $M_{1}'$ is 0 so $w(x)=0$ so $x=0$ given $w$ being injective. Hence $M_{1}$ and $\bigoplus_{j\geq2}M_{j}'$ are direct sums. Moreover if $x\in M$, then the surjectivity of $w$ implies $\exists x_{1}\in M_{1}$ s.t. $x-x_{1}\in\sum_{j\geq2}M_{j}'$ so $M=M_{1}\oplus\sum_{j\geq2}M_{j}'$. We get $M=M_{1}\oplus\sum_{j\geq2}M_{j}'$ and $M_{1}\cong M_{1}'$, by the inductive hypothesis, $M/M_{1}\cong\bigoplus_{j=2}^{r}M_{j}\cong\bigoplus_{j=2}^{s}M_{j}'$. 

\end{itemize}

\begin{proposition}Let $P$ be a projective $R$-module, $M$ be a $R$-module and $I$ be an ideal of $R$. Then the reduction modulo $I$ map\begin{equation*}\hom_{R}(P,M)\rightarrow\hom_{R}(P/IP,M/IM)=\hom_{R/I}(P/IP,M/IM)\end{equation*}is surjective.\end{proposition}
pf. Denote by $\left\{\begin{array}{r}\pi:P\rightarrow P/IP\\\pi':M\rightarrow M/IM\end{array}\right.$ the projections. Then $\forall u\in \hom_{R}(P/IP,M/IM),$ \begin{tikzcd}P\arrow[r,dotted,"u'"]\arrow[d,"\pi"]\arrow[dr,"u\pi"]&M\arrow[d,"\pi'"]\\P/IP\arrow[r,"u"]&M/IM\end{tikzcd}. As $\pi'$ is surjective and $P$ is projective, $\exists u':P\rightarrow M$ s.t. $\pi'u'=u\pi$. 

\begin{lemma}If $M$ is a nonzero $R$-module of finite length, $rad(M)\neq M$.\end{lemma}
pf. Let $M'\subsetneq M$ be a maximal proper submodule. Then $M/M'$ is simple so $rad(R)$ acts trivially on $M/M'$ so $rad(M)=rad(R)M\subset M'$. 


\begin{corollary}Assume $R$ is left-Artinian. Let $P$ be a projective $R$-module of finite length. Then $P$ is indecomposable iff $P/rad(P)$ is simple.\end{corollary}
pf. If $P$ is not indecomposable, $P=P_{1}\oplus P_{2},$ with $P_{1},P_{2}\neq0$. Then $P/rad(P)=P_{1}/rad(P_{1})\oplus P_{2}/rad(P_{2})$ and each factor is not zero by the previous lemma. So $P/rad(P)$ is not simple. 
\\Assume $M\equiv P/rad(P)$ is not simple. As $R$ is left-Artinian, $R/rad(R)$ is semisimple, so $M$ is a semisimple module and therefore $M=M_{1}\oplus M_{2}$ with $M_{1},M_{2}\neq0$. Let $\pi_{1}\in\text{End}_{R}(M_{1})$ be the composition $M\rightarrow M_{1}\hookrightarrow M$ where the maps are the obvious projection and inclusion. Then by the previous proposition, $\exists \pi\in \text{End}_{R}(P)$ s.t. $\pi\mod rad(P)=\pi_{1}$. Then neither $\pi$ nor $1-\pi$ are invertible so $\text{End}_{R}(P)$ is not local, so $P$ cannot be indecomposable by the corollary A.2.11.
\\\\\indent Write $PI(R)$ for the set of isomorphism classes of finite length projective indecomposable $R$-modules and $S(R)$ for the set of isomorphism classes of simple $R$-modules.
\begin{proposition}Suppose $R$ is left-Artinian and left-Noetherian. Then the map $P\mapsto P/rad(P)$ induces a bijection $PI(R)\rightarrow S(R)$.\end{proposition}
pf. \begin{itemize}\item Well-definedness follows by the previous corollary. 
\item Let $M$ be a simple $R$-module. Any $x\in M\setminus\{0\}$ gives a surjective map $\begin{array}{rrr}R&\rightarrow&M\\a&\mapsto&ax\end{array}$. As $R$ is a $R$-module of finite length, by the Krull-Schmidt-Remak Theorem, we can write $R=P_{1}\oplus\dots\oplus P_{n}$ where $P_{i}$ are indecomposable module that are automatically projective as direct summands of a free module. Then $R/rad(R)=\bigoplus P_{i}/rad(P_{i})$ surjects to $M$ and every $P_{i}/rad(P_{i})$ is simple by the previous corollary. So $M$ is isomorphic to one of $P_{i}/rad(P_{i})$ by the Schur's lemma. The map above is surjective.
\item Let $P,P'$ be two projective indecomposable modules of finite length, and suppose we have a $R$-module isomorphism $u:P/rad(P)\rightarrow P'/rad( P')$. Then by the previous proposition, there exists a $R$-module map $u':P\rightarrow P'$ s.t. $u=u'\mod rad(R)$. Let $N\subset P'$ be a proper maximal submodule. Then $P'/N$ is simple so $rad(R)(P'/N)=0$ so $rad(P')=rad(R)P'\subset N$. As $P'/rad(P')$ is simple, $N=rad(P')$ and $rad(P')$ is the unique proper maximal submodule of $P'$. As $u'(P)\not\subset rad(P'),\,\,\,u'(P)=P'$ so $u'$ is surjective. As $P'$ is projective, $P\cong P'\oplus\ker u'$ but since $P$ is indecomposable, $\ker u'=0$ and $u'$ is an isomorphism.

\end{itemize}


If $M$ is a simple $R$-module, its inverse image in $PI(R)$ is a minimal projective $R$-module s.t. $P$ surjects onto $M$. This is called a \textbf{projective envelope}  or \textbf{projective cover of }$M$. 

\begin{theorem}Let $S$ be a ring, and $I$ be an ideal of $S$ s.t. every element of $I$ is nilpotent. Let $\bar{e}\in S/I$ be idempotent. Then there exists $e\in S$ idempotent s.t. $\bar{e}=e\mod I$.\end{theorem}
pf. $\bar{f}=1-\bar{e}$ is also idempotent and $\bar{e}\bar{f}=\bar{f}\bar{e}=0$. Let $e$ be any lift of $\bar{e}$ and let $f=1-e$. Then $ef=fe\in I$ and $e+f=1\mod I$. Then by the assumption on $I,\,\,\,\,\exists k\geq1$ s.t. $(ef)^{k}=e^{k}f^{k}=0$. Here $e^{k}=\bar{e}^{k}=\bar{e}\mod I$. Let $\left\{\begin{array}{r}e'=e^{k}\\f'=f^{k}\end{array}\right.$. Then $e'f'=f'e'=0$ and $e'+f'=e^{k}+f^{k}=\bar{e}+\bar{f}=1\mod I$. Let $x=1-(e'+f')$. As $x\in I,\,\,\,\,\exists n\geq1$ s.t. $x^{n}=0$. So $u\equiv 1+x+\dots+x^{n-1}$ is an inverse of $1-x=e'+f'$ and it commutes with $e',f'$. Let $\left\{\begin{array}{r}e''=ue'\\f''=uf'\end{array}\right.$. Then $e''\equiv\bar{e}\mod I$ and $\left\{\begin{array}{r}e''f''=f''e''=0\\f''=uf'\end{array}\right.$. Then $e''=\bar{e}\mod I$ so $e''f''=f''e''=0$ and $e''+f''=u(e'+f')=1$. So $(e'')^{2}=(e'')^{2}+e''f''=e''(e''+f'')=e''$ and we have the idempotent lift of $\bar{e}$. 

\begin{corollary}Let $S$ be a ring and $I$ be an ideal of $S$. Suppose that the obvious map $S\rightarrow\hat{S}\equiv\lim S/I^{n}$ is an isomorphism. Then any idempotent of $S/I$ lifts to an idempotent of $S$.\end{corollary}
pf. Let $\bar{e}\in S/I$ be idempotent. Then by the previous theorem and induction on $n\geq1,$ we can construct a sequence of idempotents $e_{n}\in S/I^{n}$ s.t. $e_{1}=\bar{e}$ and $e_{n+1}$ lifts to $e_{n},\,\,\,\,\forall n$. Then $e\equiv(e_{n})\in\hat{S}$ and its preimage in $S$ is an idempotent lifting $\bar{e}$. 
\begin{corollary}\begin{enumerate}\item $PK(R)$ is a free $\mathbb{Z}$-module with basis $([P])_{P\in PI(R)}$.\item If $P,P'$ are projective $R$-modules of finite length, then $P\cong P'$ as $R$-modules iff $[P]=[P']$ in $PK(R)$.\end{enumerate}\end{corollary}
pf. Write $\left\{\begin{array}{r}P=P_{1}\oplus\dots\oplus P_{r}\\P'=P_{1}'\oplus\dots\oplus P_{s}'\end{array}\right.$ with $P_{i},P_{j}'$ indecomposable by the Krull-Schmidt-Remak theorem. Then $P_{i},P_{j}'$ are projective since it is a direct summand of projective modules, and therefore they are in $PI(R)$. Then by the same theorem, $P\cong P'$ iff $\exists \sigma:\{1,\dots,r\}\rightarrow\{1,\dots,s\}$ s.t. $P_{i}\cong P_{\sigma(i)}',\,\,\,\forall i$. Then by $i$, this is equivalent to $[P]=[P'].$
\\\\\indent Let $\Lambda$ be a commutative ring and $G$ be a finite group. Assume every $\Lambda[G]$-module is of finite type over $\Lambda$.
\begin{definition}We denote by $P_{\Lambda}(G)$ the quotient of the free $\mathbb{Z}$-module on all the $[P]$, for $P$ projective $\Lambda[G]$-module that is of finite type as a $\Lambda$-module, by the relations $[P]=[P']+[P'']$, for every exact sequence $0\rightarrow P'\rightarrow P\rightarrow P''\rightarrow0$.\end{definition}
$\rightarrow \Lambda\rightarrow\Lambda'$ a morphism of commutative rings, then $P\mapsto P\otimes_{\Lambda}\Lambda'$ induces a morphism of groups $P_{\Lambda}(G)\rightarrow P_{\Lambda'}(G)$ ( if $P$ is a projective $
\Lambda[G]$-module, it is a direct summand of some free $\Lambda[G]$-module $F$ and then $P\otimes_{\Lambda}\Lambda'$ is a direct summand of the free $\Lambda'[G]$-module $F\otimes_{\Lambda}\Lambda'$ and the morphism of commutative rings give $\Lambda'$ as $\Lambda$-modules ).
\begin{proposition}Let $P$ be a $\Lambda[G]$-module. Then TFAE\begin{enumerate}\item $P$ is a projective $\Lambda[G]$-module\item $P$ is projective as a $\Lambda$-module and there exists $u\in\text{End}_{\Lambda}(P)$ s.t. $\forall x\in P,\,\,\,\,\sum_{g\in G}gu(g^{-1}x)=x$.\end{enumerate}\end{proposition}
pf. Write $P_{0}$ for $P$ seen as a $\Lambda$-module. Let $Q=\Lambda[G]\otimes_{\Lambda}P_{0}$. Then there exists a surjective map $\begin{array}{rrr}q:Q&\rightarrow&P\\x\otimes y&\mapsto&xy\end{array}$.
\begin{itemize}
\item \textbf{CLAIM : $\begin{array}{rrr}\phi:\text{End}_{\Lambda}(P_{0})&\rightarrow&\hom_{\Lambda[G]}(P,Q)\\u&\mapsto&\sum_{g\in G}g\otimes ug^{-1}\end{array}$ is an well-defined isomorphism of $\Lambda$-modules.}
\\\\pf. $\phi(u)$ is $\Lambda[G]$-linear $\forall u\in \text{End}_{\Lambda}(P_{0})$ and therefore $\phi$ is well-defined.
\\Let $u\in \text{End}_{\Lambda}(P_{0})$. Here $Q=\bigoplus_{g\in G}g\otimes P_{0}$ as a $\Lambda$-module. So if $x\in P$ s.t. $0=\phi(u)(x)=\sum g\otimes u(g^{-1}x),\,\,\,\,u(g^{-1}x)=0,\,\,\,\forall g\in G$ and in particular $u(x)=0$. Hence $u=0$ if $\phi(u)=0$ and $\phi$ is injective. 
\\Let $v\in\hom_{\Lambda[G]}(P,Q)$. Then $v(x)=\sum_{g\in G}g\otimes u_{g}(x),\,\,\,\forall x\in P$ with $u_{g}\in \text{End}_{\Lambda}(P_{0})$. By $\Lambda[G]$-linearity of $v,\,\,\,\,\,\forall h\in G,\,\,\,\forall x\in P,\,\,\,v(h^{-1}x)=\sum g\otimes u_{g}(h^{-1}x)=\sum (h^{-1}g)\otimes u_{g}(x)$ so $u_{g}(x)=u_{h^{-1}g}(h^{-1}x),\,\,\,\,\forall h,g\in G,\,\,\,\forall x\in P$. In particular, $u_{g}(x)=u_{1}(g^{-1}x),\,\,\,\\\forall g\in G,\,\,\,\forall x\in P$ and therefore $v=\phi(u_{1})$ and $\phi$ is surjective.

\item $i\rightarrow ii$
\\\\pf. If $P$ is projective, it is a direct summand of some free module $\Lambda[G]^{(I)}$. As $\Lambda[G]$ is a free $\Lambda$-module, $\Lambda[G]^{(I)}$ is also a free $\Lambda$-module, so $P_{0}$ is projective. Also as $q:Q\rightarrow P$ is surjective and $\Lambda$-linear, $\exists \Lambda$-linear map $s:P\rightarrow Q$ s.t. $qs=id_{P}$. Write $s=\phi(u)$ with $u\in\text{End}_{\Lambda}(P_{0})$. Then $id_{P}=q\sum gug^{-1}$ which gives $ii$.

\item $ii\rightarrow i$
\\\\pf. If $P_{0}$ is projective $\Lambda$-module, $Q$ is projective $\Lambda[G]$-module. Also $s\equiv\phi(u):P\rightarrow Q$ satisfies $qs=id_{P}$, so $P$ is a direct summand of $Q$, hence is also a projective $\Lambda[G]$-module.

\end{itemize}
\begin{theorem}Suppose $\Lambda$ is a discrete valuation ring with residue field $k$ and maximal ideal $\mathbf{m}$. 
\begin{enumerate}\item if $P$ is a $\Lambda[G]$-module that is free of finite type over $\Lambda$, then $P$ is projective $\Lambda[G]$-module iff $\bar{P}\equiv P\otimes_{\Lambda}k$ is a projective $k[G]$-module.\item if $P,P'$ are projective $\Lambda[G]$-modules, then $P\cong P'$ as $\Lambda[G]$-module iff $P\otimes_{\Lambda}k\cong P'\otimes_{\Lambda}k$ as $k[G]$-modules.\item suppose the DVR $\Lambda$ is complete. If $\bar{P}$ is a projective $k[G]$-module, then there exists a unique projective $\Lambda[G]$-module $P$ s.t. $\bar{P}\cong P\otimes_{\Lambda}k$.\end{enumerate}\end{theorem}
pf.\begin{itemize}\item If $P$ is projective as a $\Lambda[G]$-module, then $\bar{P}$ is projective as a $k[G]$-module. 
\\Suppose $\bar{P}$ is a projective $k[G]$-module. Then by the previous proposition, $\exists \bar{u}\in \text{End}_{k}(\bar{P})$ s.t. $\sum g\bar{u}g^{-1}=id_{\bar{P}}$. As $P$ is a projective $\Lambda$-module, $\exists u\in\text{End}_{\Lambda}(P)$ lifting $\bar{u}$ by the proposition A.2.16. Then $u'\equiv\sum gug^{-1}\in\text{End}_{\Lambda[G]}(P)$ is equal to $id_{\bar{P}}$ modulo $\mathbf{m}$. So $\det(u')=1\mod\mathbf{m}$, so $\det(u')\in\Lambda^{\times}$ so $u'$ is invertible and $\sum g(u(u')^{-1})g^{-1}=id_{P}$. Then by the previous proposition, $P$ is projective $\Lambda[G]$-module.
\item Let $\bar{u}:P\otimes_{\Lambda}k\rightarrow P'\otimes_{\Lambda}k$ be an isomorphism of $k[G]$-modules. Then by the previous proposition A.2.26., there exists a $\Lambda[G]$-module map $u:P\rightarrow P'$ lifting $\bar{u}$. We know $P,P'$ are projective $\Lambda$-modules of finite type, hence they are free $\Lambda$-modules of finite type. Their ranks are equal, because they are equal to the dimension over $k$ of $P\otimes_{\Lambda}k\cong P'\otimes_{\Lambda}k$. If we choose $\Lambda$-bases of $P,P'$, then the matrix $A$ of $u$ is square, and $\det(A)\neq0\mod\mathbf{m}$ because $\bar{u}$ is invertible. Hence $\det(A)\in\Lambda^{\times}$ and $u$ is an isomorphism of $\Lambda$-modules.
\item The uniqueness follows from $ii$. For existence, write $\bar{F}=\bar{P}\oplus\bar{P}'$ with $\bar{F}=k[G]^{\oplus n}$, and let $F=\Lambda[G]^{\oplus n},$ and $B=\text{End}_{\Lambda[G]}(F)$. Then the map $B\rightarrow\text{End}_{k[G]}(\bar{F})$ is surjective by the proposition A.2.26., so it identifies $\text{End}_{k[G]}(\bar{F})$ with $B/\mathbf{m}B$. Let $\bar{e}\in\text{End}_{k[G]}(\bar{F})$ be the projection on $\bar{P}$ with kernel $\bar{P}'$. Then there exists an idempotent $e\in B$ lifting $\bar{e}$. Then $F=\text{Im}(e)\oplus \ker(e)$ and $P\equiv \text{Im}(e)$ is a projective $\Lambda[G]$-module s.t. $P\otimes_{\Lambda}k=\text{Im}(\bar{e})=\bar{P}$.




\end{itemize}
\indent Fix a complete DVR $\Lambda$ with uniformizer $\omega$, maximal ideal $\mathbf{m}=(\omega)$, residue field $k$ and fraction field $K$ and assume $char(K)=0$. \begin{tikzcd}P_{k}(G)\arrow[rr,"c"]\arrow[rd,"e"]&&R_{k}(G)\\&R_{K}(G)\arrow[ru,"d"]&\end{tikzcd} where $d$ is the decomposition \begin{tikzcd}P_{k}(G)\arrow[r,"\psi^{-1}"]&P_{\Lambda}(G)\arrow[r,"\otimes_{\Lambda}K"]&P_{K}(G)=R_{K}(G)\end{tikzcd} and $c$ is the obvious map. Let $V$ be a $K[G]$-module. Let $M_{1}\subset V$ be a $\Lambda$-lattice, i.e. a finite type $\Lambda$-submodule of $V$ s.t. $KM_{1}=V$. After replacing $M_{1}$ by $\sum_{g\in G}gM_{1}$, assume $M_{1}$ is stable by $G$. Then $\bar{M}_{1}\equiv M_{1}\otimes_{\Lambda}k$ is a $k[G]$-module. 
\begin{theorem}$[\bar{M}_{1}]\in R_{k}(G)$ only depends on $V$.\end{theorem}
pf. Let $M_{2}$ be another $G$-stable $\Lambda$-lattice of $V$. 
\begin{itemize}
\item $\omega M_{1}\subset M_{2}\subset M_{1}
\\\\$ let $N=M_{1}/M_{2}$, which is a $k[G]$-module because $\omega M_{1}\subset M_{2}$. From $\omega M_{2}\subset \omega M_{1}\subset M_{2}\subset M_{1},$ there exists an exact sequence of $k[G]$-modules $0\rightarrow N\rightarrow \bar{M}_{2}\equiv M_{2}\otimes_{\Lambda}k=M_{2}/\omega M_{2}\rightarrow\bar{M}_{1}\rightarrow N\rightarrow0$. So in $R_{k}(G),\,\,\,\,[N]-[\bar{M}_{2}]+[\bar{M}_{2}]-[N]=0\rightarrow [\bar{M}_{1}]=[\bar{M}_{2}]$
\item general case
\\\\ multiplying $M_{2}$ by a high enough power of $\omega,$ we may assume $M_{2}\subset M_{1}$. There also exists $n\geq1$ s.t. $\omega^{n}M_{1}\subset M_{2}$. Then when n=1, by the previous statement the theorem follows. If $n\geq2,$ let $M_{3}=\omega^{n-1}M_{1}+M_{2}$. Then $\omega^{n-1}M_{1}\subset M_{3}\subset M_{1}$ so $[M_{3}\otimes_{\Lambda}k]=[\bar{M}_{1}]$ by the inductive hypothesis. Also $\omega M_{3}\subset M_{2}\subset M_{3}$ so $[M_{3}\otimes_{\Lambda}k]=[M_{2}\otimes_{\Lambda}k]$ by the previous statement. Then the theorem holds for $n$.

\end{itemize}
\begin{theorem}Assume $p= char(k)$ and $p\not| |G|$.\begin{enumerate}\item every $k[G]$-module is semisimple.\item every $\Lambda[G]$-module that is projective as a $\Lambda$-module is projective as a $\Lambda[G]$-module.\item the map $d:R_{K}\rightarrow R_{k}(G)$ is an isomorphism and so are $c,e$. Also $d$ induces a bijection $S_{K}(G)\rightarrow S_{k}(G)$. \end{enumerate}\end{theorem}
pf.\begin{itemize}\item $i$ is trivial by the chapter 1 Maschke's theorem 
\item $ii$ follows from $i$ and the previous theorem A.2.17.
\item by $i,\,\,\,P_{k}(G)=R_{k}(G)$ and $c$ is just the identity morphism. Also $de=id_{R_{k}(G)}$ and $e([S_{k}(G)])\subset S_{K}(G)$. Let $V,V'$ be two $K[G]$-modules s.t. $d([V]-[V'])=0$, let $\left\{\begin{array}{r}M\subset V\\M'\subset V'\end{array}\right.$ be $G$-stable $\Lambda$-lattices and let $\left\{\begin{array}{r}\bar{M}=M\otimes_{\Lambda}k\\\bar{M}'=M'\otimes_{\Lambda}k\end{array}\right.$. Then $d([V]-[V'])=[\bar{M}]-[\bar{M}']=0\in R_{k}(G)$. As $p\not||G|,\,\,\,\bar{M}\cong\bar{M}'$ as $k[G]$-modules, so $M\cong M'$ as $\Lambda[G]$ by the previous theorem and therefore $V\cong V'$ and $[V]-[V']=0$.

\end{itemize}
\textbf{CLAIM : Let $G$ be a finite group and $p$ a prime number. If $k$ is algebraically closed field of characteristic $p$ and $G$ is a $p$-group and $V$ is an irreducible representation of $G$ on a finite dimensional $k$-vector space, $V$ is the trivial representation.}
\\pf.\begin{itemize}\item Suppose $G$ is abelian. Let $V$ be an irreducible representation of $V$. Then by Schur's lemma, $\text{End}_{k[G]}(V)=k$. As $G$ is abelian, $k[G]$ is commutative, so the action of any element of $k[G]$ on $V$ is in $\text{End}_{k[G]}(V)$, so there exists a $k$-algebra map $\begin{array}{rrr}u:k[G]&\rightarrow &k\\x&\mapsto&(v\mapsto u(x)v)\end{array}$. In particular, every $k$-subspace of $V$ is a subrepresentation; as $V$ is irreducible $\dim V=1$. Giving $u$ is the same as giving a morphism of groups $\rho:G\rightarrow k^{\times}$. But the only element of $k^{\times}$ that has order a power of $p$ is 1, so $\rho$ is trivial so $G$ acts trivially on $V$.
\item Assume $G$ is not abelian. Let $Z$ be the center of $G$ which is a nontrivial abelian $p$-group, so the inductive hypothesis applies to it. Let $V$ be an irreducible representation of $G$. Denote by $V^{Z}$ the $k$-subspace of $v\in V$ s.t. $gv=v,\,\,\,\,\forall g\in Z$. If $g\in G,\,\,\,v\in V^{Z}$, then $\forall h\in Z,\,\,\,h(gv)=(hg)v=(gh)v=g(hv)=gv\in V^{Z}$. Thus $V^{Z}$ is a subrepresentation of $V$. Let $V=V_{0}\supset\dots\supset V_{n}=0$ be a Jordan-H\"older series for $V$ seen as a $k[Z]$-module. Then $V_{n-1}$ is a simple $k[Z]$-module, so it is the trivial representation of $Z$ by the inductive hypothesis. This means $V_{n-1}\subset V^{Z}$ so $V^{Z}\neq0$. As $V$ is irreducible $V^{Z}=V$ so the action of $G$ on $V$ factors $G/Z$ and we can see $V$ as an irreducible representation of $G/Z$. By the inductive hypothesis, this is the trivial representation of $G/Z$ and so $V$ is the trivial representation of $G$.\end{itemize}

\begin{theorem}The maps \begin{equation*}\left\{\begin{array}{r}\text{Ind}\equiv\bigoplus_{l,\,\,\,prime}\bigoplus_{H\in X(l)}\text{Ind}_{H}^{G}:\bigoplus_{l\neq p,\,\,\,prime}\bigoplus_{H\in X(l)}R_{k}(H)\rightarrow R_{k}(G)\\\text{Ind}\equiv\bigoplus_{l,\,\,\,prime}\bigoplus_{H\in X(l)}\text{Ind}_{H}^{G}:\bigoplus_{l\neq p,\,\,\,prime}\bigoplus_{H\in X(l)}P_{k}(H)\rightarrow P_{k}(G)\end{array}\right.\end{equation*} are both surjective.\end{theorem}
pf. Let $\left\{\begin{array}{rr}\mathbb{1}_{K} &$ be the unit element in $R_{K}(G)\\\mathbb{1}_{k}&$ be the unit element in $R_{k}(G)\end{array}\right.$. Then $d(\mathbb{1}_{K})=\mathbb{1}_{k}$. By the Brauer's theorem, $\mathbb{1}_{K}=\sum_{l,\,\,\,prime}\sum_{H\in X(l)}\text{Ind}_{H}^{G}x_{H},\,\,\,\,\exists x_{H}\in R_{K}(H)$. Applying $d$ to this equality gives $\mathbb{1}_{k}=\sum_{l,\,\,\,prime}\sum_{H\in X(l)}\text{Ind}_{H}^{G}x_{H}'$ with $x_{H}'= d(x_{H})$ in $R_{k}(H)$. So if $y\in \left\{\begin{array}{r}R_{k}(G)\\P_{k}(G)\end{array}\right.,\,\,\,\,y=y\mathbb{1}_{k}=\sum_{l,\,\,\,prime}\sum_{H\in X(l)}\text{Ind}_{H}^{G}(x_{H}'\text{Res}_{H}^{G}y)$ where $x_{H}'\text{Res}_{H}^{G}(y)\in P_{k}(H)$ if $y\in P_{k}(G)$ given $\text{Res}_{H}^{G}$ sending $P_{k}(G)$ to $P_{k}(H).
\\\rightarrow $ given $k$ algebraically closed, when $l=p,$ for any $l$-elementary group $G=H\times P,$ and $P$ acts simple $k[G]$-module, so each is pulled back from $H$. Since $H$ is cyclic of order prime to $p$, it is $l'$-elementary for any prime $l'\mid|H|$; hence every class $\text{Ind}_{H}^{G}[M]$ already lies in the image without $l\neq p$, so domain can be taken as above. 
\begin{lemma}Suppose $G=P\times H$ with $P$ a $p$-group and $H$ of order prime to $p$. Then $P$ acts trivially on every semisimple $k[G]$-module.\end{lemma}
pf. Let $M$ be a simple $k[G]$-module. As $P$ is a $p$-group its only irreducible representation over $k$ is the trivial representation, so $M^{P}\neq0$. As $G=P\times H,$ the action of $G$ preserves $M^{p}$. As $M$ is simple, $M^{P}=M$ and so $P$ acts trivially on $M$.
\begin{corollary}If $k$ is algebraically closed and $K$ contains all the $|G|$th roots of the unity in $\bar{K}$, then $d:R_{K}(G)\rightarrow R_{k}(G)$ is surjective.\end{corollary}
pf. By the previous theorem, we may assume $G$ is $l$-elementary, for some prime $l$. Write $G=C\times G'$, with $C$ cyclic of order prime to $l$ and $G'$ a $l$-group. Let $\left\{\begin{array}{rr}l=p,&P=G',\,\,H=C\\l\neq p& P=C_{p},\,\,\,\,H=C^{p}\times G'\end{array}\right.$ where $C=C_{p}\times C^{p},\,\,\,C_{p}$ a $p$-group, and $C^{p}$ of order prime to $p$. Then $G=H\times P$ with $P$ a $p$-group and $H$ of order prime to $p$. Let $M$ be a simple $k[G]$-module. Then by the lemma, $P$ acts trivially on $M$. So $M$ is a simple $k[H]$-module. Then by the previous theorem A.2.17., there exists $K[H]$-module $V$ s.t. $d[V]=[M]$. We see $V$ as a $K[G]$-module by making $P$ act trivially, and then we still have $d[V]=[M]$.
\begin{proposition}If $G$ is a finite $p$-group, then $\left\{\begin{array}{r}P_{k}(G)\cong\mathbb{Z}\\R_{k}(G)\cong\mathbb{Z}\end{array}\right.$ and the map $P_{k}(G)\rightarrow R_{k}(G)$ corresponds to multiplication by $p^{n}\equiv|G|$.\end{proposition}
pf. The $\mathbb{Z}$-rank of $P_{k}(G)$ and $R_{k}(G)$ is 1, and there exists an isomorphism $\begin{array}{rrr}R_{k}(G)&\rightarrow&\mathbb{Z}\\M&\mapsto&\dim_{k}M\end{array}$ and $k[G]/rad(k[G])=k$. Thus $k[G]$ is local and every element of $k[G]$ is nilpotent or invertible. In particular, the only idempotents are 0 and 1. Let $k[G]\equiv M=M_{1}\oplus M_{2}$ with $\left\{\begin{array}{r}M_{1}=MI_{1}\\M_{2}=MI_{2}\end{array}\right.$ two left ideals $I_{1},I_{2}\subset k[G]$. Then we get two idempotents $e_{1},e_{2}\in k[G]$ s.t. $I_{1}=k[G]e_{1},I_{2}=k[G]e_{2}$, but there are only 0,1 for idempotents, this implies $k[G]$ is a projective indecomposable module, i.e. there is the unique projective indecomposable finite length $k[G]$-module that surjects to $\mathbb{1}$. Since in $R_{k}(G),\,\,\,\,\,[k[G]]=\dim_{k}(k[G])\mathbb{1}=p^{n}\mathbb{1}$, the proposition follows.
\begin{theorem}Assume $k$ is algebraically closed. Then $c:P_{k}(G)\rightarrow R_{k}(G)$ is injective and its image contains $p^{n}R_{k}(G)$ where $p^{n}$ is the biggest power of $p$ dividing $|G|$.\end{theorem}
pf. By the Brauer's theorem we wlog assume $G$ is $l$-elementary for some prime number $l$. Then we can write $G=H\times P$, with $P$ a $p$-group, $H$ of order prime to $p$. Then the trivial $k[H]$-module is projective given $k[H]$ being a semisimple ring. Thus with trivial action of $H$, $k[P]$ is a projective $k[G]$-module and by the previous proposition $k[P]=|P|\mathbb{1}\in R_{k}(G)$ is in the image of $c$. As Im$(c)$ is an ideal of $R_{k}(G),\,\,\,,\,p^{n}R_{k}(G)\subset \text{Im}(c)$. 
\\Since $R_{k}(G)/\text{Im}(c)$ is a torsion group and $P_{k}(G),R_{k}(G)$ are free $\mathbb{Z}$-modules of the same finite rank, $c$ needs to be injective. 
\begin{corollary}The map $e:P_{k}(G)\rightarrow R_{K}(G)$ is injective.\end{corollary}
pf. by the injectivity of $c$ from the previous theorem and the surjectivity of $d$ and commutative diagram above.
\begin{definition}An element $x\in G$ is called $p$-singular if $x$ is not $p$-regular, i.e. if $p$ divides the order of $x$.\end{definition}


\section{Compact Group Representation}

\begin{definition}A topological group is a topological space $G$ with a group structure s.t. the maps $\begin{array}{rrr}G&\rightarrow&G\\x&\mapsto&x^{-1}\end{array},\,\,\,\begin{array}{rrr}G\times G&\rightarrow&G\\(x,y)&\mapsto&xy\end{array}$ are both continuous.\end{definition}
Let $G$ be a Hausdorff locally compact topological group. By a measure on $G$, we mean a measure on the Borel $\sigma$-algebra of $B$. 
\begin{definition}A nonzero measure $\mu$ on $G$ is called a left Haar measure if for every measurable function $f:G\rightarrow\mathbb{C},\,\,\,\forall g\in G,$\begin{equation*}\int_{G}f(x)d\mu=\int_{G}f(gx)d\mu\end{equation*}\end{definition}
Denote by $\mathcal{C}$ the space of continuous functions with compact support $G\rightarrow\mathbb{C}$,equipped with the compact open topology ( generated by the sets $\{f\in\mathcal{C}|f(K)\subset U\}$ for $\left\{\begin{array}{rr}K\subset G&$ compact $\\U\subset\mathbb{C}&$ open $\end{array}\right.$ ). Let $\mathcal{C}^{*}$ be the space of continuous linear maps $\mathcal{C}\rightarrow\mathbb{C}$, equipped with the weak topology ( generated by $\{\lambda\in\mathcal{C}^{*}:|(\lambda-\lambda_{1})(f_{1})|<r_{1},\dots,|(\lambda-\lambda_{n})(f_{n})|<r_{n}\},\,\,\,\forall \lambda_{i}\in\mathcal{C}^{*},\,\,\,\forall f_{i}\in\mathcal{C},\,\,\,\forall r_{n}\in\mathbb{R}_{++}$ ). If $G$ is compact and metrisable, for every probability measure $\mu$ on $G$, the linear form $f\mapsto \int_{G}f(x)d\mu$ on $\mathcal{C}$ is continuous, hence in $\mathcal{C}^{*}$. This implies we can the set $\mathcal{P}(G)$ of probability measures on $G$ with a subset of $\mathcal{C}^{*}$. 
\begin{itemize}\item\textbf{The compact open topology on $\mathcal{C}$ is the same as the topology of uniform convergence, i.e. the topology induced by the norm $||\cdot||_{\infty}$.}
\\\\pf. Let $f_{0}\in\mathcal{C}$ 
\\Assume the open nhood of $f_{0}$ is of the form $X\equiv\{f\in\mathcal{C}:f(K_{1})\subset U_{1},\dots,f(K_{m})\subset U_{m}\}$ where $\forall K_{i}\subset G$ are compact and $\forall U_{j}\subset\mathbb{C}$ are open. Then $\forall \epsilon>0,\,\,\,V_{\epsilon}\equiv\{f\in\mathcal{C}:||f-f_{0}||_{\infty}<\epsilon\}$. Let $i\in\{1,\dots,m\}$. As $f_{0}(K_{i})\subset U_{i}$ is compact, $\exists \epsilon_{i}>0$ s.t. $\{x\in\mathbb{C}:\exists y\in f_{0}(K_{i}),\,\,\,|x-y|<\epsilon\}\subset U_{i}$. Let $\epsilon=\min_{i}\epsilon_{i}$. Then $V_{\epsilon}\subset X$. 
\\Let $\epsilon>0$ and $V_{\epsilon}$ be defined as above. Let $\eta=\epsilon/2$. Then $\forall x\in G,\,\,\,U_{x}\equiv\{y\in G:|f_{0}(x)-f_{0}(y)|<\eta\}\subset K_{x}\equiv\{y\in G:|f_{0}(x)-f_{0}(y)|\leq \eta\}$. Then $U_{x}$ is open and $K_{x}$ is compact. As $G=\cup_{c\in g}U_{x}$ is compact, $\exists x_{1},\dots,x_{n}\in G$ s.t. $G=\cup_{i=1}^{n}U_{x_{i}}=\cup_{i=1}^{n}K_{x_{i}}.\,\,\,\,\forall i,\,\,\,\,K_{i}=K_{x_{i}}$ for simplicity and let $U_{i}=\{a\in\mathbb{C}:|f_{0}(x_{i})-a|<\eta\}$. Then $X\equiv\{f\in\mathbb{C}:\forall i,\,\,\,\,\,f(K_{i})\subset U_{i}\}\subset V_{\epsilon}$.
\item \textbf{Let $\rho$ be the representation of $G$ on $\mathcal{C}$ defined by $(\rho(g)f)(x)=f(g^{-1}x),\,\,\,\forall f\in\mathcal{C},\,\,\,\,x,g\in G$. Then $\begin{array}{rrr}G\times \mathcal{C}&\rightarrow&\mathcal{C}\\(g,f)&\mapsto&\rho(g)f\end{array}$ is continuous.}
\\\\pf.  Denote by $\mu$ the map $G\times\mathcal{C}\rightarrow\mathcal{C}$. Let $\left\{\begin{array}{r}g_{0}\in G\\f_{0}\in\mathcal{C}\end{array}\right.$. As $G$ is compact, $f_{0}:G\rightarrow\mathbb{C}$ is uniformly continuous, so there exist two collections $\left\{\begin{array}{r}U_{1},\dots,U_{n}\\V_{1},\dots,V_{n}\end{array}\right.$ of open subsets of $G$ s.t. $\left\{\begin{array}{r}K_{i}\equiv\overline{U}_{i}\subset V_{i},\,\,\,\forall i\\G=\cup_{i=1}^{n}U_{i}$ and $\forall i,\,\,\,\forall x,y\in V_{i},\,\,\,|f_{0}(x)-f_{0}(y)|<\epsilon/2\end{array}\right.$. Let $i\in\{1,\dots,n\}.\,\,\,\,\forall x\in K_{i},\,\,\,\,\exists j$ s.t. $g_{0}^{-1}x\in V_{j}$ and open nhood $\left\{\begin{array}{r}U_{x}$ of $x$ in $G\\W_{x}$ of $g_{0}$ in $ G\end{array}\right.$ s.t. $\forall g\in W_{x},\,\,\,\forall y\in U_{x},\,\,\,g^{-1}y\in V_{j}$. As $K_{i}$ is compact, we can choose $x_{1},\dots,x_{m}$ s.t. $K_{i}\subset U_{x_{1}}\cup\dots\cup U_{x_{m}}$. Then the open nhood $W_{i}\equiv W_{x_{1}}\cap\dots\cap W_{x_{m}}$ of $g_{0}$ has the property that $\forall g\in W_{i},\,\,\,\forall x\in K_{i},\,\,\,\exists j\in\{1,\dots,n\}$ s.t. $g_{0}^{-1}x,g^{-1}x\in V_{j}$. Define $W\equiv\cap_{i=1}^{n}W_{i}$. Let $g\in W$ and $f\in\mathcal{C}$ s.t. $||f-f_{0}||_{\infty}<\epsilon/2$. Then $\forall x\in G,$ choose $i$ s.t. $x\in K_{i}$. Then $\exists j\in\{1,\dots,n\}$ s.t. $g_{0}^{-1}x,g^{-1}x\in V_{j}$. Then $|(gf-g_{0}f_{0})(x)|=|f(g^{-1}x)-f_{0}(g_{0}^{-1}x)|=|(f(g^{-1}x)-f_{0}(g^{-1}x))+(f_{0}(g^{-1}x)-f_{0}(g_{0}^{-1}x))|\leq |(f(g^{-1}x)-f_{0}(g^{-1}x))|+|((f_{0}(g^{-1}x)-f_{0}(g_{0}^{-1}x))|<\epsilon/2+\epsilon/2=\epsilon,$ so it is continuous.
\item \textbf{The contragredient representation $\rho^{*}$ s.t. $\rho^{*}(g)\lambda=\lambda \circ \rho(g^{-1}),\,\,\forall g\in G,\,\,\,\lambda\in\mathcal{C}^{*},$ is continuous.}
\\\\pf. Let $g_{0}\in G,\,\,\,\lambda_{0}\in \mathcal{C}^{*},\,\,\,f\in \mathcal{C},\,\,\,r\in\mathbb{R}_{++}$. Let $||\lambda_{0}||=\sup\{||\lambda_{0}(f_{1})||_{\infty},\,\,\,\,f_{1}\in\mathcal{C},\,\,\,||f_{1}||_{\infty}=1\}$. Then $||\lambda_{0}||<\infty$ since $\lambda_{0}$ is continuous and $||\lambda_{0}(f_{1})||_{\infty}\leq||\lambda_{0}||||f_{1}||_{\infty},\,\,\,\forall f_{1}\in\mathcal{C}$. Then by the previous statement, there exists a nhood $W$ of $g_{0}$ in $G$ s.t. $\forall g\in W,\,\,\,||g^{-1}f-g_{0}^{-1}f||_{\infty}<\frac{r}{2(1+||\lambda_{0}||)}$. Let $U\equiv\{\lambda\in\mathcal{C}^{*}:|(\lambda-\lambda_{0})(g^{-1}f)|<r/2\}$. Then $U$ is a nhood of $\lambda_{0}$ in $\mathcal{C}^{*}$ and $\forall g\in W,\,\,\,\,\,\forall\lambda\in U,\,\,\,\,|(g\lambda-g_{0}\lambda_{0})(f)|=|\lambda(g^{-1}f)-\lambda_{0}(g^{-1}f)|+|\lambda_{0}(g^{-1}f)-\lambda_{0}(g_{0}^{-1}f)|<r$. 
\item \textbf{$\mathcal{P}(G)$ is a convex compact $G$-invariant subset of $\mathcal{C}^{*}$.}
\\\\pf. Denote by $B$ its unit ball of $\mathcal{C}$ for the sup-norm. Let $B_{+}$ be the set of $f\in B$ s.t. $f(G)\subset\mathbb{R}_{+}$. Then $B_{+}$ generates $\mathcal{C}$ as a vector space so the map $\begin{array}{rrr}\mathcal{C}^{*}&\rightarrow&\Pi_{B_{+}}\mathbb{C}\\\lambda&\mapsto&(\lambda(f))_{f\in\mathcal{C}}\end{array}$ is injective, and we can use it to identify $\mathcal{C}^{*}$ to a subspace of $\Pi_{B_{+}}\mathbb{C}$. By the definition, the weak topology on $\mathcal{C}^{*}$ is the topology induced by the product topology on $\Pi_{B}\mathbb{C}$. Let $K=\Pi_{B}[0,1]\subset\Pi_{B}\mathbb{C}$. Then by Tychonoff's theorem, $K$ is compact. By the Riesz representation theorem, $\mathcal{C}^{*}\cap K$ is closed in $K$. Indeed, if $(a_{f})_{f\in B_{+}}$ is in the closure of $\mathcal{C}^{*}\cap K,\,\,\,\,\begin{array}{rrr}B_{+}&\rightarrow&\mathbb{C}\\f&\mapsto&a_{f}\end{array}$ extends to a linear map $\mathcal{C}\rightarrow\mathbb{C}$ and this linear map is positive, given $a_{f}\in\mathbb{R}_{+}$. Thus it is of the form $f\mapsto \int fd\mu$, where $\mu$ is a regular Borel measure on $G$ and hence it is continuous, i.e. an element of $\mathcal{C}^{*}$. So $K\cap \mathcal{C}^{*}$ is compact and therefore $\mathcal{P}(G)$ is compact given $\mathcal{P}(G)=\{\lambda\in K\cap\mathcal{C}^{*}|\lambda(1)=1\}$ is closed in $L\cap\mathcal{C}^{*}$. The other properties are trivial. 
\item \textbf{$G$ has a left Haar measure.}
\\\\pf. $\mathcal{P}(G)$ is nonempty because $f\mapsto f(1)$ is in $\mathcal{P}(G)$. $\mathcal{C}^{*}$ is locally convex since $\{\lambda\in\mathcal{C}^{*} :|\lambda(f_{1})<r_{1},\dots,|\lambda(f_{s})|<r_{s}\},\,\,\,\forall f_{i}\in\mathcal{C},\,\,\,\forall r_{i}\in\mathbb{R}_{++}$ form a basis of nhoods of 0 in $\mathcal{C}^{*}$ and they are convex. Also the action of $G$ on $\mathcal{P}(G)$ is equicontinuous, i.e. for every nhood of $W$ of 0 in $\mathcal{C}^{*}$, there exists a nhood $V$ of 0 in $\mathcal{C}^{*}$ s.t. $\forall \mu_{1},\mu_{2}\in\mathcal{P}(G)$ s.t. $\mu_{1}-\mu_{2}\in V,\,\,\,\,g\cdot(\mu_{1}-\mu_{2})\in W,\,\,\forall g\in G$. Then by the Markov-Kakutani fixed point theorem, $\exists \mu\in\mathcal{P}(G)$ s.t. $g\mu=\mu,\,\,\forall g\in G$. This is a left Haar measure. 

\end{itemize}


\begin{theorem}Let $G$ be a compact Hausdorff topological group. Let $\mathcal{C}(G,\mathbb{C})$ be the $\mathbb{C}$-algebra of continuous functions from $G$ to $\mathbb{C}$ with the norm $||\cdot||_{\infty}$, the supremum norm. Then there exists a unique $\mathbb{C}$-linear map $\lambda:\mathcal{C}(G,\mathbb{C})\rightarrow\mathbb{C}$ s.t. $\left\{\begin{array}{r}\lambda $ is positive, i.e. $\lambda(f)\geq0,\,\,\forall f(G)\subset\mathbb{R}_{+}\\\lambda$ is left invariant, i.e. $\lambda(f)=\lambda(f(g\cdot _ )),\,\,\,\forall f\in\mathcal{C}(G,\mathbb{C}),\forall g\in G\\\lambda$ is right invariant, i.e. $\lambda(f)=\lambda(f( _\cdot g)),\,\,\,\forall f\in\mathcal{C}(G,\mathbb{C}),\forall g\in G\\\lambda(1)=1\end{array}\right.$.Moreover, if $\lambda$ is continuous, there exists a unique probability measure $dg$ on $G$ s.t. $\forall f\in\mathcal{C}(G,\mathbb{C}),\,\,\,\lambda(f)=\int_{G}f(g)dg$ and this measure satisfies \begin{equation*}\int_{G}f(g)dg=\int_{G}f(g^{-1})dg\,\,\,\,\,\,\forall f:G\rightarrow\mathbb{C},\text{measurable.}\end{equation*}\end{theorem}
$\rightarrow dg$ is called a bi-invariant probability Haar measure on $G$.
\\\\\textbf{There exists a function $c:G\rightarrow\mathbb{R}_{++}$ s.t $\forall g\in G,\,\,\,\forall f\in\mathcal{C}(G,\mathbb{C}),\,\,\,\int_{G}f(xg)d\mu(x)=c(g)\int_{G}f(x)d\mu(x)$.}
\\\\pf. Let $g\in G$. Then $d\mu(x),\,\,\,d\mu(xg^{-1})$ are both left Haar measures on $G$ so $\exists c(g)\in\mathbb{R}_{++}$ s.t. $d\mu(xg^{-1})=c(g)d\mu(x)$ i.e. $\forall f\in \mathcal{C}(G,\mathbb{C}),$ by the previous theorem uniqueness and therefore the above equation holds.
\\$\rightarrow \forall g_{1},g_{2}\in G,\,\,\,\,c(g_{1}g_{2})d\mu(x)=d\mu(x(g_{1}g_{2})^{-1})=d\mu(xg_{2}^{-1}g_{1}^{-1})=c(g_{1})d\mu(xg_{2}^{-1})=c(g_{1})c(g_{2})d\mu(x)$ so $c(g_{1}g_{2})=c(g_{1})c(g_{2})$, hence $c:G\rightarrow(\mathbb{R}_{++},\times)$ is a morphism of groups. Let $\epsilon>0$. Choose $f\in \mathcal{C}(G,\mathbb{R}_{+})$ s.t. $\int_{G}f(x)d\mu(x)=1$. Let $U$ be a nhood of 1 in $G$ s.t. $\forall x\in G,\,\,\,\forall y\in U,\,\,\,|f(x)-f(xy)|\leq\epsilon$ and let $K$ be the suport of $f$. Then if $g,g'\in G$ s.t. $g^{-1}g'\in U,\,\,\,\,|c(g')-c(g)|=|(c(g')-g(g))\int_{G}f(x)d\mu(x)|=|\int_{G}f(xg')d\mu(x)-\int_{G}f(xg)d\mu(x)|\leq\epsilon\mu(K)$ hence it is continuous.
\begin{definition} The function $c:G\rightarrow(\mathbb{R}_{++},\times)$ is called the modular function of $G$ and we say $G$ is unimodular if $c(G)=1$.\end{definition}
\textbf{If $G$ is compact, it is unimodular.}
\\\\pf. If $G$ is compact, $c(G)$ is a compact subgroup of $\mathbb{R}_{++}$. The only compact subgroup of $\mathbb{R}_{++}$ is $\{1\}$, so $c(G)=\{1\}$, i.e. $G$ is unimodular. 
\\\\\indent\textbf{Let $f\in\mathcal{C}(G,\mathbb{C})$. Then $\int_{G}f(x^{-1})d\mu(x)=\int_{G}c(x)f(x)d\mu(x)$.}
\\\\pf. $\int_{G}\int_{G}g(yx)c(x)^{-1}f(x^{-1})d\mu(x)d\mu(y)
\\=\left\{\begin{array}{rr}\int_{G}f(x^{-1})c(x)^{-1}(\int_{G}g(yx)d\mu(y))d\mu(x)=\int_{G}f(x^{-1})d\mu(x)\int_{G}g(y)d\mu(y)&$ by the definition of $c\\\int_{G}\int_{G}g(z)c(y^{-1}z)^{-1}f(z^{-1}y)d\mu(z)d\mu(y)=\int_{G}g(z)d\mu(z)\int_{G}c(t)f(t)d\mu(t)&$ by $z=yx,\,\,\,t=z^{-1}y\end{array}\right.$. So the equivalence follows by choosing $\int_{G}g(x)d\mu(x)=1$. 
\\\\\indent If $G$ is just locally compact, we can find nonzero $\mathbb{C}$-linear map satisfying the first and second or the first and third and they are unique upto multiplication by a positive real number. These are also continuous and the corresponding measures on $G$ are called left-invariant or right-invariant Haar measure. Then $\int_{G}f(g)d_{r}g=\int_{G}f(g^{-1})d_{l}g,\,\,\,\left\{\begin{array}{rr}d_{l}& $ left invariant$\\d_{r}&$ right invariant $\end{array}\right.$ for every measurable function $f:G\rightarrow\mathbb{C}$. 
\\\\\\e.g.\begin{itemize}\item a finite group $G$ with the discrete topology is a topology group. A left and right-invariant Haar measure on $G$ is given by $dg(A)=|A|/|G|$. 
\item $G=U(1)\equiv\{z\in\mathbb{C} : |z|=1\}$ with the topology induced by that of $\mathbb{C}$. Then it is a topological group and we have an isomorphism of topological groups $\begin{array}{rrr}\phi:\mathbb{R}/\mathbb{Z}&\rightarrow&U(1)\\t&\mapsto&e^{2i\pi t}\end{array}$. Then $\int_{G}f(g)dg=\int_{0}^{1}(f\circ\phi)(t)dt=\int_{0}^{1}f(e^{2i\pi t})dt$ where $dt$ is the Lebesgue measure on $\mathbb{R}$.\end{itemize}
\begin{definition}Let $G$ be a topological group and $V$ be a normed $\mathbb{C}$-vector space. Then a continuous representation of $G$ on $V$ is an abstract representation of $G$ on $V$ s.t. the action map $\begin{array}{rrr}G\times V&\rightarrow&V\\(g,v)&\mapsto&gv\end{array}$ is continuous.\end{definition}
$\rightarrow $ let End$(V)$ be the $\mathbb{C}$-algebra of continuous endomorphisms of $V$. Put the operator norm on End$(V)$ and consider a continuous representation of $G$ on $V$. Then the action of each $ g\in G$ on $V$ is continuous endomorphism of $V$, so $G\rightarrow \text{End}(V)$ exists. But it is not necessarily continuous. 
\begin{definition}If $V$ is a normed $\mathbb{C}$-vector space, we denote by End$(V)$ the $\mathbb{C}$-algebra of continuous endomorphisms of $V$ and we put on it the topology given by the operator norm. Write $GL(V)$ for End$(V)^{\times}$, with the topology induced by that of End$(V)$.\end{definition}
\begin{proposition}Let $V$ be a normed $\mathbb{C}$-vector space and $\rho:G\rightarrow GL(V)$ be a morphism of groups. Consider the following conditions: \begin{enumerate}\item the map $\begin{array}{rrr}G\times V&\rightarrow&V\\(g,v)&\mapsto&\rho(g)v\end{array}$ is continuous\item $\forall v\in V,$ the map $\begin{array}{rrr}G&\rightarrow&V\\g&\mapsto&\rho(g)(v)\end{array}$ is continuous\item the map $\rho:G\rightarrow GL(V)$ is continuous\end{enumerate} Then $3\rightarrow 1\rightarrow 2$. If moreover $V$ is finite-dimensional, then all three conditions are equivalent.\end{proposition}
pf.\begin{itemize}\item $1\rightarrow 2$ is trivial.
\item $3\rightarrow 1$
\\\\Let $\left\{\begin{array}{r}g_{0}\in G\\v_{0}\in V\end{array}\right.$ and $\epsilon>0$. Choose a $\delta$ s.t. $0<\delta\leq\frac{\epsilon}{2||\rho(g_{0})||}$ and let $U$ be a nhood of $g_{0}$ in $G$ s.t. $g\in U$ implies $||\rho(g)-\rho(g_{0})||<\frac{\epsilon}{2(||v_{0}||+\delta)}$. Then if $\left\{\begin{array}{r}g\in U\\||v-v_{0}||<\delta\end{array}\right., \\\begin{array}{rr}||\rho(g)(v)-\rho(g_{0})(v_{0})||\leq ||\rho(g)(v)-\rho(g_{0})(v)||+||\rho(g_{0})(v)-\rho(g_{0})(v_{0})||& $ by the triangular inequality $\\\leq ||\rho(g)-\rho(g_{0})||||v||+||\rho(g_{0})||||v-v_{0}||&\\<\frac{\epsilon}{2(||v_{0}||+\delta)}(||v_{0}||+\delta)+||\rho(g_{0})||\delta& $ since $\left\{\begin{array}{r}||v||\leq||v_{0}||+\delta\\$ the choice of $g,\delta\end{array}\right.\\\leq \epsilon/2+\epsilon/2=\epsilon&\end{array}$. 

\item given $V$ finite-dimensional, $2\rightarrow 3$
\\\\Let $(e_{1},\dots,e_{n})$ be a basis of $V$ and let $||\cdot||$ be the supremum norm. We use the corresponding operator norm on End$(V)$ and denote it by $||\cdot||$. Let $g_{0}\in G$ and $\epsilon>0$. For every $i,$ the function $\begin{array}{rrr}G&\rightarrow&V\\g&\mapsto&\rho(g)(e_{i})\end{array}$ is continuous by 2, so there exists a nhood $U_{i}$ of $g_{0}$ in $G$ s.t. $g\in U_{i}$ implies $||\rho(g)(e_{i})-\rho(g_{0})(e_{i})||\leq \epsilon/n$. Let $U=\cap_{i=1}^{n}U_{i}$. Then if $g\in U,$ then for all $v=\sum_{i=1}^{n}x_{i}e_{i}\in V,\,\,\,||\rho(g)(v)-\rho(g_{0})(v)||\leq\sum|x_{i}|||\rho(g)(e_{i})-\rho(g_{0})(e_{i})||\leq\sum|x_{i}|\epsilon/n\leq \epsilon||v||$ so $||\rho(g)-\rho(g_{0})||\leq \epsilon$.
\end{itemize}
A complex Hilbert space is a $\mathbb{C}$-vector space $V$ with a Hermitian inner product s.t. $V$ is complete for the corresponding norm. If $V$ is a finite dimensional $\mathbb{C}$-vector space with a Hermitian inner product, it is complete and therefore a Hilbert space. Write $U(V)$ for the group of unitary endomorphisms of $V$. 
\begin{proposition}If $V$ is a Hilbert space and $\rho:G\rightarrow U(V)$ is a morphism of groups, then TFAE\begin{enumerate}\item the map $\begin{array}{rrr}G\times V&\rightarrow&V\\(g,v)&\mapsto&\rho(g)(v)\end{array}$ is continuous\item $\forall v\in V,$ the map $\begin{array}{rrr}G&\rightarrow&V\\g&\mapsto &\rho(g)(v)\end{array}$ is continuous.\end{enumerate}\end{proposition}
pf. \begin{itemize}\item by the previous proposition, $2\rightarrow 1$.
\item Let $\left\{\begin{array}{r}g_{0}\in G\\v_{0}\in V\end{array}\right.$ and $\epsilon>0$. Choose a nhood $U$ of $g_{0}$ in $G$ s.t. $g\in U$ implies $||\rho(g)(v_{0})-\rho(g_{0})(v_{0})||<\epsilon/2$ and take $\delta=\epsilon/2$. Then if $g\in U\\||v-v_{0}||<\delta,\,\,\,\,\begin{array}{r}||\rho(g)(v)-\rho(g_{0})(v_{0})||\leq||\rho(g)(v)-\rho(g)(v_{0})||+||\rho(g)(v_{0})-\rho(g_{0})(v_{0})||\\<||\rho(g)||||v-v_{0}||+\epsilon/2\\<\epsilon/2+\epsilon/2=\epsilon\end{array}$, since $\rho(g)\in U(V)$.
\end{itemize}

\textbf{CLAIM : if $\rho:G\rightarrow U(V)$ is a unitary representation of $G$ on a Hilbert space $(V,<\cdot,\cdot>),$ for every $G$-invariant subspace $W$ of $V$, the subspace $W^{\perp}$ is also $G$-invariant.}
\\pf. if $w\in W^{\perp},\,\,\,g\in G$, then $\forall v\in W,\,\,\,<v,\rho(g)w>=<\rho(g)^{-1}v,w>=0$ since $\rho(g)^{-1}v\in W$ by the $G$-invariance of $W$. This implies $\rho(g)W\in W^{\perp}$ and $W^{\perp}$ is $G$-invariant. 
\\$\rightarrow $ if $W$ is a closed $G$-invariant subspace of $V,\,\,\,V=W\oplus W^{\perp}$ with $W^{\perp}$ a closed $G$-invariant subspace. If $V$ is finite dimensional, every subspace is closed and therefore every unitary representation of $G$ on $V$ is semisimple. 

\begin{theorem}Assume that the group $G$ is compact Hausdorff. Let $V$ be a finite-dimensional $\mathbb{C}$-vector space and $\rho:G\rightarrow GL(V)$ be a continuous representation of $G$ on $V$. Then there exists a Hermitian inner product on $V$ that makes $\rho$ a unitary representation.\end{theorem}
pf. Let $<\cdot,\cdot>_{0}$ be a Hermitian inner product on $V$. Define $\\\begin{array}{rrr}<\cdot,\cdot>:V\times V&\rightarrow&\mathbb{C}\\(v,w)&\mapsto&\int_{G}<\rho(g)v,\rho(g)w>_{0}dg\end{array}$. Then $<\rho(g)v,\rho(g)w>=<v,w>,\,\,\,\forall v,w\in V$ and $\forall g\in G$. It is a Hermitian product trivially. Let $v\in V\setminus\{0\}$. Then the function $\begin{array}{rrr}G&\rightarrow\mathbb{R}\\g&\mapsto&<\rho(g)v,\rho(g)v>_{0}\end{array}$ is continuous and takes positive values. As $G$ is compact, $\exists \epsilon>0$ s.t. $<\rho(g)v,\rho(g)v>_{0}>\epsilon,\,\,\,\forall g\in G$. Then $<v,v>\geq\epsilon>0$ and it is positive definite. 

\begin{corollary}If $G$ is compact Hausdorff, then every finite-dimensional continuous representation of $G$ is semisimple.\end{corollary}
\begin{definition}Let $V,W$ be Hermitian inner product spaces. If $v\in V,\,\,\,w\in W$, we define the $\mathbb{C}$-linear map $\begin{array}{rrr}T_{v,w}^{0}:V&\rightarrow&W\\x&\mapsto&<x,v>w\end{array}$.\end{definition}
$\rightarrow\begin{array}{rrr}V&\rightarrow&\hom_{\mathbb{C}}(V,W)\\v&\mapsto&T_{v,w}^{0}\end{array}$ is semi-linear.
\begin{proposition}\begin{enumerate}\item $\forall v\in V,\,\,\,\forall w\in W,\,\,\,(T_{v,w}^{0})^{*}=T_{w,v}^{0}$\item $\forall v\in V,\,\,\,\forall w\in W,\,\,\,T_{v,w}^{0}$ generate $\hom_{\mathbb{C}}(V,W)$ as a $\mathbb{C}$-vector space\item if $V=W,\,\,\,\,\forall v,w\in V,\,\,\,T(T_{v,w}^{0})=<w,v>$\item $\forall v_{1},v_{2}\in V,\,\,\,\forall w_{1},w_{2}\in W,\,\,\,T((T_{v_{1},w_{1}}^{0}T_{v_{2},w_{2}}^{0})^{*})=<w_{1},w_{2}><v_{2},v_{1}>$\end{enumerate}\end{proposition}
pf.\begin{itemize}\item $\forall x\in V,\,\,\,y\in W,\,\,\,<T_{v,w}^{0}(x),y>=<<x,v>w,y>=<x,v><w,y>=<x,v>\overline{<y,w>}\\=<x,<y,w>v>=<x,T_{w,v}^{0}(y)>$
\item if we choose the orthonormal bases $(e_{i})$ of $V$ and $(f_{j})$ of $W,\,\,\,\,\forall i,j$, the matrix of $T_{e_{i},f_{j}}^{0}$ in these bases is the $m\times n$ with $(i,j)$-entry equal to 1 and all other entries equal to 0. These clearly generate the $\mathbb{C}$-vector space $M_{mn}(\mathbb{C})$. 
\item let $(v_{1},\dots,v_{n})$ be an orthogonal basis of $V$ s.t. $v_{1}=v$. Then $T_{v,w}^{0}(v_{i})=\left\{\begin{array}{rr}<v_{1},v_{1}>w&i=1\\0&otherwise\end{array}\right.$. As $w=\sum \frac{<w,v_{i}>}{<v_{i},v_{i}>}v_{i}$, the proposition follows
\item $\forall y\in W,\,\,\,T_{v_{1},w_{1}}^{0}(T_{v_{2},w_{2}}^{0})^{*}=T_{v_{1},w_{1}}^{0}T_{w_{2},v_{2}}^{0}(y)=T_{v_{1},w_{1}}^{0}(<y,w_{2}>v_{2})=<y,w_{2}><v_{2},v_{1}>w_{1}$. Choose an orthogonal basis $(y_{1},\dots,y_{n})$ of $W$ s.t. $y_{1}=w_{2}$. Then $T_{v_{1},w_{1}}^{0}(T_{v_{2},w_{2}}^{0}(y))^{*}(y_{i})=\left\{\begin{array}{rr}<y_{1},y_{1}><v_{2},v_{1}>w_{1}&i=1\\0&otherwise\end{array}\right.$. As $w_{1}=\sum\frac{<w_{1},y_{i}>}{<y_{i},y_{i}>}y_{i}$, the proposition follows.
\end{itemize}
Let $G$ be a Hausdorff compact group. 
\begin{definition}Let $\rho:G\rightarrow GL(V)$ be a continuous representation of $G$ on a finite dimensional complex vector space. Remember that the map $\begin{array}{rrr}\chi_{V}:G&\rightarrow&\mathbb{C}\\g&\mapsto&T(\rho(g))\end{array}$ is called the character of the representation $(V,\rho)$.\end{definition}
$\rightarrow \chi_{V}:G\rightarrow\mathbb{C}$ is a continuous map by the previous proposition A.3.1.
\begin{theorem}Let $V,W$ be continuous representations of $G$ on finite-dimensional complex vector spaces. Assume $V,W$ are both irreducible. Choose $G$-invariant Hermitian inner products on $V,W$ that will both be denoted by $<\cdot,\cdot>$. Then $\left\{\begin{array}{r}\forall v_{1},v_{2}\in V\\\forall w_{1},w_{2}\in W\end{array}\right.$, \begin{equation*}\int_{G}<gv_{1},v_{2}>\overline{<gw_{1},w_{2}>}dg=\left\{\begin{array}{rr}\frac{1}{\dim_{\mathbb{C}}V}<v_{1},w_{1}>\overline{<v_{2},w_{2}>}&V\cong W\\0&otherwise\end{array}\right.\end{equation*}\end{theorem}
pf. If $v\in V,w\in W,$ let $T_{v,w}=\int_{G}gT_{v,w}^{0}g^{-1}dg\in\hom_{\mathbb{C}}(V,W)$ where $T_{v,w}^{0}$ is defined above. Then $T_{v,w}$ is $G$-equivariant, so by the Schur's lemma, $T_{v,w}=\left\{\begin{array}{rr}0&V\not\cong W\\c(v,w)id_{V}&V\cong W\end{array}\right.$ with $c(v,w)\in\mathbb{C}$. Suppose $V\cong W$, then $c(v,w)\dim_{\mathbb{C}}V=T(T_{v,w})=\int_{G}T(gT_{v,w}^{0}g^{-1})dg=T(T_{v,w}^{0})$ so by the previous proposition, $c(v,w)=\frac{1}{\dim_{\mathbb{C}}V}<w,v>$. Let $I=\int_{G}<gv_{1},v_{2}>\overline{<gw_{1},w_{2}>}dg$. Then $I=\int_{G}<gv_{1},v_{2}><g^{-1}w_{2},w_{1}>dg=\int_{G}<<g^{-1}w_{2},w_{1}>gv_{1},v_{2}>dg=\int_{G}<gT_{w_{1},v_{1}}^{0}g^{-1}w_{2},v_{2}>dg=<T_{w_{1},v_{1}}(w_{2}),v_{2}>$ so $I=\left\{\begin{array}{rr}0&V\not\cong W\\c(w_{1},v_{1})<w_{2},v_{2}>&V\cong W\end{array}\right.$ and the theorem follows.
\begin{corollary}[Schur's Orthogonality] Let $V,W$ be continuous representations of $G$ on finite dimensional complex vector spaces. Assume $V,W$ are both irreducible. Then \begin{equation*}\int_{G}\chi_{V}(g)\overline{\chi_{W}(g)}dg=\left\{\begin{array}{rr}0&V\not\cong W\\1&V\cong W\end{array}\right.\end{equation*}\end{corollary}
pf. Choose $G$-invariant Hermitian inner products on $V,W$ and fix orthonormal bases $\left\{\begin{array}{rr}(v_{1},\dots,v_{n}& $ for $V\\(w_{1},\dots,w_{n})&$ for $W\end{array}\right.$. Then for any $\mathbb{C}$-linear endomorphism $u$ of $V$ or $W$, we have $T(u)=\sum<u(v_{i}),v_{i}>$ or $\sum<u(w_{j}),w_{j}>$. In particular, $\forall g\in G,\,\,\,\left\{\begin{array}{r}\chi_{V}(g)=\sum<gv_{i},v_{i}>\\\chi_{W}(g)=\sum<gw_{j},w_{j}>\end{array}\right.$ and therefore $\int_{G}\chi_{V}(g)\overline{\chi_{W}(g)}dg=\sum_{i,j}\int_{G}<gv_{i},v_{i}>\overline{<gw_{j},w_{j}>}dg=\left\{\begin{array}{rr}0&V\not\cong W\\\frac{1}{\dim_{\mathbb{C}}V}\sum_{i,j}<v_{i},v_{j}>\overline{<v_{i},v_{j}>}=1&V\cong W\end{array}\right.$ and the corollary follows.
\begin{definition}$L^{2}(G)$ is the quotient \begin{equation*}\{f:G\rightarrow\mathbb{C},\text{ measurable }|\int_{G}|f(g)|^{2}dg<\infty\}/\{f:G\rightarrow\mathbb{C},\text{ measurable }|\int_{G}|f(g)|^{2}dg=0\}\end{equation*}
This is a Hilbert space for the Hermitian inner product given by $<f_{1},f_{2}>=\int_{g\in G}f_{1}(g)\overline{f_{2}(g)}dg$. \end{definition}
\begin{definition}If $x\in G,$ and $f$ is a function from $G$ to $\mathbb{C}$, define $\\\begin{array}{rrr}L_{x}f:G&\rightarrow&\mathbb{C}\\g&\mapsto&f(xg)\end{array}, \,\,\,\begin{array}{rrr}R_{x}f:G&\rightarrow&\mathbb{C}\\g&\mapsto&f(gx)\end{array}$. \end{definition}
\begin{proposition}$\forall x,y\in G$, the operators $L_{x},R_{x}$ induce endomorphisms $L_{x},R_{x}$ of $L^{2}(G)$ with the following properties; $\begin{array}{r}L_{x}R_{y}=R_{y}L_{x},\,\,\,R_{x}R_{y}=R_{xy},\,\,\,\,L_{x}L_{y}=L_{yx}\\R_{x}^{*}=R_{x}^{-1}=R_{x^{-1}},\,\,\,L_{x}^{*}=L_{x}^{-1}=L_{x^{-1}}\end{array}$. In particular, $R_{x},L_{x}$ are unitary endomorphisms of $L^{2}(G)$.\end{proposition}
pf. \begin{itemize}
\item Let $f_{1},f_{2}:G\rightarrow\mathbb{C}$ be measurable functions. Then $\\\left\{\begin{array}{rr}<R_{x}f_{1},R_{x}f_{2}>=\int_{G}f_{1}(gx)\overline{f_{2}(gx)}dg=\int_{G}f_{1}(g)\overline{f_{2}(g)}dg=<f_{1},f_{2}>&$ by the right invariance $\\<L_{x}f_{1},L_{x}f_{2}>=\int_{G}f_{1}(xg)\overline{f_{2}(xg)}dg=\int_{G}f_{1}(g)\overline{f_{2}(g)}dg=<f_{1},f_{2}>&$ by the left-invariance$ \end{array}\right.$. Thus $R_{x},L_{x}$ preserves both spaces in the quotient defining $L^{2}(G)$ and therefore they induce endomorphisms of $L^{2}(G)$. 
\item Let $f\in L^{2}(G),\,\,\,g\in G$. Then $\left\{\begin{array}{r}(L_{x}R_{y}f)(g)=(R_{y}f)(xg)=f(xgy)=(L_{x}f)(gy)=(R_{y}L_{x}f)(g)\\(L_{x}L_{y}f)(g)=(L_{y}f)(xg)=f(yxg)=(L_{yx}f)(g)\\(R_{x}R_{y}f)(g)=(R_{y}f)(gx)=f(gyx)=(R_{xy}f)(g)\end{array}\right.$
\item by the first stmt, $L_{x},R_{x}$ are unitary endomorphisms and by the second stmt, $L_{x}L_{x^{-1}}=L_{xx^{-1}}=L_{1}=Id\rightarrow L_{x}^{-1}=L_{x^{-1}}$ and similarly for $R_{x}$. \end{itemize}
\begin{definition}If $f$ is a function $G\rightarrow\mathbb{C}$, define a function $\begin{array}{rrr}\tilde{f}:G&\rightarrow&\mathbb{C}\\g&\mapsto&\overline{f(g^{-1})}\end{array}$.\end{definition}
\begin{proposition}The above function defines a unitary endomorphism $f\mapsto \tilde{f}$ of $L^{2}(G)$.\end{proposition}
pf. Let $f:G\rightarrow\mathbb{C}$ be a measurable function. Then, $\int_{G}\overline{|f(g^{-1})|}^{2}dg=\int_{G}|f(g)|^{2}dg$ because $G$ being compact implies the Haar measure $dg$ is equal to its own pullback by $g\mapsto g^{-1}$. Then the proposition follows. 
\begin{definition}If $f_{1},f_{2}\in L^{2}(G)$, define $\begin{array}{rrr}f_{1}*f_{2}:G&\rightarrow&\mathbb{C}\\g&\mapsto&\int_{G}f_{1}(gx)f_{2}(x^{-1})dx\end{array}$, called the convolution product of $f_{1},f_{2}$.\end{definition}
\begin{proposition}\begin{enumerate}\item $\forall f_{1},f_{2}\in L^{2}(G),\,\,\,\,f_{1}*f_{2}$ is a bounded function on $G$ and we have $||f_{1}*f_{2}||_{\infty}\leq||f_{1}||_{2}||f_{2}||_{2}$\item the convolution product is associative and distributive w.r.t. addition\end{enumerate}
\end{proposition}
pf.\begin{itemize}\item $\forall g\in G,\,\,\,\,|f_{1}*f_{2}(g)|=|\int_{G}f_{1}(gx)f_{2}(x^{-1})dx|=|\int_{G}(L_{g}f_{1})\overline{\tilde{f}_{2}(x)}dx|\leq||L_{g}f_{1}||_{2}||\tilde{f}_{2}||_{2}=||f_{1}||_{2}||f_{2}||_{2}$ by CS inequality. 
\item $\forall f_{1}f_{2},f_{3}\in L^{2}(G),\,\,\,\forall g\in G
\\\left\{\begin{array}{r}((f_{1}*f_{2})*f_{3})(g)=\int_{G}(f_{1}*f_{2})(gx)f_{3}(x^{-1})dx=\int_{G}(\int_{G}f_{1}(gxy)f_{2}(y^{-1})dy)f_{3}(x^{-1})dx\\(f_{1}*(f_{2}*f_{3}))(g)=\int_{G}f_{1}(gz)(f_{2}*f_{3})(z^{-1})dz=\int_{G}f_{1}(gz)(\int_{G}f_{2}(z^{-1}t)f_{3}(t^{-1})dt)dz\end{array}\right.$. By the change of variables $\left\{\begin{array}{r}z\mapsto y\\z\mapsto y^{-1}x\end{array}\right.$ and Fubini's theorem and bi-invariance Haar measure $dg$, they are equivalent, hence it is associative. 


\end{itemize}
$\rightarrow f_{1}*f_{2}$ are continuous for all $f_{1},f_{2}\in L^{2}(G)$ since if $f_{1},f_{2}$ are the limit in $L^{2}(G)$ of sequences $f_{(1,n)},f_{(2,n)},$ then the sequence $(f_{1,n}*f_{2,n})\rightarrow f_{1}*f_{2}$ in $L^{\infty}(G)$. Since the continuous functions are dense in $L^{2}(G)$ and the convolution of the two continuous functions is continuous, the continuity follows. 
\begin{definition}Let $f\in L^{2}(G)$. Define endomorphisms $L_{f},R_{f}$ of $L^{2}(G)$ by $\left\{\begin{array}{r}L_{f}(f_{1})=f*f_{1}\\R_{f}(f_{1})=f_{1}*f\end{array}\right.$.\end{definition}
$\rightarrow $ well-defined and continuous by the previous proposition. 
\begin{proposition}\begin{enumerate}\item $\forall f\in L^{2}(G),\,\,\,R_{f}^{*}=R_{\tilde{f}},\,\,\,L_{f}^{*}=L_{\tilde{f}}$.\item $\forall f\in L^{2}(G),\,\,\,\,\forall x\in G,\,\,\,R_{x}L_{f}=L_{f}R_{x},\,\,\,\,\,R_{f}L_{x}=L_{x}R_{f}$\end{enumerate}\end{proposition}
pf.\begin{itemize}\item $\forall f_{1},f_{2}\in L^{2}(G),\\\begin{array}{rr}<R_{f}(f_{1}),f_{2}>=\int_{G}(f_{1}*f)\overline{f_{2}(g)}dg&=\int_{G}\int_{G}f_{1}(gx^{-1})f(x)\overline{f_{2}(g)}dxdg\\=\int\int f_{1}(y)\overline{f_{2}(yx)}f(x)dxdy & $ by the change of variables $y=gx^{-1}\\=\int\int f_{1}(y)\overline{f_{2}(yz^{-1})}\overline{\tilde{f}(z)}dzdy&$ by the change of variables $z=x^{-1}\\=\int_{G}f_{1}(y)\overline{(f_{2}*\tilde{f})(y)}dy=<f_{1},R_{\tilde{f}}f_{2}>&\end{array}$.
\\Similarly for $L_{f}^{*}$.
\item $\forall h\in L^{2}(G),\,\,\,\forall g\in G,\,\,\,\left\{\begin{array}{r}(R_{x}L_{f}(h))(g)=(f*h)(gx)=\int_{G}f(gxy)h(y^{-1})dy\\(L_{f}R_{x}(h))(g)=(f*(R_{x}h))(g)=\int_{G}f(gz)h(z^{-1}x)dz\end{array}\right.$ are equivalent by the change of variables $z=xy$. 
\\Similarly for $R_{f}L_{x}=L_{x}R_{f}$.

\end{itemize}

\begin{definition}If $V_{1},V_{2}$ are two Hilbert spaces over $\mathbb{C}$ and $T:V_{1}\rightarrow V_{2}$ is a continuous $\mathbb{C}$-linear map, we say $T$ is compact if the set $\overline{T(B)}\subset V_{2}$ is compact, where $B=\{v\in V_{1}|||v||=1\}$.\end{definition}
\begin{lemma}\begin{enumerate}\item every finite rank operator is compact
\item the space of compact operators $T:V_{1}\rightarrow V_{2}$ is closed in $\hom(V_{1},V_{2})$\end{enumerate}\end{lemma}
pf.\begin{itemize}
\item The closed unit ball of a finite dimensional $\mathbb{C}$-vector space is compact, and therefore every finite rank operator is compact. 
\item Let $(T_{n})$ be a sequence of compact operators in $\hom(V_{1},V_{2})$ and suppose it converges to $T\in\hom(V_{1},V_{2})$. Let $(x_{n})$ be a sequence in $B$. Construct a sequence $x^{(r)}=(x_{n}^{(r)})$ of subsequences of $(x_{n})$ in the following inductive way; take $x^{(0)}=(x_{n})$. If $r\geq1,\,\,$ suppose $x^{r-1}$ is constructed and take for $x^{(r)}$ a subsequence of $x^{(r-1)}$ s.t. $(T_{r}(x_{n}^{(r)})$ is a Cauchy sequence. Set $y_{n}=x_{n}^{(n)}$. Let $\epsilon>0$. Choose $m\geq1$ s.t. $||T-T_{m}||\leq \epsilon$ and $N\geq m$ s.t. $\forall r,s\geq N,\,\,\,\,||T_{m}(y_{r})-T_{m}(y_{s})||\leq\epsilon$. Then if $r,s\geq N,\,\,\,||T(y_{r})-T(y_{s})||\leq||(T-T_{m})(y_{r})||+||T_{m}(y_{r}-y_{s})||+||(T_{m}-T)(y_{s})||\leq 3\epsilon,$ given $y_{s},y_{r}\in B$. Hence $(T(y_{n}))$ is a Cauchy sequence. Since $V_{2}$ is complete, the lemma follows. 


\end{itemize}
\begin{theorem}Let $V_{1},V_{2}$ be Hilbert spaces. Write $\hom(V_{1},V_{2})$ for the space of continuous $\mathbb{C}$-linear maps from $V_{1}$ to $V_{2}$ and consider the topology on it defined by the operator norm. Then $\forall T\in\hom(V_{1},V_{2})$, TFAE\begin{enumerate}\item $T$ is compact\item $T$ is a limit of finite rank element of $\hom(V_{1},V_{2})$\end{enumerate}\end{theorem}
pf. \begin{itemize}\item $2\rightarrow 1$ follows by the previous lemma.
\item $1\rightarrow 2$
\\\\pf. Let $T\in \hom(V_{1},V_{2})$ be a compact operator. Let $(W_{n})$ be a sequence of finite dimensional subspace of $V_{2}$ s.t. $V_{2}=\cup W_{n}$, and $\forall n\geq0,\,\,\,\pi_{n}:V_{2}\rightarrow W_{n}$ be the orthogonal projection. Set $T_{n}=\pi_{n}\circ T$. The $\forall x\in V_{1},\,\,\,\,T_{n}(x)\rightarrow T(x)$ as $n\rightarrow\infty$. Let $\epsilon>0$. As $\overline{T(B)}$ is compact, $\exists x_{1},\dots,x_{r}\in B$ s.t. $\forall y\in B,\,\,\,\exists i\in\{1,\dots,r\}$ s.t. $||T(x_{i})-T(y)||\leq\epsilon$. Then $\forall n\geq0,\,\,\,||T_{n}(y)-T_{n}(x_{i})||=||\pi_{n}(T(y)-T(x_{i}))||\leq||T(y)-T(x_{i})||\leq\epsilon$. Choose $N\geq0$ s.t. $\forall i,\,\,\,\forall n\geq N,\,\,\,\,\,||T_{n}(x_{i})-T(x_{i})||\leq\epsilon$. Then fix $n\geq N,$ and if $y\in B$ choose $i\in\{1,\dots,r\}$ s.t. $||T(y)-T(x_{i})||\leq \epsilon$. Then $||T(y)-T_{n}(y)||\leq||T(y)-T(x_{i})||+||T(x_{i})-T_{n}(x_{i})||+||T_{n}(x_{i})-T_{n}(y)||\leq3\epsilon$. 

\end{itemize}
\begin{definition}Let $V$ be a Hilbert space, and $T$ be a continuous endomorphism of $V$. We say $\lambda\in\mathbb{C}$ is an eigenvalue of $T$ if $\ker(T-\lambda id_{V})\neq\{0\}$. We denote by Spec$(T)\subset\mathbb{C}$ the set of eigenvalues of $T$ and call it the spectrum of $T$.\end{definition}
\begin{theorem}[Spectral Theorem]Let $V$ be a Hilbert space over $\mathbb{C}$, and let $T:V\rightarrow V$ be a continuous endomorphism of $V$. Assume $T$ is compact and self-adjoint, and write $V_{\lambda}=\ker(T-\lambda id_{V}),\,\,\,\,\,\forall \lambda\in \mathbb{C}$. Then $\left\{\begin{array}{r}\text{Spec}(T)\subset\mathbb{R}\\$ if $\lambda,\mu\in\text{Spec}(T),\,\,\,\lambda\neq\mu,$ then $V_{\mu}\subset V_{\lambda}^{\perp}\\ $ if $\lambda\in\text{Spec}(T)\setminus\{0\},\,\,\,\dim_{\mathbb{C}}V_{\lambda}<\infty\\\text{Spec}(T)$ is finite or countable$\\\bigoplus_{\lambda\in\text{Spec}(T)}V_{\lambda}$ is dense in $V\end{array}\right.$.\end{theorem}
 

\begin{theorem}$\forall f\in L^{2}(G),\,\,\,\,\,\exists (f_{n})_{n\geq0}$ a sequence of functions in $L^{2}(G)$ s.t. $\\\left\{\begin{array}{r}\tilde{f}_{n}=f_{n},\,\,\,\forall n\geq0\\||f_{n}||_{2}=1,\,\,\,\forall n\geq0\\f*f_{n}\rightarrow f$ in $L^{2}(G)$ as $n\rightarrow\infty\end{array}\right.$.\end{theorem}
pf. Let $\epsilon>0$. Choose a nhood $U$ of 1 in $G$ s.t. $U=U^{-1}$ and $\forall x,y\in U,\,\,\,||R_{x}f-f||_{2}<\epsilon$. Let $h_{\epsilon}=\frac{1}{vol(U)}\mathbb{1}_{U}$. Then $\tilde{h}_{\epsilon}=h_{\epsilon}$ and $||h_{\epsilon}||_{2}=1$. Also $\\(f*h_{\epsilon})(g)=\int_{G}f(gx^{-1})h_{\epsilon}(x)dx=\frac{1}{vol(U)}\int_{U}f(gx^{-1})dx\\\rightarrow (f*h_{\epsilon})(g)-f(g)=\frac{1}{vol(U)}\int_{U}(R_{x^{-1}}(f)(g)-f(g))dx$. This gives $\\||(f*h_{\epsilon})-f||_{2}^{2}=\int_{G}|(f*h_{\epsilon})(g)-f(g)|^{2}dg=\frac{1}{vol(U)^{2}}\int_{G\times U\times U}(R_{x^{-1}}f-f)(g)\overline{R_{y^{-1}}f-f(g)}dxdydg\\=\frac{1}{vol(U)^{2}}\int_{U\times U}<R_{x^{-1}}f-f,R_{y^{-1}}f-f>dxdy\leq\frac{1}{vol(U)^{2}}\int||R_{x^{-1}}f-f||_{2}||R_{y^{-1}}f-f||_{2}dxdy\\\leq\frac{1}{vol(U)^{2}}\int \epsilon^{2}dxdy=\epsilon^{2}$ by the CS inequality and the choice of $U$. If we take $f_{n}=h_{1/2^{n}}$, the theorem follows. 
\begin{lemma}If $f_{1},f_{2}:G\rightarrow \mathbb{C}$, define $\begin{array}{rrr}f_{1}\otimes f_{2}:G\times G&\rightarrow&\mathbb{C}\\(g_{1},g_{2})&\mapsto&f_{1}(g_{1})f_{2}(g_{2})\end{array}$. This induces a $\mathbb{C}$-linear map $L^{2}(G)\otimes_{\mathbb{C}}L^{2}(G)\rightarrow L^{2}(G\times G)$ which is injective with dense image.\end{lemma}
pf. Take the Hilbert basis $(e_{i})_{i\in I}$ of $L^{2}(G)$. Then the family $(e_{i}\otimes e_{j})$ of $L^{2}(G\times G)$ is orthonormal and any function on $L^{2}(G\times G)$ that is orthogonal to all $e_{i}\otimes e_{j}$ is 0 a.e. hence 0. 
\begin{lemma}$\forall K\in L^{2}(G\times G)$ define $\begin{array}{rrr}T_{K}:L^{2}(G)&\rightarrow&L^{2}(G)\\f(x)&\mapsto&\int_{G}K(x,y)f(y)dy\end{array}$. Then $||T_{K}||\leq||K||_{2}.$\end{lemma}
\begin{theorem}$\forall f\in L^{2}(G)$, the endomorphisms $R_{f}$ and $L_{f}$ of $L^{2}(G)$ are compact.\end{theorem}
pf. Consider $\begin{array}{rrr}K:G\times G&\rightarrow&\mathbb{C}\\(x,y)&\mapsto&f(xy^{-1})\end{array}$. Then $K\in L^{2}(G\times G)$, so by the previous lemma, $\exists (g_{i})_{i\in I},\,\,\,(h_{i})_{i\in I}$ in $L^{2}(G)$ s.t. $\sum g_{i}\otimes h_{i}$ converges to $K$ in $L^{2}(G\times G)$. Here $T_{K}=L_{f}$. For every finite subset $J$ of $I$, let $S_{J}=\sum_{i,j\in J}g_{i}\otimes h_{j}\in L^{2}(G\times G)$ and $T_{J}=T_{S_{J}}$. Then $\forall h\in L^{2}(G),\,\,\,\forall x\in G,\,\,\,\,(T_{J}h)(x)=\sum_{(i,j)\in J^{2}}\int_{G}h(y)g_{i}(x)h_{j}(y)dy=\sum_{i}\sum_{j}\int_{G}h(y)h_{j}(y)dyg_{i}(x)$, i.e. $\forall h\in L^{2}(G),\,\,\,T_{J}h$ is in the finite-dimensional subspace of $L^{2}(G)$ spanned by $g_{i},\,\,\,i\in J$. Hence the operator $T_{J}$ has finite rank. By the previous lemma, $||T_{J}||_{2}\leq||S_{J}||_{2}\rightarrow ||K||_{2}$ so $L_{f}$ is the limit of the operator $T_{J}$ and therefore $L_{f}$ is compact by the previous theorem. 
\begin{definition}We make $G\times G$ act on $L^{2}(G)$ by $(x,y)\cdot f=L_{x^{-1}}R_{y}f=R_{y}L_{x^{-1}}f$, i.e. $((x,y)\cdot f)(g)=f(x^{-1}gy),\,\,\,\forall g\in G$.\end{definition}
\begin{proposition}The representation of $G\times G$ on $L^{2}(G)$ defined above is continuous and so it is a unitary representation.\end{proposition}
pf. $L_{x^{-1}}R_{y}$ is a unitary endomorphism of $L^{2}(G),\,\,\,\forall x,y\in G$, so if the representation is continuous, the proposition follows. Fix $f_{0}\in L^{2}(G)$ and let $x_{0},y_{0}\in G,$ and $\epsilon>0$. Choose $f:G\rightarrow\mathbb{C}$ continuous s.t. $||f_{0}-f||_{2}\leq\epsilon/3$. As $G$ is compact and $f$ is continuous, $f$ is uniformly continuous, so there exists a nhood $U$ of 1 in $G$ s.t. $\forall x,y\in U,\,\,\,\forall g\in G,\,\,\,|f(x^{-1}gy)-f(g)|\leq\epsilon/3$. If $x\in x_{0}U,\,\,\,y\in y_{0}U,\,\,\,\,||L_{x^{-1}}R_{y}f-L_{x_{0}^{-1}}R_{y_{0}}f||_{2}=\sqrt{\int_{G}|f(x^{-1}gy)-f(x_{0}^{-1}gy_{0})|^{2}dg}\leq \epsilon/3$. Then $||L_{x^{-1}}R_{y}f_{0}-L_{x_{0}^{-1}}R_{y_{0}}f_{0}||_{2}\leq ||L_{x^{-1}}R_{y}f_{0}-L_{x^{-1}}R_{y}f||_{2}+|L_{x^{-1}}R_{y}f-L_{x_{0}^{-1}}R_{y_{0}}f||_{2}+||L_{x_{0}^{-1}}R_{y_{0}}f-L_{x_{0}^{-1}}R_{y_{0}}f_{0}||_{2}\leq \epsilon$, so the map $\begin{array}{rrr}G\times G&\rightarrow&L^{2}(G)\\(x,y)&\mapsto&L_{x^{-1}}R_{y}f_{0}\end{array}$ is continuous. Then by the previous proposition A.3.1., the representation is continuous. 
\\\\\indent Assume $G$ is a compact Hausdorff group.
\begin{definition}Denote by $\hat{G}$ the set of isomorphism classes of irreducible continuous representations of $G$ on finite dimensional complex vector spaces.\end{definition}
Let $(V_{\rho},\rho)\in\hat{G}$. Fix a $G$-invariant Hermitian inner product $<\cdot,\cdot>$ on $V_{\rho}$. Then $G\times G$ acts on End$_{\mathbb{C}}(V_{\rho})\cong V_{\rho}^{*}\otimes_{\mathbb{C}}V_{\rho}$ by $(g_{1},g_{2})\cdot u=\rho(g_{1})u\rho(g_{2})^{-1}$. 
\begin{definition}The Hilbert-Schmidt inner product on End$_{\mathbb{C}}(V_{\rho})$ is given by the formula \begin{equation*}<u,v>_{HS}=T(uv^{*})\end{equation*} where $v^{*}$ is the adjoint of $v$ for the chosen $G$-invariant inner product $<\cdot, \cdot>$ on $V_{\rho}$.\end{definition}
$\rightarrow $Hermitian inner product and $G\times G$-invariant given $\rho(g)$ begin unitary $\forall g\in G$.


\begin{lemma}Let $V$ be a Hilbert space with a unitary representation of $G$ and let $V_{1},V_{2}\subset V$ be finite-dimensional irreducible subrepresentations of $G$. Then either $V_{1}\cong V_{2}$ as representations of $G$ or $<v_{1},v_{2}>=0,\,\,\,\forall v_{1}\in V_{1},\,\,\,v_{2}\in V_{2}$.\end{lemma}
pf. $V_{2}^{\perp}$ is stable by $G$ and $V= V_{2}\oplus V_{2}^{\perp}$. So the orthogonal projection $\pi:V\rightarrow V_{2}$ is $G$-equivariant, hence the composition $V_{1}\subset V\rightarrow V_{2}$ is $G$-equivariant. If $V_{1}\not\cong V_{2}$, then this $G$-equivariant map $V_{1}\rightarrow V_{2}$ has to be 0 by Schur's lemma, which means $V_{1}\subset V_{2}^{\perp}$. 
\\\\\indent Define $\begin{array}{rrr}\iota_{\rho}:\text{End}_{\mathbb{C}}(V_{\rho})&\rightarrow &L^{2}(G)\\u&\mapsto&(g\mapsto T(\rho(g)^{-1}\circ u))\end{array}\,\,\,\,(\,\,\,\,\,\forall u\in \text{End}_{\mathbb{C}}(V_{\rho})$, the function $g\mapsto T(\rho(g)^{-1}\circ u)$ is continuous on $G$, hence on $L^{2}(G)$, given $G$ being compact. So it is well-defined ).
\begin{theorem}\begin{enumerate}\item $\forall (V_{\rho},\rho)\in\hat{G},$ End$_{\mathbb{C}}(V_{\rho})$ is an irreducible representation of $G\times G$. Also if $(V_{\rho},\rho)\not\cong(V_{\rho'},\rho')$, then End$_{\mathbb{C}}(V_{\rho})\not\cong\text{End}_{\mathbb{C}}(V_{\rho'})$\item $\forall (V_{\rho},\rho)\in\hat{G}$, the map $\iota_{\rho}:\text{End}_{\mathbb{C}}(V_{\rho})\rightarrow L^{2}(G)$ is $G\times G$-equivariant, the map $\dim(V_{\rho})^{1/2}\iota_{\rho}$ is an isometry, and $(\dim_{\mathbb{C}}V_{\rho})\iota_{\rho}$ sends the composition in End$_{\mathbb{C}}(V_{\rho})$ to the convolution product in $L^{2}(G)$\item the map $\iota=\bigoplus_{\rho\in \hat{G}}\iota_{\rho}:\bigoplus_{\rho\in\hat{G}}\text{End}_{\mathbb{C}}(V_{\rho})\rightarrow L^{2}(G)$ is injective and has dense image, i.e. as a representation of $G\times G$ and as $\mathbb{C}$-algebra, \begin{equation*}L^{2}(G)\cong\bigoplus_{\rho\in\hat{G}}\widehat{\text{End}_{\mathbb{C}}(V_{\rho})}\end{equation*}\end{enumerate}\end{theorem}
pf. 

\begin{itemize}\item let $\rho\in\hat{G}$. Then $\chi_{\text{End}_{\mathbb{C}}(V_{\rho})}(g_{1},g_{2})=\chi_{V_{\rho}}(g_{1})\overline{\chi_{V_{\rho}}(g_{2})}$ by the definition of the action of $G\times G$ on $\text{End}_{\mathbb{C}}(V_{\rho})$. By the Schur's orthogonality formula, $\int_{G\times G}|\chi_{\text{End}_{\mathbb{C}}(V_{\rho})}(g_{1},g_{2})|^{2}dg_{1}dg_{2}=(\int_{G}|\chi_{V_{\rho}}(g_{1})|^{2}dg_{1})(\int_{G}|\chi_{V_{\rho}}(g_{2})|^{2}dg_{2})=1$ so $\text{End}_{\mathbb{C}}(V_{\rho})$ is irreducible. Moreover, if $(V_{\rho'},\rho')\not\cong(V_{\rho},\rho)$, then by the Schur's orthogonality, $<\chi_{\text{End}_{\mathbb{C}}(V_{\rho})},\chi_{\text{End}_{\mathbb{C}}(V_{\rho'})}>_{L^{2}(G\times G)}=<\chi_{V_{\rho}},\chi_{V_{\rho'}}>_{L^{2}(G)}<\overline{\chi_{V_{\rho}}},\overline{\chi_{V_{\rho'}}}>_{L^{2}(G)}=0$ so $\text{End}_{\mathbb{C}}(V_{\rho})\not\cong\text{End}_{\mathbb{C}}(V_{\rho'})$.
\item fix $(V_{\rho},\rho)\in\hat{G}$. 
\\$\forall x,y\in G,\,\,\,\,\forall u\in\text{End}_{\mathbb{C}}(V_{\rho}),\,\,\,\forall g\in G,
\\\iota_{\rho}(\rho(x)u\rho(y)^{-1})(g)=T(\rho(g)^{-1}\rho(x)u\rho(y)^{-1})=T(\rho(x^{-1}gy)^{-1}u)=(L_{x^{-1}}R_{y}\iota_{\rho}(u))(g)\rightarrow \iota_{\rho}$ is $G\times G$-equivariant.
\\Let $u_{1},u_{2}\in \text{End}_{\mathbb{C}}(V_{\rho})$. End$_{\mathbb{C}}(V_{\rho})$ is spanned by the $T_{v,w}^{0},\,\,\,\forall v,w\in V_{\rho}$, so assume $\left\{\begin{array}{r}u_{1}=T_{v_{1},w_{1}}^{0}\\u_{2}=T_{v_{2},w_{2}}^{0}\end{array}\right.,\,\,\,\,\,\exists v_{1},v_{2},w_{1},w_{2}\in V_{\rho}$. Then,$\\\begin{array}{rr}<\iota_{\rho}(u_{1}),\iota_{\rho}(u_{2})>=\int_{G}T(\rho(g)^{-1}T_{v_{1},w_{1}}^{0})\overline{T(\rho(g)^{-1}T_{v_{2},w_{2}}^{0})}dg&\rho(g)T_{v,w}^{0}=T_{v,gw}^{0}\\=\int_{G}<g^{-1}v_{1},w_{1}>\overline{<g^{-1}w_{2},v_{2}>}dg&\\=\int_{G}<gv_{2},w_{2}>\overline{<gv_{1},w_{1}>}dg&\\=\frac{1}{\dim_{\mathbb{C}}V_{\rho}}<v_{2},v_{1}><w_{1},w_{2}>&$ by the theorem A.3.3. $\\=\frac{1}{\dim_{\mathbb{C}}V_{\rho}}<u_{1},u_{2}>_{HS}& $ by the proposition A.3.3. $\end{array}$. Thus $\dim(V_{\rho})^{1/2}\iota_{\rho}$ is an isometry.
\\Let $g\in G$. Then $\iota_{\rho}(u_{1}u_{2})(g)=T(\rho(g)^{-1}T_{v_{1},w_{1}}^{0}(T_{w_{2},v_{2}}^{0})^{*})=<\rho(g)^{-1}w_{1},v_{2}><w_{2},v_{1}>$ by the proposition A.3.3. On the other hand, \\$\begin{array}{rr}(\iota_{\rho}(u_{1})*\iota_{\rho}(u_{2}))(g)=\int_{G}T(\rho(x)^{-1}\rho(g)^{-1}T_{v_{1},w_{1}}^{0})T(\rho(x)T_{v_{2},w_{2}}^{0})dx&\\=\int_{G}<\rho(x)^{-1}\rho(g)^{-1}w_{2},v_{1}><\rho(x)w_{2},v_{2}>dx&\\=\int_{G}<\rho(x)w_{2},v_{2}>\overline{<\rho(x)v_{1},\rho(g)^{-1}w_{1}>}dx&\\=\frac{1}{\dim_{\mathbb{C}}(V_{\rho})}<\rho(g)^{-1}w_{1},v_{2}><w_{2},v_{1}>&$ by the theorem A.3.3. $\end{array}$. 
\item Let $u=(u_{\rho})\in\bigoplus_{\rho\in\hat{G}}\text{End}_{\mathbb{C}}(V_{\rho})$ be s.t. $\iota(u)=0$. Then $\sum\iota_{\rho}(u_{\rho})=0$. By the previous lemma, and the statement 1, the subspaces $\iota_{\rho}(\text{End}_{\mathbb{C}}(V_{\rho})$ are pairwise orthogonal to each other, so $\forall \rho\in\hat{G},\,\,\,0=<\iota_{\rho}(u_{\rho}),\iota(u)>=<\iota_{\rho}(u_{\rho}),\iota_{\rho}(u_{\rho})>=\frac{1}{\dim_{\mathbb{C}}V_{\rho}}<u_{\rho},u_{\rho}>$. Thus all $u_{\rho}=0$, i.e. $u=0$ and $\iota$ is injective.
\\\textbf{CLAIM : Every finite dimensional $G$-subrepresentation of $L^{2}(G)$ is contained in Im$(\iota)$, where we make $G$ act on $L^{2}(G)$ by the left regular action.}
\\pf. As finite dimensional representatons of $G$ are semisimple, if $\forall \rho\in\hat{G},\,\,\,\forall u:V_{\rho}\rightarrow L^{2}(G),\,\,\,\,G$-equivariant map, if Im$(u)\subset\text{Im}(\iota)$, then the claim follows. Fix $\rho\in\hat{G}$ and $u:V_{\rho}\rightarrow L^{2}(G),\,\,\,G$-equivariant map. Let $v\in V_{\rho}$. Then $\forall f\in L^{2}(G),\,\,\,\forall g\in G,\,\,\,\begin{array}{rr}(u(v)*f)(g)=\int_{G}u(v)(gx)f(x^{-1})dx&\\=\int_{G}L_{g}(u(v))(x)\overline{\tilde{f}(x)}dx&\\=<L_{g}(u(v)),\tilde{f}>_{L^{2}(G)}&\\=<u(\rho(g)^{-1}v),\tilde{f}>_{L^{2}(G)}&$ by the $G$-equivariance of $u\\=<\rho(g)^{-1}v,u^{*}(\tilde{f})>_{V_{\rho}}&\\=T(\rho(g)^{-1}T_{u^{*}(\tilde{f}),v})&$ by the definition of $T_{v,w}^{0}\\=\iota_{\rho}(T^{0}_{u^{*}(\tilde{f}),v})(g)&\end{array}$. Thus $u(v)*f\in\text{Im}(\iota_{\rho}),\,\,\,\,\forall f\in L^{2}(G)$. Then there exists a sequence $(f_{n})$ in $L^{2}(G)$ s.t. $u(v)*f_{n}\rightarrow u(v)$ as $n\rightarrow\infty$ by the theorem A.3.6. Thus $u(v)\in\overline{\text{Im}\iota_{\rho}}$. As Im$\iota_{\rho}$ is a finite-dimensional subspace of $L^{2}(G)$, it is closed in $L^{2}(G)$ so $u(v)\in \text{Im}\iota_{\rho}\subset\text{Im}\iota$.
\\Let $f\in \text{Im}(\iota)^{\perp}$. Then $f$ is orthogonal to every finite dimensional $G$-subrepresentation of $L^{2}(G)$. Let $h\in L^{2}(G)$ be s.t. $h=\tilde{h}$. Then the endomorphism $R_{h}$ of $L^{2}(G)$ is self-adjoint and compact by the theorem A.3.7. Hence by the spectral theorem, $\ker R_{h}$ is the orthogonal of the closure of the direct sum of the eigenspaces $\ker(R_{h}-\lambda id_{L^{2}(G)}),\,\,\,\lambda\in\mathbb{R}^{\times}$, which are all finite dimensional. As $R_{h}$ is $G$-equivariant and commutes with the left regular action, there eigenspaces are left stable by $G$ and so $f$ is orthogonal to all of them, hence $f\in\ker R_{h},\,\,\,\,f*h=0$. Then by the theorem A.3.6., there exists a sequence $(h_{n})_{n\geq0}$ of elements of $L^{2}(G)$ s.t. $\tilde{h}_{n}=h_{n},\,\,\,\,\forall n\geq0$ and $f*h_{n}\rightarrow f$ as $n\rightarrow\infty$. Then $f*h_{n}=0,\,\,\,\forall n\geq0$ and therefore $f=0$. So Im$\iota$ is dense in $L^{2}(G)$.  

\end{itemize}












\end{document}
